% **************************************************************************************************************
% A Classic Thesis Style
% An Homage to The Elements of Typographic Style
%
% Copyright (C) 2018 André Miede and Ivo Pletikosić
%
% If you like the style then I would appreciate a postcard. My address
% can be found in the file ClassicThesis.pdf. A collection of the
% postcards I received so far is available online at
% http://postcards.miede.de
%
% License:
% This program is free software; you can redistribute it and/or modify
% it under the terms of the GNU General Public License as published by
% the Free Software Foundation; either version 2 of the License, or
% (at your option) any later version.
%
% This program is distributed in the hope that it will be useful,
% but WITHOUT ANY WARRANTY; without even the implied warranty of
% MERCHANTABILITY or FITNESS FOR A PARTICULAR PURPOSE.  See the
% GNU General Public License for more details.
%
% You should have received a copy of the GNU General Public License
% along with this program; see the file COPYING.  If not, write to
% the Free Software Foundation, Inc., 59 Temple Place - Suite 330,
% Boston, MA 02111-1307, USA.
%
% PLEASE SEE ALSO THE AUTHORS' NOTE REGARDING THIS LICENSE
% IN THE DOCUMENTATION (ClassicThesis.pdf --> Chapter 1 / Chapter01.tex)
% **************************************************************************************************************
\RequirePackage{silence} % :-\
    \WarningFilter{scrreprt}{Usage of package `titlesec'}
    %\WarningFilter{scrreprt}{Activating an ugly workaround}
    \WarningFilter{titlesec}{Non standard sectioning command detected}
\documentclass[ twoside,openright,titlepage,numbers=noenddot,%1headlines,
                headinclude,footinclude,cleardoublepage=empty,abstract=on,
                BCOR=5mm,paper=a4,fontsize=11pt
                ]{scrreprt}

%********************************************************************
% Note: Make all your adjustments in here
%*******************************************************
\input{classicthesis-config}

%********************************************************************
% Bibliographies
%*******************************************************
\addbibresource{../../source/references.bib}

%********************************************************************
% Hyphenation
%*******************************************************
%\hyphenation{put special hyphenation here}
% For using pandoc's --csl= option with a csl file. It works pretty well but I
% wanted to preserve the bibliography section style. Finding a csl file that did
% that was hard so I went with the --biblatex flag.

\usepackage{longtable,booktabs,array}
\usepackage{calc} % for calculating minipage widths
% Correct order of tables after \paragraph or \subparagraph
\usepackage{etoolbox}
\makeatletter
\patchcmd\longtable{\par}{\if@noskipsec\mbox{}\fi\par}{}{}
\makeatother
% Allow footnotes in longtable head/foot
\IfFileExists{footnotehyper.sty}{\usepackage{footnotehyper}}{\usepackage{footnote}}
\makesavenoteenv{longtable}


\usepackage{amsthm}

% center floats
\makeatletter
\g@addto@macro\@floatboxreset\centering
\makeatother

% declare listing environment
\usepackage{newfloat}
\DeclareFloatingEnvironment{fragmento}
% reduce space between listing caption and actual listing, taken from:
% https://stackoverflow.com/questions/2318598/how-to-reduce-the-seperation-from-other-text-using-latex-minted
\setlength\partopsep{-\topsep}
\usepackage{mdframed} % TODO: used in TODO notes, remove
\usepackage{placeins}
\usepackage{algorithm}
\floatname{algorithm}{Algoritmo}
\usepackage[noabbrev,capitalize,nameinlink,spanish]{cleveref}
\creflabelformat{equation}{#2\textup{#1}#3}
\crefname{section}{sección}{secciones}
\Crefname{section}{Sección}{Secciones}
\crefname{subsection}{sección}{secciones}
\Crefname{subsection}{Sección}{Secciones}
\crefname{listing}{fragmento}{fragmentos}
\Crefname{listing}{Fragmento}{Fragmentos}
\crefname{table}{tabla}{tablas}
\Crefname{table}{Tabla}{Tablas}

% This needs to be after loading cleveref
\newtheorem{theorem}{Teorema}[chapter]
\newtheorem{definition}[theorem]{Definición}
\newtheorem{remark}[theorem]{Observación}
\newtheorem{property}[theorem]{Propiedad}
\newtheorem{notation}[theorem]{Notación}

\newcommand{\R}{\mathbb{R}}
\newcommand{\RR}[1]{\R^{#1 \times #1}}
\newcommand{\so}[1]{\mathfrak{so(#1)}}
\newcommand{\se}[1]{\mathfrak{se(#1)}}
\newcommand{\Rnn}{\R^{n \times n}}
\renewcommand{\H}{\mathbb{H}}
\newcommand{\N}{\mathbb{N}}
\newcommand{\dotprod}[2]{\langle #1, #2 \rangle}
\newcommand{\norm}[1]{\| #1 \|}
\newcommand{\passthrough}[1]{#1}
\renewcommand{\dot}[1] {\overset{\,_{\mbox{\Large .}}}{#1}}

% If using pdflatex, spanish accents will not render. You'll need this. (also you won't be able to use fira code ligatures)
% \lstset{%
%   literate={á}{{\'a}}1 {é}{{\'e}}1 {í}{{\'i}}1 {ó}{{\'o}}1 {ú}{{\'u}}1
% }

\providecommand{\tightlist}{%
  \setlength{\itemsep}{0pt}\setlength{\parskip}{0pt}}

\newcommand{\fig}[4]{%
\begin{figure}[bth]
  \myfloatalign
  \includegraphics[width=1\linewidth]{#2}
  \caption[#3]{#4}
  \label{#1}
\end{figure}
}
\newcommand{\figw}[5]{%
\begin{figure}[bth]
  \myfloatalign
  \includegraphics[width=#5]{#2}
  \caption[#3]{#4}
  \label{#1}
\end{figure}
}
\newcommand{\figref}{\autoref}
\newcommand{\mono}[1]{\passthrough{\lstinline!#1!}}
% \newcommand{\mono}{\texttt}
\renewcommand{\bold}{\spacedlowsmallcaps}
% \renewcommand{\bold}{\textbf}
% \newcommand{\italic}{\textit}
\newcommand{\italic}{\emph}

\newcounter{mparcnt}
\renewcommand\themparcnt{${}^{\arabic{footnote}}$}
\newcommand\marginnote[1]{\stepcounter{footnote}\themparcnt \marginpar{\themparcnt #1}}

\newcommand{\HRule}{\rule{.9\linewidth}{.6pt}} % New command to make the lines in the

\renewcommand{\lstlistingname}{Fragmento}
\renewcommand{\lstlistlistingname}{Lista de \lstlistingname s}
% ********************************************************************
% GO!GO!GO! MOVE IT!
%*******************************************************
\begin{document}
\ActivateVerbatimLigatures
\frenchspacing
\raggedbottom
\selectlanguage{spanish} % american ngerman
%\renewcommand*{\bibname}{new name}
%\setbibpreamble{}
\pagenumbering{roman}
\pagestyle{plain}
%********************************************************************
% Frontmatter
%*******************************************************
% \include{FrontBackmatter/DirtyTitlepage}
%*******************************************************
% Titlepage
%*******************************************************
\begin{titlepage}
    \begin{addmargin}[-0.6cm]{-4.4cm}
    \begin{center}

    \vspace*{.06\textheight}
    % {\scshape\LARGE \myUni \\ \myLocation\par}\vspace{1.5cm} % University name
    {\scshape\LARGE \myUni\par}\vspace{1.5cm} % University name
    \textsc{\Large Trabajo Especial}\\[0.5cm] % Thesis type

    \HRule \\[0.4cm] % Horizontal line
    % {\huge \bfseries \spacedallcaps{\myTitle}\par}\vspace{0.4cm} % Thesis title
    {\Large \bfseries \spacedallcaps{\myTitle}\par}\vspace{0.4cm} % Thesis title
    \HRule \\[1.5cm] % Horizontal line

    \begin{minipage}[t]{0.4\textwidth}
    \begin{flushleft} \large
    \emph{Autor:}\\
    \href{https://www.famaf.unc.edu.ar/~mdemayo366/}{\myName} % Author name - remove the \href bracket to remove the link
    \end{flushleft}
    \end{minipage}
    \begin{minipage}[t]{0.4\textwidth}
    \begin{flushright} \large
    \emph{Director:} \\
    \href{https://cs.famaf.unc.edu.ar/~nicolasw/}{\mySupervisor} % Supervisor name - remove the \href bracket to remove the link
    \end{flushright}
    \end{minipage}\\[2cm]

    \large \textit{\myDegree}\\[0.5cm] % University requirement text
    \myUni\\ \myLocation\\[2cm] % Research group name and department name

    \href{http://creativecommons.org/licenses/by/4.0/}{\includegraphics[height=1cm]{style/ccby.pdf}}\\[0.3cm]
    \textit{Esta obra está bajo una Licencia Creative Commons\\ Atribución 4.0 Internacional.}\\[1cm]


    % \vfill

    {\large \myTime}

    \vfill

    \begin{center}
        \includegraphics[height=2cm]{style/unc.pdf}
        \includegraphics[height=2cm]{style/famaf.pdf}
    \end{center}
    % \vfill


    \end{center}
    \end{addmargin}
\end{titlepage}

\cleardoublepage\thispagestyle{empty}

\hfill

\vfill

\noindent\myName: \textit{\myTitleOneLine.} \mySubtitle. \myTime.\\
\\
Esta obra está licenciada bajo la Licencia Creative Commons Atribución 4.0 Internacional.
Para ver una copia de esta licencia, visite \\
\url{http://creativecommons.org/licenses/by/4.0/}
o envíe una carta a Creative Commons, PO Box 1866, Mountain View, CA 94042, USA.
%\bigskip
%
%\noindent\spacedlowsmallcaps{Supervisors}: \\
%\myProf \\
%\myOtherProf \\
%\mySupervisor
%
%\medskip
%
%\noindent\spacedlowsmallcaps{Location}: \\
%\myLocation
%
%\medskip
%
%\noindent\spacedlowsmallcaps{Time Frame}: \\
%\myTime

% \cleardoublepage\include{FrontBackmatter/Dedication}
%\cleardoublepage\include{FrontBackmatter/Foreword}
%*******************************************************
% Abstract
%*******************************************************
%\renewcommand{\abstractname}{Abstract}
\pdfbookmark[1]{Abstract}{Abstract}
% \addcontentsline{toc}{chapter}{\tocEntry{Abstract}}
\begingroup
\let\clearpage\relax
\let\cleardoublepage\relax
\let\cleardoublepage\relax

\chapter*{Resumen}

Las aplicaciones de realidad virtual (VR), aumentada (AR) y sus derivadas,
englobadas dentro del término XR, necesitan métodos para localizar y entender
los movimientos realizados por el usuario y así poder actualizar la simulación
que ejecutan de manera acorde. El problema de localización ha sido uno de los
mayores desafíos a la hora de producir dispositivos de realidad virtual para
electrónica de consumo en los últimos años. Afortunadamente, han ocurrido
grandes avances en la localización visual-inercial mediante cámaras y sensores
inerciales; dados principalmente gracias a su continua mejora y abaratamiento
relacionados con su producción masiva para telefonía móvil. A pesar de esto,
sistemas de este tipo que puedan ser utilizados para XR se desarrollan
mayormente como parte de productos privativos. En este trabajo, se estudian
sistemas de localización visual-inercial provenientes de la academia y se
detallan las dificultades propias de integrarlos y adaptarlos para aplicaciones
de XR. En particular se los integra con Monado, la implementación de código
libre del estándar OpenXR. De esta forma, se presenta una de las primeras
soluciones completamente libres que permite localizar dispositivos XR con este
tipo de métodos visual-inerciales en sistemas operativos basados en GNU/Linux.

\vfill

\begin{otherlanguage}{american}
\chapter*{Abstract}

Applications for virtual reality, augmented reality, and their
derivatives, all of them encompassed in the term XR, need ways to keep track of
motion in the physical world to be able to reflect these movements accordingly
in the simulated world. Tracking has been one of the main
challenges in bringing VR devices into mainstream consumer electronics in recent
years. Fortunately, there has been breakthroughs in visual-inertial
tracking with cameras and inertial sensors due to their continuous improvement and
reduced costs thanks to the massive production tied to the mobile phone industry. However,
systems that perform this kind of tracking are developed for mostly privative
products. In this work, we study visual-inertial tracking systems that come
from academia, and detail the difficulties that come from integrating and
adapting them for XR applications. In particular, we integrate them to Monado,
the open-source implementation of the OpenXR standard. Therefore, we present one
of the first completely open solutions that integrates this kind of tracking for
XR applications and devices on GNU/Linux-based operating systems.

\end{otherlanguage}

\endgroup

\vfill

% \cleardoublepage\include{FrontBackmatter/Publications}
% \cleardoublepage\include{FrontBackmatter/Acknowledgments}
\cleardoublepage\include{FrontBackmatter/Contents}

% Mainmatter and backmatter start (appendix)
\cleardoublepage
\pagestyle{scrheadings}
\pagenumbering{arabic}

\cleardoublepage
\ctparttext{Esta primera parte está enfocada a intentar entender el problema
encarado en este trabajo. Además de una \italic{introducción} general a algunos
de los conceptos de XR y SLAM, veremos los \italic{fundamentos} que serán
utilizados de referencia a lo largo del escrito.}

\hypertarget{preliminares}{%
\part{Preliminares}\label{preliminares}}

\hypertarget{introducciuxf3n}{%
\chapter{Introducción}\label{introducciuxf3n}}

\hypertarget{localizaciuxf3n-en-xr}{%
\section{Localización en XR}\label{localizaciuxf3n-en-xr}}

Los sistemas de \emph{realidad virtual} y \emph{realidad aumentada}, o
\emph{VR} y \emph{AR} respectivamente por sus siglas en inglés, intentan
valerse de nuestras formas usuales de interactuar con el mundo que nos
rodea, para así generar simulaciones con características muy
particulares y difíciles de obtener de otra forma.

En el caso de VR uno usualmente tiene cascos con pantallas o, en inglés,
\emph{head mounted displays (HMD)}, y algún tipo de controles manejados
con ambas manos. Estos permiten aplicaciones con un alto nivel de
inmersión. La característica de poder adentrarse profundamente en la
simulación son deseadas para fines que varían desde el entretenimiento,
generalmente en forma de videojuegos, hasta el entrenamiento para
situaciones críticas, como lo pueden ser operaciones médicas o la
navegación de aeronaves. Por el lado de AR, se suele pensar en
dispositivos que permiten agregar o ``aumentar'' lo que uno ve con
información adicional. Estos pueden ser celulares con cámaras o incluso
lentes con pantallas transparentes que superponen información 2D o
modelos 3D sobre la escena real. Además existen otros términos como
\emph{realidad mixta} (\emph{MR}) que intentan hacer referencia a los
distintos matices en los que la realidad puede ser combinada con la
simulación; es común utilizar el término \emph{XR} para englobar a este
espectro de sistemas, a veces también referido como \emph{realidad
extendida}. En general, todos ellos comparten problemas similares que
son propios de la percepción humana y están relacionados con nuestra
psicología y fisiología \marginnote{%\
Al considerar sistemas de VR deben tenerse en cuenta, además de los desafíos
de ingeniería, los problemas propios de comprender como nuestra percepción
afecta lo que experimentamos. Lograr entender estos procesos cognitivos lo suficiente
como para valernos de ellos a la hora de diseñar sistemas de XR es fundamental
para obtener experiencias inmersivas. Esto puede pensarse como un trabajo de
``ingeniería inversa'' sobre nuestra propia biología \autocite{lavalleVirtualReality}.
}.

Para producir experiencias de XR, es necesario coordinar un gran número
de sistemas internos que se encuentran a lo largo del hardware y del
software para XR. Uno de los problemas fundamentales y particularmente
complicados en XR es el de la \emph{localización} o \emph{tracking};
esta es la capacidad del sistema de identificar la posición y
orientación del dispositivo XR, su \emph{pose}, en el entorno para poder
reflejarla en la experiencia simulada. En este trabajo nos enfocaremos
en el problema de localización. Para una lectura más comprensiva sobre
el resto de los problemas a los que XR se enfrenta, y en particular VR,
se recomienda el trabajo de \textcite{lavalleVirtualReality}.

Hay muchas formas de encarar el problema de localización y en general
todas requieren del uso inteligente de sensores físicos, ya sean
mecánicos, inerciales, magnéticos, ópticos, o incluso acústicos
\autocite{welchMotionTrackingNo2002}. Todos estos métodos de tracking
comparten un problema en común: los sensores son imperfectos y sus
mediciones son \emph{ruidosas}. Tomemos por ejemplo uno de los paquetes
de sensores más comúnmente utilizados en esta área, una \emph{unidad de
medición inercial} o \emph{IMU} por sus siglas en inglés. Estas integran
\emph{giroscopios} para la medición de velocidad angular,
\emph{acelerómetros} para la aceleración lineal, y algunas veces también
incluyen \emph{magnetómetros}. En teoría, si sus mediciones fueran
perfectas, una IMU debería proveer suficiente información para
determinar la pose en el espacio de un dispositivo que la contenga. Pero
incluso las IMU de mayor calidad acumulan tanto ruido que la integración
de sus mediciones durante cortos períodos de tiempo devuelve poses que
tienen un error de cientos de metros, ver
\Cref{tbl:imu-accumulated-error}. Para contrarrestar la naturaleza
imperfecta de sensores físicos como estos, los enfoques más exitosos de
localización emplean una combinación de múltiples sensores junto a
algoritmos de fusión ingeniosos capaces de integrar tipos de muestras
tremendamente distintos en una estimación de la pose que es
suficientemente buena.

\begin{longtable}[]{@{}llllll@{}}
\caption{\label{tbl:imu-accumulated-error} Error acumulado luego de
cierto tiempo de integrar mediciones de IMU de distinta
calidad\autocite[secc 3.3]{InertialNavigationPrimer}. Vale aclarar que
en dispositivos XR las IMU utilizadas son del tipo \emph{consumidor} por
su menor costo.}\tabularnewline
\toprule
\textbf{Categoría / Tiempo} & \textbf{1 s} & \textbf{10 s} & \textbf{60
s} & \textbf{10 min} & \textbf{1 hr} \\
\midrule
\endfirsthead
\toprule
\textbf{Categoría / Tiempo} & \textbf{1 s} & \textbf{10 s} & \textbf{60
s} & \textbf{10 min} & \textbf{1 hr} \\
\midrule
\endhead
Consumidor & 6 cm & 6.5 m & 400 m & 200 km & 39.000 km \\
Industrial & 6 mm & 0.7 m & 40 m & 20 km & 3.900 km \\
Táctico & 1 mm & 8 cm & 5 m & 2 km & 400 km \\
Navegación & \textless1 mm & 1 mm & 50 cm & 100 m & 10 km \\
\bottomrule
\end{longtable}

En años recientes, la llamada localización \emph{visual-inercial (VI)}
ha cobrado gran popularidad en el ecosistema de XR. Al emplear cámaras,
usualmente al menos dos, junto a una IMU dentro de un dispositivo XR
(p.~ej. un HMD o un celular). La localización VI estima la pose del
agente haciendo que el mismo dispositivo ``mire'' hacia su entorno para
ganar entendimiento del mismo. Uno de los principales beneficios de este
tipo de tracking es que tener todos los sensores empaquetados dentro del
dispositivo resulta muy conveniente para el usuario, ya que no necesita
colocar ni configurar ningún tipo de sensor externo en su entorno. Este
tipo de localización para aplicaciones de XR se consideraba inviable
hasta hace un poco más de una década, pero avances en la capacidad de
cómputo y en las técnicas de fusión han hecho posible que este ya no sea
el caso.

Ya existen dispositivos XR que emplean localización visual-inercial de
forma exitosa en productos comerciales como el \emph{Meta
Quest}\footnote{\url{https://www.oculus.com/quest-2/}}, los cascos
\emph{Windows Mixed Reality}\footnote{\url{https://www.microsoft.com/en-us/mixed-reality/windows-mixed-reality}}
o incluso los SDK \emph{ARCore}\footnote{\url{https://developers.google.com/ar}}
y \emph{ARKit}\footnote{\url{https://developer.apple.com/augmented-reality/arkit/}}
utilizados en celulares inteligentes. No obstante, todas estas
soluciones son propietarias, por lo tanto no hay manera de mejorarlos,
reusarlos en nuevos proyectos o productos, o incluso aprender de su
implementación a menos que se adquieran licencias especiales de sus
fabricantes.

\hypertarget{slam-y-vio-para-localizaciuxf3n-visual-inercial}{%
\section{SLAM y VIO para localización
visual-inercial}\label{slam-y-vio-para-localizaciuxf3n-visual-inercial}}

Afortunadamente, el área académica en navegación visual-inercial ha
experimentado un gran desarrollo en estas últimas décadas. No solo es un
problema central en áreas como las de robótica, sino que además es
atractiva por combinar una diversa cantidad de áreas como visión por
computadora, fusión de sensores, optimización, estimación
probabilística, entre otras. Esta ola de investigación ha dado lugar a
una gran cantidad de implementaciones de software libre, con distintos
grados de rendimiento, robustez, precisión, aplicaciones, facilidad de
uso, entre otras propiedades de interés. Recursos como
OpenSLAM\footnote{\url{http://openslam.org}} y estudios como los de las
referencias \autocite{servieresVisualVisualInertialSLAM2021} o
\autocite{taketomiVisualSLAMAlgorithms2017} pueden listar docenas de
sistemas disponibles para considerar. Más aún, cada año nuevos sistemas
aparecen mientras otros dejan de ser mantenidos. Seguir los avances del
área puede ser desafiante, pero esto es una consecuencia de su
desarrollo tan activo.

Específicamente, en este trabajo trataremos con sistemas de
\emph{localización y mapeo simultáneo (SLAM)} mediante sensores
visuales-inerciales (VI-SLAM). Como su nombre lo indica, SLAM intenta
construir un mapa del entorno en el que el agente XR se encuentra y
localizarlo en el mismo de forma simultánea; usualmente, sin contar con
información a priori sobre su pose ni el entorno en el que se encuentra.
Hay múltiples maneras de implementar SLAM y VI-SLAM, pero los sistemas
en los que nos concentraremos en este trabajo usan mediciones de la IMU
a altas frecuencias (p.~ej. 200 Hz) que miden el \emph{``movimiento
interno''} que el agente experimenta, también conocidas como mediciones
\emph{propioceptivas}, junto a muestras más lentas (p.~ej. 20 Hz) de
cámaras, usualmente un par de ellas, que dan información acerca de como
el entorno está cambiando alrededor del agente cuando este se mueve.
Estas son llamadas mediciones \emph{exteroceptivas} y ayudan a corregir
las mediciones ruidosas de la IMU.

El mapa que se forma en VI-SLAM suele crearse a partir de \emph{puntos
de interés} (\emph{landmarks}) en la escena que fueron triangulados y
detectados durante múltiples muestras como \emph{esquinas} o
\emph{features}\marginnote{%\
La terminología para definir puntos de interés es un poco confusa. Intentaremos
esclarecer en capítulos posteriores el significado de términos como keypoint,
landmark, feature, corner, etc. que algunas veces pueden parecer
intercambiables en la literatura de visión por computadora.
} en las imágenes de las cámaras. En el mejor de los casos el mapa
formado en SLAM y las actualizaciones que este sufre pueden durar la
corrida entera o incluso ser guardados en almacenamiento persistente, y
esto ayuda al dispositivo a entender cuando está viendo un lugar que ya
vio anteriormente. Sin embargo, hay soluciones más sencillas que solo
mantienen el mapa por una corta ventana de tiempo. A estos sistemas se
los suele denominar de \emph{odometría visual-inercial (VIO)}, y suelen
presentar mejor desempeño que soluciones completas de SLAM a costa de la
precisión en la trayectoria estimada.

\hypertarget{integraciuxf3n-en-monado}{%
\section{Integración en Monado}\label{integraciuxf3n-en-monado}}

\emph{Monado}\footnote{\url{https://monado.dev}} es una implementación
de código libre del estándar abierto \emph{OpenXR}
\autocite{thekhronosgroupinc.OpenXRSpecification} que intenta unificar
la forma en la que las aplicaciones interactúan con el hardware XR. Si
bien entenderemos mejor estos términos más adelante en la
\Cref{sec:thesis-context}, basta con saber que OpenXR es una
especificación desarrollada por el \emph{Khronos Group}\footnote{\url{https://www.khronos.org}},
un consorcio abierto sin fines de lucro integrado por distintos
participantes de la industria. Monado fue desarrollado y es mantenido
por \emph{Collabora}\footnote{\url{https://www.collabora.com}}, una
consultora especializada en proyectos de código abierto, la cual además
participa activamente en la especificación de OpenXR desde sus inicios.
Este trabajo fue realizado en el marco de una pasantía con Collabora.

Monado, además de implementar la especificación de OpenXR, provee varias
herramientas y funcionalidades para XR, además de controladores para
dispositivos populares. El proyecto, antes de este trabajo, contaba con
un puñado de métodos de tracking, pero la localización visual-inercial
por SLAM o VIO era una característica faltante.

Este trabajo, entonces, resultó en la integración de tres sistemas de
SLAM/VIO de código libre con Monado. En particular, los sistemas
integrados fueron Kimera-VIO
\autocite{rosinolKimeraOpenSourceLibrary2020}, ORB-SLAM3
\autocite{camposORBSLAM3AccurateOpenSource2021} y Basalt
\autocite{usenkoBasaltVisualInertialMapping2020}. Además, se adaptaron
controladores para poder utilizar dispositivos con este tipo de tracking
en Monado, incluyendo cascos de la plataforma Windows Mixed Reality.
Estos, según el mejor entendimiento del autor, son ahora los primeros
cascos comerciales capaces de ser localizados por SLAM/VIO en una
plataforma basada completamente en código libre. Se puede ver un video
demostrando el tracking funcionando en la URL referenciada al pie de
página \footnote{\url{https://youtu.be/g1o2xADr5Fw}}.

\hypertarget{estructura-de-la-tesis}{%
\section{Estructura de la tesis}\label{estructura-de-la-tesis}}

En la primera parte de este trabajo veremos algunos conceptos
fundamentales para entender el funcionamiento de los sistemas de SLAM y
VIO. Veremos representaciones usuales de la pose de un cuerpo en el
espacio y de las transformaciones que se le pueden aplicar al mismo.
Esto incluye conceptos como \emph{ángulos Euler}, \emph{cuaterniones} y
algunas ideas básicas de \emph{grupos de Lie}. En el capítulo siguiente
desarrollaremos el método de \emph{optimización no lineal} de
\emph{Gauss Newton}. Este tipo de optimización es el núcleo
computacional de los sistemas modernos, ya que es utilizada para
realizar la fusión de muestras de sensores maximizando la verosimilitud
de un estimador estadístico. Además, se utiliza este tipo de
optimización para resolver varios otros subproblemas como lo son la
calibración de sensores, la proyección y reproyección de puntos hacia
las cámaras y viceversa, la creación y el mantenimiento del mapa del
entorno, entre otros.

La siguiente parte del trabajo nos enfocaremos en una implementación
particular de uno de estos sistemas: Basalt. Esto permitirá introducir y
ver muchos de los conceptos y algoritmia que se utilizan de forma
concreta y contextualizadas. Veremos cómo Basalt utiliza algoritmos de
\emph{optical flow} para hacer el seguimiento de \emph{features}, cómo
decide cuáles cuadros serán \emph{keyframes}, de qué forma utiliza y
\emph{preintegra} las muestras de la IMU, la gestión de las
\emph{landmarks} que realiza y finalmente cómo construye la función de
\emph{error} a optimizar.

En la tercera parte presentaremos los detalles de la integración con
Monado. Veremos las decisiones tomadas a la hora de diseñar la interfaz
fundamental de la integración con este tipo de sistemas. Hablaremos
algunos problemas particulares que debieron resolverse para aplicarlos a
XR. Además hablaremos de los controladores de dispositivos que se
extendieron y los cuidados que hubo que tener para generar datos
aceptables para SLAM.

Para cerrar el trabajo, en la cuarta parte presentamos algunos
resultados cualitativos del rendimiento y la precisión de los sistemas
integrados. Finalmente cerramos con algunas conclusiones y líneas de
trabajo a considerar para el futuro.

\hypertarget{fundamentos}{%
\chapter{Fundamentos}\label{fundamentos}}

\hypertarget{transformaciones}{%
\section{Transformaciones}\label{transformaciones}}

Usualmente necesitaremos manipular y referirnos a transformaciones
tridimensionales entre distintos sistemas de referencias como se muestra
en la \figref{fig:transforms}. Estos sistemas surgen de las entidades
que forman parte de nuestro proceso de optimización en SLAM. Algunos
ejemplos de transformaciones en los que podríamos estar interesados son
las que describen que transformación es necesaria realizar sobre la
cámara izquierda para llegar a la cámara derecha de nuestro dispositivo,
o la transformación que le ocurrió al agente localizado entre el
instante anterior y el actual.

\fig{fig:transforms}{source/figures/transforms.pdf}{Transformaciones}{%
Intuición de lo que representa una transformación $T$ que rota y traslada el
sistema de referencia $A$ hacia $B$. En general estas transformaciones tendrán
una inversa $T^{-1}$ que puede deshacer la acción original. Notar que un
punto en el espacio $p \in \R^3$ tiene diferentes coordenadas dependiendo de que
sistema se tome como referencia.
}

El tipo de transformación en el que estamos interesados son los
\emph{movimientos de cuerpo rígido}\marginnote{%\
No profundizaremos en la definición formal de este concepto, pero intuitivamente,
son transformaciones que al aplicarlas sobre un conjunto de puntos, preservan su volumen.
Esto implica que preservan la norma y el producto cruz.
}. Estos pueden describirse mediante \emph{traslaciones} y
\emph{rotaciones}. Además, también nos gustaría poder expresar la
\emph{ubicación} y \emph{orientación} de las entidades en nuestro
sistema. Estas dos características conforman la \emph{pose} en el
espacio de tal entidad. Es posible describir este concepto como una
transformación aplicada sobre un sistema de referencia global fijo. Por
ejemplo en la \figref{fig:transforms} si pensamos en \(A\) como este
origen global, tenemos entonces que \(T\) está describiendo la pose de
\(B\) con respecto a \(A\). A continuación, desarrollaremos algunas
definiciones que nos ayudarán a describir la idea de transformación más
formalmente, y utilizaremos este mismo concepto para identificar las
poses de nuestras entidades.

\hypertarget{sec:linearalg-prelim}{%
\subsection{Preliminares del álgebra
lineal}\label{sec:linearalg-prelim}}

Comenzaremos construyendo sobre algunas ideas básicas del álgebra
lineal.

\begin{definition}Una transformación lineal \(L\) entre dos espacios
vectoriales \(V, W\) es una función \(L : V \rightarrow W\) tal que:
\bigbreak

\begin{itemize}
  \item $L(x + y) = L(x) + L(y) \quad \forall x, y \in V$
  \item $L(ax) = aL(x) \quad \forall a \in \R$
\end{itemize}

\end{definition}

\begin{property}[Matriz de una transformación]Si \(V \subseteq \R^n\)
tiene base canónica \(e_1, ..., e_n\) con \(e_i \in \R^m\) tenemos que
la matriz \(A = [L(e_1), ..., L(e_n)] \in \R^{m \times n}\) cumple
\(Ax = L(x) \quad \forall x \in V\)\end{property}

\begin{property}[Matrices cuadradas]El conjunto de matrices en
\(\R^{n\times n}\), con las operaciones de adición y multiplicación
matricial usuales, forman un anillo\footnote{\url{https://en.wikipedia.org/wiki/Ring_(mathematics)\#Definition}}
sobre el cuerpo de los reales. En particular, es un conjunto que se
encuentra cerrado bajo estas operaciones.\end{property}

\begin{property}\label{rmk:sqnormofsum}Sean \(a, b \in \R^n\), tenemos
que \(\| a + b \| ^ 2 = \|a\|^2 + \|b\|^2 + 2 a^T b\).\end{property}

\begin{proof}Desarrollemos: \begin{align}
\| a + b \| ^ 2 &= (a_1 + b_1)^2 + \dots + (a_n + b_n)^2 \\
&= (a_1^2 + b_1^2 + 2 a_1 b_1) + ... + (a_n^2 + b_n^2 + 2 a_n b_n) \\
&= \| a \| ^2 + \| b \|^2 + 2 \langle a, b \rangle \\
&= \| a \| ^2 + \| b \|^2 + 2 a^T b
\end{align}\end{proof}

\hypertarget{representaciuxf3n-de-rotaciones}{%
\subsection{Representación de
rotaciones}\label{representaciuxf3n-de-rotaciones}}

Para las rotaciones, y orientaciones, existen múltiples representaciones
válidas, cada una con sus ventajas y desventajas. Veremos algunas que
han sido utilizadas en este trabajo.

\hypertarget{uxe1ngulos-euler}{%
\subsubsection{Ángulos Euler}\label{uxe1ngulos-euler}}

La representación por ángulos Euler utiliza un vector
\([x, y, z]^T \in \R^3\) en donde cada componente representa el ángulo
de rotación que aplicar alrededor de tres ejes seleccionados por
convención. Es decir, el vector describe tres rotaciones que aplicar.
Esta representación tiene la principal ventaja de resultar intuitiva
cuando solo se necesita describir una rotación, pero se vuelve
inconveniente en otros contextos que necesitaremos.

Uno de sus problemas es que existen docenas de convenciones válidas que
varían en: la elección de los tres ejes, orden de los mismos, aplicación
global o local en cada uno. Aunque en la práctica solo se utiliza un
puñado de ellas, se puede volver fácilmente un punto de
confusión\marginnote{%\
En la práctica, la manipulación de rotaciones y sistemas de coordenadas son
fuentes usuales de muchos dolores de cabeza. Usualmente producto del uso
no documentado de distintas convenciones.
}. El otro problema fundamental es el llamado \emph{gimbal
lock}\footnote{\url{https://en.wikipedia.org/wiki/Gimbal_lock}} en donde
la combinación de ciertas rotaciones puede causar la pérdida de un grado
de libertad y describir nuevas rotaciones arbitrarias se hace imposible
a partir de ese punto. Otros problemas de los ángulos Euler se detallan
en \textcite{shoemakeAnimatingRotationQuaternion1985}.

En general, evitaremos esta representación.

\hypertarget{uxe1ngulo-axial}{%
\subsubsection{Ángulo axial}\label{uxe1ngulo-axial}}

Toda rotación puede representarse con un ángulo \(\omega\) y un eje
axial, representado por un vector unitario \(a \in \R^3\), sobre el cual
rotar \(\omega\) radianes. La representación de ángulo axial de esta
rotación queda definida por el vector \(v = \omega a\). Tenemos entonces
que \(a = \frac{v}{\| v \|}\) y \(\omega = \| v\|\).

Esta representación surge naturalmente del uso de giroscopios. Estos
módulos, internamente se componen de tres sensores capaces de reportar
cada uno la velocidad angular sobre un eje y se los ubica sobre tres
ejes ortogonales. De esta forma, la velocidad que reportan queda
determinada por la cantidad de rotación que ocurrió en cada eje durante
el instante de la muestra capturada. Si se ven los valores de los tres
sensores como componentes de un vector tridimensional, este coincide con
la representación angular axial \(v\) de la rotación capturada.

Esta representación se utilizará para el almacenamiento de valores de
velocidad angular provistos por muestras de giroscopio, ya que no es
necesaria ningún tipo de conversión. Además, veremos más abajo que esta
forma de describir rotaciones surge naturalmente en el estudio de otras
representaciones.

\hypertarget{cuaterniones-unitarios}{%
\subsubsection{Cuaterniones unitarios}\label{cuaterniones-unitarios}}

Los cuaterniones están definidos sobre un álgebra \(\H\). Son una
construcción que extiende a \(\R\) con tres números imaginarios \(i\),
\(j\) y \(k\) que satisfacen las siguientes propiedades: \begin{align}
i^2 = j^2 = k^2 = -1 \label{eq:quat-imaginary-props-start} \\
i = jk, \quad -i = kj \\
k = ij, \quad -k = ji \\
j = ki, \quad -j = ik \label{eq:quat-imaginary-props-end}
\end{align}

Un cuaternión \(q \in \H\) tiene la forma \begin{align}
q = q_w + q_x i + q_y j + q_z k \quad \text{con }\  q_w, q_x, q_y, q_z \in \R
\end{align}

El producto y la adición se extienden de \(\R\) naturalmente incluyendo
las propiedades de las partes imaginarias listadas en
\Crefrange{eq:quat-imaginary-props-start}{eq:quat-imaginary-props-end}

Nos interesarán tres operadores de los cuaterniones:

\bigbreak

\begin{itemize}
\item
  Conjugado: \(q^* = q_w - (q_x i + q_y j + q_z k)\)
\item
  Norma:
  \(\| q \| = \sqrt{q_w ^2 + q_x^2 + q_y^2 + q_z^2} = \sqrt{qq^*}\)
\item
  Inverso: \(q^{-1} = \frac{q^*}{\|q\|^2}\)
\end{itemize}

\bigbreak

Representaremos rotaciones y orientaciones con \textbf{cuaterniones
unitarios}, es decir cuaterniones \(q\) tal que \(\| q \| = 1\). Más aún
cualquier cuaternión \(q'\) puede hacerse unitario dividiéndolo por su
norma, o sea \(q = q' / \|q'\|\) tiene norma 1.

Una rotación representada en \textbf{ángulos axiales} por el vector
unitario \(a = [x, y, z]^T\) y ángulo \(\omega\) se representa con el
siguiente cuaternión unitario \(q\): \begin{align}
q = cos\frac{\omega}{2} + x\ sin\frac{\omega}{2}\ i + y\ sin\frac{\omega}{2}\ j + z\ sin\frac{\omega}{2}\ k
\end{align} Más aún, todo cuaternión unitario puede expresarse de esa
forma.

Para \textbf{rotar un vector} \(v = [x, y, z]^T \in \R^3\) con una
rotación representada por el cuaternión \(q\) basta con definir el
cuaternión \(p = 0 + xi + yj + zk\) y computar \(p'\) de la siguiente
manera: \begin{align}
p' = qpq^{-1} = 0 + p'_x i + p'_y j + p'_z k
\end{align}

El cuaternión resultante \(p'\) tendrá parte real nula y el vector
\(v' = [p'_x, p'_y, p'_z]^T \in \R^3\) es el vector \(v\) rotado por
\(q\).

La \textbf{composición de rotaciones} queda bien definida por la
multiplicación de cuaterniones. Es decir, dados los cuaterniones \(q\) y
\(p\), el cuaternión \(pq\) describe la rotación que se produce al
aplicar primero \(q\) y luego \(p\).

La \textbf{rotación identidad} o neutra, una rotación que ``no rota'',
coincide con \(1 \in \R\), es decir no tiene parte imaginaria, es
\(1 + 0i + 0j + 0k\).

El \textbf{inverso multiplicativo} de un cuaternión \(q^{-1}\) al
componerlo con \(q\), como es de esperarse, produce la identidad, o sea:
\begin{align}
q^{-1}q = 1
\end{align}

Muchas veces necesitaremos computar el \textbf{delta de rotación} que
lleva de una orientación \(p\) a otra \(q\). Esto es similar a lo que
obtenemos cuando restamos vectores, así que sobrecargaremos el operador
de resta de cuaterniones según el contexto de la siguiente manera:
\begin{align}
q - p = p^{-1} q
\end{align}

Será también útil poder interpolar entre dos orientaciones \(p\) y \(q\)
con un factor \(t \in [0, 1]\) para conseguir orientaciones intermedias
de \(p\) a \(q\). La interpolación lineal \(lerp\) genera cuaterniones
no unitarios, y por esto utilizaremos en su lugar la operación de
\textbf{interpolación esférica} \(slerp\) presentada en
\textcite{shoemakeAnimatingRotationQuaternion1985} y definida de la
siguiente manera: \begin{align}
slerp(p, q, t) = p(p^{-1}q)^t \label{eq:slerp-def}
\end{align}

\hypertarget{sec:rotation-matrices}{%
\subsubsection{Matrices de rotación}\label{sec:rotation-matrices}}

Las rotaciones pueden representarse también como una matriz
\(R \in \R^{3x3}\) con ciertas restricciones.

\begin{definition}Sean \(v, w \in \R^3\), decimos que \(v\) y \(w\) son
ortonormales si \(\norm{v} = \norm{w} = 1\) y son ortogonales, es decir
\(\dotprod{v}{w} = 0\)\end{definition}

Sean \(R^{(1)}, R^{(2)}, R^{(3)} \in \R^3\) las columnas de \(R\).

\begin{definition}\(R\) se dice ortogonal u ortonormal si sus columnas
son ortonormales de a pares, es decir,
\(\dotprod{R^{(i)}}{R^{(j)}} = 0\) y \(\norm{R^{(i)}} = 1\) para
\(i \neq j; \; i, j \in \{1, 2, 3\}\).\end{definition}

Las matrices ortonormales pueden tener determinante \(\pm 1\).
Llamaremos \textbf{matrices de rotación} solo a las \(R\) tal que
\(det(R) = +1\) \marginnote{%\
El hecho de que el determinante sea $+1$, intuitivamente, quiere decir que la
transformación lineal de $R$ no escala los vectores en $\R^3$,
o sea preserva volúmenes.
} \marginnote{%\
Las transformaciones representadas por $R$ con $det(R) = -1$ suelen llamarse
rotaciones impropias o rotorreflexiones, ya que además de rotar, permiten la reflexión
de vectores en $\R^3$.
}. El conjunto de estas matrices se denomina \(SO(3)\) y veremos en la
sección siguiente el por qué de este nombre. Estas matrices presentan
propiedades agradables para ser manipuladas como rotaciones en el
álgebra matricial usual:

\bigbreak

\begin{itemize}
\item
  \(R\) es la transformación lineal que rota vectores \(v \in \R^3\), es
  decir \(Rv\) es el vector \(v\) rotado por la rotación representada en
  \(R\).
\item
  La composición de rotaciones \(R, S \in SO(3)\) queda representada con
  la multiplicación de matrices usual \(RS\).
\item
  La inversa de \(R\) es \(R^{-1} = R^T\), más aún esta propiedad define
  a las matrices ortogonales.
\end{itemize}

\bigbreak

\begin{proof}Tenemos por la ortonormalidad de las columnas de \(R\) y
definición de \(R^T\) que \begin{align}
  & R^T R = I_{3 \times 3} \\
  & \Leftrightarrow \begin{bmatrix}
  \dotprod{R^{(1)}}{R^{(1)}} & \dotprod{R^{(1)}}{R^{(2)}} & \dotprod{R^{(1)}}{R^{(3)}} \\
  \dotprod{R^{(2)}}{R^{(1)}} & \dotprod{R^{(2)}}{R^{(2)}} & \dotprod{R^{(2)}}{R^{(3)}} \\
  \dotprod{R^{(3)}}{R^{(1)}} & \dotprod{R^{(3)}}{R^{(2)}} & \dotprod{R^{(3)}}{R^{(3)}} \\
  \end{bmatrix} = \begin{bmatrix}
1 & 0 & 0 \\
0 & 1 & 0 \\
0 & 0 & 1 \\
  \end{bmatrix} \\
  & \Leftrightarrow \| R^{(i)} \| = 1 \text{ y } \dotprod{R^{(i)}}{R^{(j)}} = 0 \quad \forall i \neq j \in \{1,2,3\} \\
  & \Leftrightarrow \text{R es ortogonal}
\end{align}\end{proof}

\bigbreak

\hypertarget{representaciuxf3n-de-transformaciones}{%
\subsection{Representación de
transformaciones}\label{representaciuxf3n-de-transformaciones}}

Tener rotaciones expresadas como matrices, o equivalentemente como
transformaciones lineales, prueba ser muy útil en la práctica, ya que
podemos recurrir al arsenal de funcionalidad preexistente para estas
estructuras. A continuación, desarrollaremos una serie de definiciones
que nos permitirán extender esta idea de utilizar matrices para cubrir
el concepto de transformación en su totalidad, con traslaciones
incluidas.

\hypertarget{grupos}{%
\subsubsection{Grupos}\label{grupos}}

Vimos en la \Cref{sec:linearalg-prelim} que \(\Rnn\) forma un anillo con
multiplicación y suma cerradas bajo \(\R\). Además, mostramos que las
matrices en \(\Rnn\) corresponden a transformaciones lineales. Veremos
ahora una serie de entidades que restringen a las transformaciones en
\(\Rnn\) con el objetivo de llegar a describir las transformaciones de
cuerpo rígido sobre \(\R^3\) que nos interesan.

\begin{definition}Un grupo es un conjunto \(G\) con una operación
binaria \newline \(\circ : G \times G \rightarrow G\) tal que
\(\forall a, b, c \in G\): \bigbreak

\begin{enumerate}
  \item Es cerrada: $a \circ b \in G$
  \item Es asociativa: $(a \circ b) \circ c = a \circ (b \circ c)$
  \item Tiene neutro: $\exists!\ e \in G: e \circ a = a \circ e = a$
  \item Tiene inverso: $\exists a^{-1} \in G: a \circ a^{-1} = a^{-1} \circ a = e$
\end{enumerate}

\end{definition}

Es posible empezar a intuir que \emph{tanto las rotaciones como las
transformaciones} que describimos al principio del capítulo
\emph{coinciden con esta idea de grupo}. ``Mezclar'' transformaciones
debería dar como resultado otra transformación (cerrada) y no debería
importar el orden de mezcla (asociativa). Debería existir una
transformación neutra que no haga nada, y una transformación inversa que
deshagan la transformación original. Si además pudiésemos describir
estas transformaciones con matrices, seríamos capaces de recurrir a
todas las herramientas preexistentes. Veamos como hacer eso ahora.

\begin{definition}El conjunto de matrices invertibles en \(\Rnn\) con la
operación de multiplicación matricial forman un grupo denominado Grupo
Lineal General \(GL(n)\). Es decir: \begin{align}
  GL(n) = \{ A \in \Rnn : det(A) \neq 0 \}
\end{align}\end{definition}

\begin{definition}El subconjunto de matrices de GL(n) con determinante 1
se denomina el Grupo Lineal Especial \(SL(n)\). Es decir: \begin{align}
  SL(n) = \{ A \in GL(n) : det(A) = 1 \}
\end{align}\end{definition}

\begin{remark}\(A \in SL(n) \Rightarrow A^{-1} \in SL(n)\) ya que
\(det(A^{-1}) = \frac{1}{det(A)}\).\end{remark}

\begin{definition}Un grupo \(G\) tiene una representación matricial si
existe un homomorfismo inyectivo \(F:G \rightarrow GL(n)\). Es decir,
una función inyectiva para la cual la multiplicación matricial preserva
la estructura del operador \(\circ\). En símbolos,
\(\forall a, b \in G\): \begin{align}
  F(e) &= I_{n \times n} \\
  F(a \circ b) &= F(a)F(b)
\end{align}\end{definition}

\begin{definition}El subconjunto de matrices ortogonales de \(GL(n)\) se
denomina el Grupo Ortogonal \(O(n)\). Es decir: \begin{align}
O(n) = \{ R \in GL(n) : R^T R = I \}
\end{align}\end{definition}

\begin{property}Una matriz ortogonal \(R \in \Rnn\) preserva el producto
interno.\end{property}

\begin{proof}\(\forall x, y \in \R^{n}\) tenemos: \begin{align}
\dotprod{Rx}{Ry} = (Rx)^T Ry = x^T R^T R y = x^T y = \dotprod{x}{y}
\end{align}\end{proof}

\begin{property}Una matriz ortogonal \(R \in O(n)\) tiene
\(det(R) = \pm 1\)\end{property}

\begin{proof}\begin{align}
& 1 = det(I) = det(R^T R) = det(R^T) det(R) = det(R)^2 \\
& \Leftrightarrow det(R) = \pm 1
\end{align}\end{proof}

\begin{definition}El subconjunto de matrices ortogonales de \(O(n)\) con
determinante \(+1\) se denomina el Grupo Ortogonal Especial \(SO(n)\).
Es decir: \begin{align}
  SO(n) = \{ R \in O(n) : det(R) = +1 \} = O(n) \cap SL(n)
\end{align}\end{definition}

Como vimos en la \Cref{sec:rotation-matrices}, este grupo es el que
define a las \textbf{matrices de rotación}. Y se corresponde a la
representación matricial que presentamos allí. Continuemos ahora con las
transformaciones.

\hypertarget{transformaciones-como-grupos}{%
\subsubsection{Transformaciones como
grupos}\label{transformaciones-como-grupos}}

\begin{definition}Una transformación lineal afín
\(L : \R^n \rightarrow \R^n\) es tal que existe \(B \in GL(n)\) y
\(b \in \R^n\) que determinan \(L(x) = Bx + b\).\end{definition}

\begin{definition}El grupo de tales transformaciones se denomina Grupo
Afín de dimensión \(n\) \(A(n)\).\end{definition}

En general, una transformación afín \(L(x) = Bx + b \in A(n)\) no es una
transformación lineal a menos que \(b = 0\). Introduciremos las llamadas
\emph{coordenadas homogéneas} para expresar transformaciones afines en
\(A(n)\) como transformaciones lineales en \(GL(n + 1)\). Extenderemos
\(L\) de la siguiente manera: \begin{align}
L': \R^{n+1} &\rightarrow \R^{n+1} \\
L' \begin{pmatrix}\begin{bmatrix}x \\ 1\end{bmatrix}\end{pmatrix} &= \begin{bmatrix}
B & b \\
0 & 1
\end{bmatrix}
\begin{bmatrix}x \\ 1\end{bmatrix}
= \begin{bmatrix} Bx + b \\ 1 \end{bmatrix}
= \begin{bmatrix} L(x) \\ 1 \end{bmatrix}
\end{align}

Notar que hay un isomorfismo entre \(L\) y \(L'\), ya que \(0\) y \(1\)
son constantes y por esto diremos que son la misma transformación. La
matriz \(\begin{bmatrix} B & b \\ 0 & 1 \end{bmatrix}\) se dice una
matriz afín y pertenece a \(GL(n + 1)\). Tenemos entonces que las
matrices afines forman un subgrupo de \(GL(n + 1)\)

\begin{definition}Una transformación euclídea
\(L: \R^n \rightarrow \R^n\) se define con una matriz ortogonal
\(R \in O(n)\) y un vector \(t \in \R^n\) como
\(L(x) = Rx + t\).\end{definition}

\begin{definition}El grupo de tales transformaciones se denomina Grupo
Euclídeo \(E(n)\) y es un subgrupo de \(A(n)\). Es decir \begin{align}
  E(n) = \begin{Bmatrix} \begin{bmatrix} R & t \\ 0 & 1 \end{bmatrix} \in
  \R^{(n+1) \times (n+1)} : R \in O(n), t \in \R^n \end{Bmatrix}
\end{align}\end{definition}

\begin{definition}El subconjunto de transformaciones euclídeas con
\(R \in SO(n)\) se denomina el Grupo Euclídeo Especial \(SE(n)\), es
decir: \begin{align}
  E(n) = \begin{Bmatrix} \begin{bmatrix} R & t \\ 0 & 1 \end{bmatrix} \in
  \R^{(n+1) \times (n+1)} : R \in SO(n), t \in \R^n  \end{Bmatrix}
\end{align}\end{definition}

Y es este el grupo que buscábamos, ya que \(SE(3)\) es capaz de
representar las transformaciones de cuerpo rígido en \(\R^3\) que
necesitábamos. Tenemos entonces que representaremos rotaciones con
\(R \in SO(3)\), osea matrices cuadradas \(3 \times 3\); y
representaremos transformaciones con \(T \in SE(3)\), osea matrices
cuadradas \(4 \times 4\). Estos grupos también funcionan con dos
dimensiones en \(\R^2\) de la misma manera para \(SO(2)\) y \(SE(2)\).
Finalmente, se pueden visualizar los conjuntos definidos en esta sección
en la \figref{fig:hasse-groups}.

\fig{fig:hasse-groups}{source/figures/hasse-groups.pdf}{Árbol de grupos}{%
Diagramas de Hasse de los grupos desarrollados con el orden parcial
$\subset$ (\italic{``es subconjunto propio de''}).
}

\hypertarget{representaciuxf3n-de-infinitesimales}{%
\subsection{Representación de
infinitesimales}\label{representaciuxf3n-de-infinitesimales}}

Ahora nos gustaría entender como tratar cambios infinitesimales en las
rotaciones de \(SO(3)\) y las transformaciones de \(SE(3)\). Esto será
necesario a la hora de aplicar algoritmos de optimización que dependen
de las derivadas de estas operaciones.

\hypertarget{matrices-antisimuxe9tricas}{%
\subsubsection{Matrices
antisimétricas}\label{matrices-antisimuxe9tricas}}

\begin{definition}[Producto cruz]Dados \(v, w \in \R^3\) el producto
cruz de \(v\) y \(w\) es un vector ortogonal a ambos tal que:
\begin{align}
v \times w = \begin{bmatrix}
v_y w_z - v_z w_y \\
v_z w_x - v_x w_z \\
v_x w_y - v_y w_x
\end{bmatrix} \in \R^3
\end{align}\end{definition}

\begin{definition}[Matriz antismétrica]\(R \in \RR3\) se dice
antisimétrica si \(R = -R^T\)\end{definition}

\begin{definition}Definimos el operador
\(\hat{\cdot} : \R^3 \rightarrow \RR3\) que devuelve la siguiente matriz
antisimétrica: \begin{align}
\hat{v} = \begin{bmatrix}
0 & -v_z & v_y \\
v_z & 0 & -v_x \\
-v_y & v_x & 0
\end{bmatrix} \in \RR3
\end{align}\end{definition}

La construcción del operador \(\hat{v}\) está diseñada para tener la
siguiente propiedad:

\begin{property}\(\hat{v} w = v \times w\)\end{property}

Además, por la definición de matriz antisimétrica es sencillo ver que:

\begin{property}\label{prop:skew-mat-vec}Toda matriz antisimétrica \(R\)
está unívocamente definida por un vector \(v \in \R^3\) tal que
\(R = \hat{v}\)\end{property}

\begin{definition}El conjunto de todas las matrices antisimétricas en
\(\R^3\) se denomina Álgebra de Lie Ortogonal Especial \(\so3\), es
decir: \begin{align}
\so3 = \{ \hat{v} \in \RR3 : v \in \R^3 \}
\end{align}\end{definition}

Veremos a continuación cuál es su relación con \(SO(3)\) y el por qué de
su nombre.

\hypertarget{rotaciones-infinitesimales}{%
\subsubsection{Rotaciones
infinitesimales}\label{rotaciones-infinitesimales}}

Consideremos una familia de rotaciones \(R(t) \in SO(3)\) con
\(t \in \R\) que describen una rotación continua aplicada sobre un punto
\(X(0) \in \R^3\) hacia otro \(X(t)\), es decir: \begin{align}
  R(0) = I \\
  X(t) = R(t)X(0)
\end{align}

\begin{notation}El parámetro \(t\) lo consideraremos implícito en
algunas ocasiones, o sea \(R = R(t)\). Además, utilizaremos la notación
\(\dot{R} = \frac{dR}{dt}\).\end{notation}

Como \(RR^T = I\) tenemos que \begin{align}
& 0 = \frac{d}{dt}I = \frac{d}{dt}(RR^T) = \dot{R}R^T + R \dot{R}^T \\
& \Rightarrow \dot{R}R^T = -R \dot{R}^T
\end{align}

O sea que \(\dot{R}R^T \in \so3\). Por la \Cref{prop:skew-mat-vec}
tenemos que existe un vector único \(v(t) \in \R^3\) tal que
\begin{align}
\dot{R}(t) R(t)^T = \hat{v}(t) \Leftrightarrow \dot{R}(t) = \hat{v}(t) R(t)
\end{align} Y como \(R(0) = I\), entonces \(\dot{R}(0) = \hat{v}(0)\).
Por lo tanto, tenemos que la matriz antisimétrica \(\hat{v}(0)\) nos da
la aproximacíon de primer orden de una rotación: \begin{align}
R(dt) = R(0) + (dR)(0) = I + \hat{v}(0)dt \label{eq:hatop-is-rotvel}
\end{align}

Notar que si pensamos a \(t\) en términos de tiempo, la
\Cref{eq:hatop-is-rotvel} deja ver a \(\hat{v}\) como una matriz que
describe la velocidad de la rotación.

\clearpage

Tenemos entonces que el efecto de una rotación infinitesimal en
\(SO(3)\) puede ser aproximado por matrices en \(\so3\). Es necesario
mencionar que \(SO(3)\) es lo que se denomina un \emph{grupo de Lie}
\marginnote{%\
Un grupo de Lie es una ``variedad diferenciable'' en el cual la operación del grupo,
y su inversa, también son diferenciables. Intuitivamente, esto significa que la
aplicación de rotaciones o transformaciones es ``suave'' o continua y esto nos
permite trabajar con el concepto de límite y derivadas.
} mientras que \(\so3\) es su correspondiente \emph{álgebra de Lie}
\marginnote{%\
Un grupo de Lie tiene un álgebra de Lie relacionada. Esta última es el
``espacio tangente'' en la identidad $I$ del grupo (en particular, de la variedad).
Intuitivamente, este puede pensarse como el espacio de todas las posibles
velocidades alrededor de $I$.
Esto es precisamente lo que desarrollamos al calcular $R(dt)$ en la \Cref{eq:hatop-is-rotvel}.
}. No necesitaremos adentrarnos en estos conceptos, pero es común
encontrarlos en la literatura y gran parte del desarrollo que estamos
realizando está ligado a ellos.

Luego de haber presentado estas ideas fundamentales, a partir de aquí
procederemos a dar una \emph{vista general} de la definición de dos
operadores, \(exp\) y \(log\) que nos permiten pasar del grupo al
álgebra de Lie y viceversa. Más adelante, también veremos rápidamente
como estos conceptos se traducen de forma muy similar para las
transformaciones en \(\R^3\). Al final del capítulo listaremos algunas
referencias para el lector que quiera profundizar en los conceptos y las
derivaciones.

\bigbreak

\begin{center}\rule{0.5\linewidth}{0.5pt}\end{center}

\bigbreak

Vimos en el desarrollo anterior que existe un vector
\(\hat{v}(t) \in \so3\) que actúa como un término de velocidad
rotacional. Si asumimos velocidad constante, es decir \(\hat{v}\)
constante, nos interesaría saber cual es la rotación total luego de
rotar desde \(R(0) = I\) con esta velocidad durante un tiempo \(t\). En
otras palabras, queremos computar \(R(t)\) dado \(\hat{v}\). Teniendo en
cuenta el desarrollo anterior basta con plantear el siguiente sistema de
ecuaciones diferenciales: \begin{align} \begin{cases}
\dot{R}(t) = \hat{v} R(t) \\
R(0) = I
\end{cases}
\end{align}

Es posible ver que este sistema tiene como solución la siguiente
expresión. \begin{align}
R(t) = exp(\hat{v}t) = \sum_{n=0}^{\infty}\frac{(\hat{v}t)^n}{n!} \label{eq:exp-def}
\end{align}

Notar que los exponentes de la \Cref{eq:exp-def} son respecto a la
multiplicación matricial.

Esta rotación se corresponde con, dado \(w = v t \in \R^3\), la rotación
de \(\omega = \norm{w}\) radianes alrededor del eje dado por el vector
unitario \(a = w / \norm{w}\). Dicho de otro modo, \emph{la
representación ángulo-axial \(w = \omega a\) puede ser convertida en la
matriz de rotación \(R\) equivalente mediante el operador \(exp\) con
\(R = exp(\hat{w})\)}. Más aún, \(\so3\) contiene todos estos vectores
ángulo-axiales.

El operador \(exp : \so3 \rightarrow SO(3)\) se denomina \emph{mapa o
aplicación exponencial} y su inversa es el \emph{mapa logarítmico}
\(log : SO(3) \rightarrow \so3\). Este, dado una matriz de rotación
\(R \in SO(3)\) devuelve \(\hat{w} \in \so3\) tal que
\(R = exp(\hat{w})\) computando \(w\) de la siguiente manera
\autocite{eadeDerivativeExponentialMap}: \begin{align}
  \norm{w} = cos^{-1} \left( \frac{traza(R) - 1}{2} \right) \label{eq:log-def-start}\\
  w = \frac{\norm{w}}{2sin(\norm{w})} \begin{bmatrix} R_{3,2} - R_{2,3} \\
  R_{1,3} - R_{3,1} \\ R_{2,1} - R_{1,2} \end{bmatrix} \label{eq:log-def-end}
\end{align}

Notar que definimos \(exp\) en la \Cref{eq:exp-def} como una suma
infinita mientras que \(log\) pudo definirse con una expresión cerrada
en las \Crefrange{eq:log-def-start}{eq:log-def-end}. Existe la
\emph{fórmula de Rodrigues} que permite expresar también \(exp\) de
forma cerrada: \begin{align}
exp(\hat{w}) = I + \frac{sin(\norm{w})}{\norm{w}} \hat{w} + \frac{1 - cos(\norm{w})}{\norm{w}^2} \hat{w}^2
\end{align}

\hypertarget{transformaciones-infinitesimales}{%
\subsubsection{Transformaciones
infinitesimales}\label{transformaciones-infinitesimales}}

Para las transformaciones en \(SE(3)\) tenemos un desarrollo muy similar
al de las rotaciones en \(SO(3)\). Consideremos una familia de
transformaciones \(T(t) \in SE(3)\) compuestas por rotaciones
\(R(t) \in SO(3)\) y traslaciones \(b(t) \in \R^3\) con \(t \in \R\) que
representan una transformación continua aplicada al punto
\(X(0) \in \R^3\) con \(T(0) = I\). Es decir, \begin{align}
X(t) &= T(t)X(0) \\
T(t) &= \begin{bmatrix}
R(t) & b(t) \\
0 & 1
\end{bmatrix}
\end{align}

Considerando que esta vez la inversa de \(T\) es \(T^{-1}\) y no su
transpuesta como era el caso con las rotaciones. Podemos aplicar un
desarrollo similar al de la sección anterior y llegar a que:
\begin{align} \label{eq:translation-deriv}
\dot{T}T^{-1} = \begin{bmatrix}
\dot{R}R^T & \dot{b} - \dot{R} R^T b \\
0 & 0
\end{bmatrix} \in \RR4
\end{align}

De vuelta, \(\dot{R}R^T\) corresponde a alguna matriz antisimétrica
\(\hat{v} \in \so3\). Definiendo un vector
\(y(t) = \dot{b}(t) - \hat{v}(t) b(t)\) podemos reescribir la
\Cref{eq:translation-deriv} e introducir el concepto de \emph{giro o
twist} \(\hat{\xi}(t)\): \begin{align}
\hat{\xi}(t) = \dot{T}(t)T^{-1}(t) = \begin{bmatrix}
\hat{v}(t) & y(t) \\
0 & 0
\end{bmatrix}
\end{align}

La matriz de giro \(\hat{\xi}\) pertenece al álgebra de Lie \(\se3\) del
grupo de Lie \(SE(3)\) y puede ser parametrizada por las coordenadas de
giro \(\xi \in \R^6\). Para esto se utiliza un operador \emph{hat}
\(\cdot^{\wedge}\) y su inversa \emph{vee} \(\cdot^{\vee}\) de la
siguiente manera: \begin{align}
\hat{\xi} &= {\begin{bmatrix}y \\ v\end{bmatrix}}^{\wedge} = \begin{bmatrix}\hat{v} &
y \\ 0 & 0 \end{bmatrix} \in \RR4 \\
{\begin{bmatrix}\hat{v} & y \\ 0 & 0 \end{bmatrix}}^{\vee} &= \begin{bmatrix}y \\
v\end{bmatrix} = \xi \in \R^6
\end{align}

Es decir, podemos codificar el cambio infinitesimal de una
transformación en un vector de giro \(\xi\) de seis dimensiones en
donde: \bigbreak

\begin{itemize}
\item
  Los últimos tres componentes dados por \(v\) son la representación
  \textbf{ángulo-axial} de la velocidad rotacional.
\item
  Los primeros tres componentes dados por \(y\) describen el cambio de
  traslación teniendo en cuenta esta velocidad rotacional instantánea
  \(\hat{v}\).
\end{itemize}

\bigbreak

Finalmente, tenemos el mapa exponencial y mapa logarítmico entre
\(SE(3)\) y \(\se3\). Similar a \(SO(3)\), cualquier transformación
\(T \in SE(3)\) va a poder ser representada por un vector de giro
\(\xi\) con \(T = exp(\hat{\xi})\).

Existen expresiones cerradas para ambos \(exp : \se3 \rightarrow SE(3)\)
y \(log : SE(3) \rightarrow \se3\).

Para \(exp(\hat{\xi})\) con \(\xi = [y, v]^T\) tenemos: \begin{align}
exp(\hat{\xi}) = \begin{bmatrix}
exp(\hat{v}) & Jy \\
0 & 1
\end{bmatrix}
\end{align}

Con \(J\) el llamado \emph{jacobiano izquierdo} de \(SO(3)\) que se
puede computar con el ángulo \(\omega\) de \(v\), es decir con
\(\omega = \norm{v}\) de esta manera
\autocite{eadeDerivativeExponentialMap}: \begin{align}
J = I + \frac{1-cos(\omega)}{\omega^2} \hat{v} + \frac{\omega -
sin(\omega)}{\omega^3} \hat{v}^2.
\end{align}

Para \(log(T)\) con
\(T = \begin{bmatrix}R & b \\ 0 & 1\end{bmatrix}\in SE(3)\) tenemos:
\begin{align}
log(T) = \begin{bmatrix}log(R)^{\vee} \\ J^{-1}b\end{bmatrix}^{\wedge}
\end{align}

Con el jacobiano inverso \(J^{-1}\) dado por
\autocite{eadeDerivativeExponentialMap}: \begin{align}
J^{-1} = I - \frac{1}{2} \hat{v} +
\left( \frac{1}{\omega^2} - \frac{1 + cos(\omega)}{2 \omega sin(\omega)} \right) \hat{v}^2
\end{align}

\hypertarget{cierre-y-literatura-recomendada}{%
\subsubsection{Cierre y literatura
recomendada}\label{cierre-y-literatura-recomendada}}

Tenemos ahora una mirada suficientemente formal como para ser capaces de
utilizar y comprender las herramientas usualmente utilizadas en sistemas
que modelan rotaciones y transformaciones en dos y tres dimensiones por
computadora. Nos tomamos el trabajo de entender varias formalizaciones
para las rotaciones, ya que todas ellas aparecen de una u otra forma en
el pipeline que se desarrolla en este trabajo. Las formalizaciones
infinitesimales que hemos presentado en la última subsección, y sus
representaciones en el álgebra de Lie, permiten además modelar los
estados de las entidades que un algoritmo de SLAM necesita describir de
una forma particularmente elegante. Esto es así porque permiten la
aplicación de algoritmos de optimización, como los que veremos en el
próximo capítulo, \emph{sin restricciones}\marginnote{%\
Representaciones de cuaterniones o matriciales requieren mantener los errores
numéricos a raya mediante la constante renormalización de sus valores.
Las representaciones dadas en las álgebras de Lie no tienen este problema.
}.

Respecto a la última sección de infinitesimales, se pueden encontrar
algunos de estos conceptos explorados en mayor profundidad en
\textcite{barfootStateEstimationRobotics2017} cap. 7. Una excelente
introducción a los grupos de Lie con una teoría reducida enfocada en
aplicaciones de robótica es presentada en
\textcite{solaMicroLieTheory2021}. Finalmente, gran parte de las
derivaciones de las expresiones que se vieron se encuentran detalladas
en \textcite{eadeDerivativeExponentialMap}.

\hypertarget{optimizaciuxf3n-por-cuadrados-muxednimos}{%
\section{Optimización por cuadrados
mínimos}\label{optimizaciuxf3n-por-cuadrados-muxednimos}}

Como se verá más adelante, muchos de los problemas fundamentales de SLAM
son problemas de optimización. Querremos minimizar una \emph{función de
error} o \emph{energía} no-lineal que integre las múltiples mediciones
que contienen información conflictiva, ruido, outliers, y demás
problemáticas propias de la toma de datos mediante sensores físicos. Se
utiliza como algoritmo principal para la optimización de estos sistemas
el algoritmo de Gauss-Newton, y es el que explicaremos en esta sección.
Cabe aclarar que existen algoritmos ligeramente más sofisticados como
\emph{Levenberg-Marquardt} o el \emph{método dogleg} que pueden
aplicarse; no nos explayaremos en ellos en este trabajo, pero serán
mencionados cuando sea pertinente.

\hypertarget{cuadrados-muxednimos-lineales}{%
\subsection{Cuadrados mínimos
lineales}\label{cuadrados-muxednimos-lineales}}

Comencemos con el caso lineal que será necesario para luego aproximar el
no-lineal.

Querremos encontrar algún punto \(x \in \R^n\) que cumpla una serie de
\(m\) restricciones \textbf{lineales} \(f_i(x) = b_i\) con
\(f_i \in \R^n \rightarrow \R\) combinación lineal de los componentes de
\(x\); \(\ b \in \R^m\) y con \(i = 1, ..., m\). Además definimos
\(f(x) = [f_1(x) \dots f_m(x)]^T\) y al ser \(f_i\) combinaciones
lineales de \(x\), tenemos que existe una transformación lineal
\(A \in \R^{m \times n}\) tal que \(f(x) = Ax\). De esta forma podemos
replantear nuestras ecuaciones como el sistema \(Ax = b\).

Será usual que \(m > n\) y que \(A\) tenga columnas linealmente
independientes así que solo consideraremos este caso. Al tener más
ecuaciones que incógnitas tenemos un sistema \emph{sobre-determinado} en
el cual, para valores usuales de \(b\), no tendremos solución.
Buscaremos entonces el \(x\) que ``mejor'' cumpla con estas
restricciones en algún sentido de la palabra. Utilizaremos el criterio
de minimizar los \emph{residuales} \(r_i(x) = f_i(x) - b_i\) y definimos
al vector residual como \(r(x) = [r_1(x) \dots r_m(x)]^T\). Querremos
encontrar un \(x\) que minimice la siguiente función de \emph{error} o
\emph{energía} \(E(x)\) basada en los residuales. \begin{align}
E(x) = \sum_{i=1}^m{r_i(x) ^ 2} = \| r(x) \| ^ 2
\end{align}

Notar que la definición de \(E(x)\) coincide con la de la \emph{norma
euclídea} al cuadrado de \(r(x)\). Encontrar \(x\) que minimice \(E(x)\)
con residuales lineales es el \emph{problema de los cuadrados mínimos
lineales}; a \(x\) le decimos \emph{solución} de tal problema. Notar que
\(r_i = (Ax - b)_i\), o sea el componente \(i\) del vector
\(Ax - b \in \R^m\), y por ende \(r(x) = Ax - b\). \begin{align}
E(x) = \sum_{i=1}^m{(Ax - b)_i^2} = \| Ax - b \| ^ 2
\end{align}

Lo que queremos minimizar es entonces equivalente a la norma cuadrada
del vector residual \(Ax - b\). Introduciremos a continuación una serie
de conceptos y teoremas necesarios para poder presentar una forma
sucinta de minimizar esta norma.

Aceptaremos el siguiente teorema sin demostración (ver
\textcite{shoresAppliedLinearAlgebra2007}, teoremas 2.7 y 3.7).

\begin{theorem}\label{thm:li1inv2nots3}Son equivalentes:

\begin{enumerate}
  \item Las columnas de $A$ son linealmente independientes.
  \item $A$ es invertible.
  \item $Ax = 0 \Leftrightarrow x = 0$.
\end{enumerate}

\end{theorem}

Otras definiciones y resultados que necesitaremos a continuación.

\begin{definition}\label{def:grammat}La matriz de Gram de \(A\) es
\(A^T A \in \R^{n \times n}\).\end{definition}

\begin{theorem}\label{thm:aligraminv}La matriz de Gram de \(A\) es
invertible si y solo si \(A\) tiene columnas linealmente
independientes.\end{theorem}

\begin{proof}Por \Cref{thm:li1inv2nots3} (3) \(A\) tiene columnas
linealmente independientes si y solo si
\(x = 0 \Leftrightarrow Ax = 0\), entonces nos basta con ver que
\(Ax = 0 \Leftrightarrow A^T A x = 0\). \bigbreak

\begin{itemize}
  \item $\Rightarrow)$ Si $Ax = 0$ entonces $A^T A x = A^T (A x) = A^T 0 = 0$
  \item $\Leftarrow)$ Si $A^T A x = 0$ entonces $0 = x^T A^T A x = (x^T A^T) A x =
(A x)^T (Ax) = \| Ax \|^2 = 0$ y esta norma es $0$ si y solo si $Ax = 0$
\end{itemize}

\end{proof}

\begin{definition}\label{def:pseudoinverse}Sea \(A\) con columnas
linealmente independientes. La matriz pseudo-inversa de \(A\) es
\(A^{\dagger} = (A^T A)^{-1} A^T\)\end{definition}

Notar que \((A^T A)^{-1}\) existe por \Cref{thm:aligraminv}.

Con estas herramientas, estamos en posición de presentar la solución
directa al problema de los cuadrados mínimos lineales.

\begin{theorem}\label{thm:linleastsquaresol}\(\hat{x} = A^{\dagger} b\)
es \textbf{la} solución del problema de los cuadrados mínimos lineales.
Es decir
\(\forall x \in \R^n: E(\hat{x}) = \| A\hat{x} - b \| ^ 2 \leq \| Ax - b \| ^ 2 = E(x)\).\end{theorem}

\begin{proof}Tenemos primero que \begin{align}
  & \hat{x} = A^\dagger b \\
  &\Leftrightarrow \hat{x} = (A^T A)^{-1} A^T b \\
  &\Leftrightarrow A^T A \hat{x} = A^T b \\
  &\Leftrightarrow A^T (A \hat{x} - b) = 0 \label{eq:ataxmbis0}
\end{align} Las últimas dos ecuaciones reciben nombres particulares
\marginnote{%\
La expresión $A^T A x = A^T b$ suele ser referida como las
“ecuaciones normales”. La derivación de la solución de cuadrados mínimos
presentada en este trabajo no es la única. De hecho, es común encarar el
problema mediante cálculo diferencial. En ese contexto, las ecuaciones normales
surgen naturalmente a la hora de optimizar $E(x)$ igualando su gradiente a $0$.
Para una demostración mediante cálculo referirse a
\textcite{nocedalNumericalOptimization2006}, cap. 10.
} \marginnote{%\
La expresión $A^T (A\hat{x} - b) = 0$ es interesante,
ya que muestra que las columnas de $A$ son ortogonales al residual óptimo
$A\hat{x} - b$ al ser su producto interno $0$. Esto suele llamarse el “principio
de ortogonalidad”.
}. Veamos ahora para un \(x \in \R^n\) cualquiera. \begin{align}
  & \| Ax - b \|^2 = \| (Ax - A\hat{x}) + (A\hat{x} - b) \|^2 \\
  & \text{Por \cref{rmk:sqnormofsum}:}\notag\\
  &= \| A(x - \hat{x}) \|^2 +  \| A\hat{x} - b \|^2 + 2 (A(x - \hat{x}))^T
  (A\hat{x} - b) \\
  &= \| A(x - \hat{x}) \|^2 +  \| A\hat{x} - b \|^2 + 2 (x - \hat{x})^T A^T
  (A\hat{x} - b) \\
  & \text{Por \cref{eq:ataxmbis0}:}\notag\\
  &= \| A(x - \hat{x}) \|^2 +  \| A\hat{x} - b \|^2 \\
  &\therefore \quad \|Ax-b\|^2 \geq \|A\hat{x}-b\|^2
\end{align} Más aún, esta solución es única, ya que la igualdad en la
última ecuación solo se da si \(\|A(x-\hat{x})\|^2 = 0\), y como \(A\)
es no-singular, esto solo pasa cuando \(x=\hat{x}\).\end{proof}

\hypertarget{cuadrados-muxednimos-no-lineales}{%
\subsection{Cuadrados mínimos no
lineales}\label{cuadrados-muxednimos-no-lineales}}

Generalizaremos el problema anterior para que también considere
funciones no lineales. Reutilizaremos la notación introducida en la
sección anterior.

En este caso permitimos ahora que las \(f_i\) sean no lineales, aunque
requerimos que continúen siendo diferenciables. Además, ignoraremos el
vector \(b\) haciendo que \(f(x) = r(x)\). Es decir querremos que
\(f(x) = 0\) en lugar de \(f(x) = b\). En el caso de necesitar la
segunda restricción, podemos construir nuevas funciones
\(\tilde{f}_i(x) = f_i(x) - b_i\) y utilizar estas en su lugar. Nuestra
función de error es entonces: \begin{align}
E(x) = \sum_{i=1}^m{f_i(x)^2} = \| f(x) \| ^ 2
\end{align}

Encontrar \(x\) solución que minimice tal error es el \emph{problema de
los cuadrados mínimos no lineales}. En este caso, al no ser necesario
que los residuales sean \emph{funciones afines} (lineales más un
escalar), no podemos dar una matriz \(A\) y utilizar ideas como las de
la pseudo inversa \(A^{\dagger}\). Más aún, en el caso lineal, si bien
no lo mostramos, se cumple que la solución no solo es única, sino que es
el único punto con gradiente \(\nabla f(x) = 0\), es decir, el único
punto optimizador, es un mínimo global. Para una demostración de esto,
se utilizan argumentos de \emph{convexidad} sobre \(E(x)\); referirse a
\textcite{nocedalNumericalOptimization2006}, cap. 10. En el caso no
lineal sin embargo, no tenemos ninguna de estas garantías, pueden
existir infinitos máximos y mínimos globales, locales y cualquier
combinación de estos.

Lo que se hace entonces es utilizar algoritmos heurísticos; en nuestro
caso el algoritmo de \emph{Gauss-Newton} que busca un mínimo de forma
iterativa comenzando desde un punto que, si se encuentra lo
suficientemente cerca a la solución, convergerá a ella. Como veremos más
adelante, este requerimiento de proveer un punto inicial adecuado es muy
razonable en sistemas de SLAM/VIO, ya que usualmente contaremos con
dicha información.

Gauss-Newton genera la secuencia de puntos \(x^{(1)},\ x^{(2)}, ...\) al
ser un algoritmo iterativo. Es posible juzgar cada iteración \(k\)
evaluando \(E(x^{(k)}) = \| f(x^{(k)})\|^2\). Como es usual existen
varios criterios de terminación apropiados: se alcanza un cierto número
fijo de iteraciones \(k^{max}\), el error del último iterando
\(E(x^{(k)})\) es suficientemente cercano a \(0\), o las últimas dos
iteraciones tienen un error muy similar a los fines prácticos
\(E(x^{(k)}) \sim E(x^{(k + 1)})\).

En cada iteración \(k\) el algoritmo cuenta con dos etapas. La etapa de
\textbf{linealización} aproxima a \(f(x)\) linealmente con su expansión
de Taylor de primer orden centrada en el iterando actual \(x^{(k)}\);
llamaremos a esta aproximación \(f^{(k)}(x)\). En la etapa de
\textbf{actualización} utilizamos el hecho de que \(f^{(k)}(x)\) es una
función afín para encontrar una solución al problema de cuadrados
mínimos \emph{lineales} sobre ella. Esta solución será el valor de
nuestro próximo iterando \(x^{(k+1)}\) y, por ende, la estimación que
damos como solución del problema no-lineal. Es decir, como sabemos que
\(x^{(k+1)}\) minimiza a \(f^{(k)}\), esperaríamos que sea una buena
aproximación al mínimo de \(f\) sabiendo que comenzamos el algoritmo
cerca de su mínimo.

\textbf{Linealización}. Definimos \(f^{(k)}(x)\) como la expansión de
Taylor de primer orden de \(f(x)\) centrada en \(x^{(k)}\). Es decir:
\begin{align}
f^{(k)}(x) = f(x^{(k)}) + A(x^{(k)}) (x - x^{(k)})
\end{align}

con \(A: \R^n \rightarrow \R^{m \times n}\) la matriz jacobiana de
\(f\). \begin{align}
A = \begin{bmatrix}
  \frac{\partial f_1}{\partial x_1} & \dots & \frac{\partial f_1}{\partial x_n} \\
  \vdots & \ddots & \vdots \\
  \frac{\partial f_m}{\partial x_1} & \dots & \frac{\partial f_m}{\partial x_n}
\end{bmatrix}
\end{align}

Notar que \(f^{(k)}\) es afín; una transformación lineal \(A(x^{(k)})\)
aplicada sobre \(x\), más un vector. Reacomodemos los términos para
dejarlo explícito e introduzcamos la notación \(A_k = A(x^{(k)})\).
\begin{align}
f^{(k)}(x) = f(x^{(k)}) + A_k (x - x^{(k)}) \\
f^{(k)}(x) = A_k x - (A_k x^{(k)} - f(x^{(k)}))
\end{align}

\textbf{Actualización}. Teniendo en mente que queremos encontrar el
mínimo de \(\| f(x) \| ^ 2\), lo aproximamos con el mínimo (y próximo
iterando) \(x^{(k+1 )}\) de \(\| f^{(k)}(x) \| ^ 2\) pero como
\(f^{(k)}\) es afín, esto se reduce a buscar la solución del problema de
cuadrados mínimos \emph{lineales}. Sabemos por
\Cref{thm:linleastsquaresol} la solución exacta para este caso con base
en la matriz pseudo-inversa \(A_k^{\dagger} = (A_k^T A_k)^{-1} A_k^T\).
\begin{align}
\| f^{(k)}(x) \| ^ 2 = \| A_k x - (A_k x^{(k)} - f(x^{(k)})) \|^2
\end{align}

Se minimiza por \Cref{thm:linleastsquaresol} cuando \begin{align}
x &= A_k^{\dagger} (A_k x^{(k)} - f(x^{(k)})) \\
&=A_k^{\dagger} A_k x^{(k)} - A_k^{\dagger} f(x^{(k)}) \\
&= x^{(k)} - A_k^{\dagger} f(x^{(k)})
\end{align}

Entonces el algoritmo de Gauss-Newton queda definido de la siguiente
manera:

\begin{algorithm}[H]
\caption{Gauss-Newton para cuadrados mínimos no lineales}

Dada $f : \R^n \rightarrow \R^m$ diferenciable y un punto inicial $x^{(1)}$.
Con $k = 1, 2, ..., k^{max}$. \newline

1. Linealizar $f$ alrededor de $x^{(k)}$ computando el jacobiano $A_k$:
\begin{align}
f^{(k)}(x) &= f(x^{(k)}) + A_k (x - x^{(k)})
\end{align}

1. Actualizar el iterador a $x^{(k+1)}$ con el minimizador de $\| f^{(k)}(x) \|^2$
descripto en el \cref{thm:linleastsquaresol}. Se necesitará computar $A_k^{\dagger}$
con la inversa de la matriz de Gram de $A_k$:
\begin{align}
x^{(k+1)} = x^{(k)} - A_k^{\dagger} f(x^{(k)}) =
x^{(k)} - (A_k^T A_k)^{-1} A_k^T f(x^{(k)})
\end{align}

\end{algorithm}

Gauss-Newton y minimización lineal con la pseudo-inversa serán de gran
utilidad para expresar numerosos tipos de problemas de optimización. Se
utilizará para problemas como ajustar la posición de un punto de interés
tridimensional que es observado por múltiples cámaras; desproyectar
puntos 2D de cámaras hacia puntos 3D en la escena cuando los modelos de
proyección no tienen expresiones cerradas para su inversa; incluso
veremos que la optimización central de los sistemas de SLAM/VIO que
combina todas las mediciones necesarias para generar las estimaciones de
poses, usualmente llamada \emph{bundle
adjustment}\label{def:bundle-adjustment}, suele expresarse y resolverse
como una minimización de cuadrados no lineales. El desarrollo de esta
sección está basada en \textcite{nocedalNumericalOptimization2006}, cap.
1, 2 y 10 y \textcite{boydIntroductionAppliedLinear2018}, cap. 11, 12,
15 y 18; en esos trabajos se puede encontrar derivaciones alternativas e
información de otras técnicas más sofisticadas que Gauss-Newton como
Levenberg-Marquardt o el método dogleg.

\cleardoublepage
\ctparttext{En esta parte contextualizaremos los distintos sistemas estudiados y
profundizaremos en uno de ellos: \italic{Basalt}. El estudio detallado de una
implementación será particularmente esclarecedor al hilar sin generalizaciones
multitud de métodos y algoritmos utilizados en contextos concretos y con
objetivos bien definidos.}

\hypertarget{implementaciuxf3n-de-un-sistema}{%
\part{Implementación de un
sistema}\label{implementaciuxf3n-de-un-sistema}}

\hypertarget{basalt}{%
\chapter{Basalt}\label{basalt}}

\hypertarget{preliminares-1}{%
\section{Preliminares}\label{preliminares-1}}

La decisión sobre cuál sistema de localización visual-inercial integrar
con Monado no fue sencilla. Fue necesario descartar docenas de
implementaciones e incluso así, no fue trivial entender si los sistemas
elegidos finalmente resultaron los más adecuados para el contexto de XR.
Cada implementación presentaba ventajas y desventajas, aunque es usual
que las respectivas publicaciones se concentren en destacar solo las
métricas favorables. Más aún, es necesario conocimiento experto para
comprender cómo la elección de ciertos fundamentos teóricos, técnicas
algorítmicas, decisiones arquitecturales o de tecnología afectan a la
calidad del tracking dedicado a XR. Gran parte de este trabajo se basó
en el estudio de los conceptos necesarios para poder tomar este tipo de
decisiones.

A grandes rasgos y sin un orden en particular, las propiedades deseables
que se consideraron a la hora de elegir sistemas fueron:

\begin{enumerate}
\def\labelenumi{\arabic{enumi}.}
\item
  Versatilidad en la configuración sensores: Monado necesita soportar
  una gran variedad de dispositivos, es por esto que se prefirieron
  sistemas que soporten la mayor cantidad de combinaciones y tipos de
  sensores. En el mejor de los casos, Monado debería ser capaz de
  localizar desde cascos con cámaras estéreo y una IMU, hasta celulares
  con una única cámara y sin giroscopio.
\item
  Licencia permisiva: La licencia y filosofía de Monado da gran libertad
  al programador que lo utilice de hacer lo que desee con su código
  fuente. Esto incluye poder utilizarlo en proyectos en dónde se
  prefiere no distribuir dicho código. Licencias de código libre
  ``virales'' como la GPL \autocite{GNUGeneralPublic} no permiten esto y
  enlazar Monado a sistemas con este tipo de licencias contagiaría a los
  proyectos que dependen de Monado quitándoles la libertad de decidir no
  publicar su código.
\item
  Desarrolladores activos: Estos sistemas suelen surgir de grupos de
  investigación y es muy común que se abandonen luego de que el trabajo
  sea publicado. Será necesario encontrar, dentro de lo posible,
  sistemas con un desarrollo activo, con mantenedores accesibles y que,
  en el mejor de los casos, acepten contribuciones a su proyecto.
\item
  Estabilidad, facilidad de instalación y buenas prácticas de
  desarrollo: Los sistemas de SLAM son muy complejos y es usual que se
  optimicen para funcionar únicamente en los conjuntos de datos de
  prueba dejando de lado estas características que son fundamentales a
  la hora de querer brindar un sistema para ser utilizado por usuarios
  finales.
\item
  Rendimiento: El área de XR tiende a buscar la reducción del tamaño de
  los dispositivos para mejorar su ergonomía y practicidad. Un sistema
  de tracking debe ser capaz de operar en contextos con recursos
  acotados, debería utilizar poca memoria, poca energía y ser capaz de
  estimar poses a altas frecuencias y utilizando poca capacidad de
  cómputo.
\item
  Precisión: Por último y no menos importante, queremos que la precisión
  de la localización sea adecuada. El nivel de precisión requerido
  dependerá del tipo de aplicación. Es común que se necesite precisión
  submilimétrica en contextos de VR donde la simulación cubre
  completamente el campo de visión del usuario, y fallos en el tracking
  puedan inducir mareos. Mientras que para otros contextos como AR
  realizado por un celular, no es tan vital contar con ese nivel de
  exactitud.
\end{enumerate}

\hypertarget{sistemas-integrados}{%
\subsection{Sistemas integrados}\label{sistemas-integrados}}

En este trabajo se integraron con Monado tres sistemas de código libre
distintos, primero \emph{Kimera-VIO}
\autocite{rosinolKimeraOpenSourceLibrary2020}, luego \emph{ORB-SLAM3}
\autocite{camposORBSLAM3AccurateOpenSource2021} y finalmente
\emph{Basalt} \autocite{usenkoBasaltVisualInertialMapping2020}.

Kimera-VIO\footnote{\url{https://github.com/MIT-SPARK/Kimera-VIO}} es
una implementación desarrollada en el Instituto de Tecnología de
Massachusetts (MIT) con una licencia permisiva BSD-2
\autocite{2ClauseBSDLicense}. Fue publicada originalmente en octubre de
2019 y su última actualización significativa\footnote{\url{https://github.com/MIT-SPARK/Kimera-VIO/pull/152}},
en dónde se introduce la posibilidad utilizar una única cámara, ocurrió
en abril de 2021. Kimera resultó inicialmente atractivo por su licencia
y su promesa de correr en tiempo real. Además de esto posee
funcionalidades muy interesantes para XR como la reconstrucción y
análisis semántico del entorno en el que el dispositivo se encuentra.
Desgraciadamente, no logró ser adecuado mostrando grandes dificultades a
la hora de correrlo en sensores distintos a los presentados en su
publicación. Más aún el término ``tiempo real'' en el contexto de la
publicación es ambiguo, ya que está lejos de ser capaz de ejecutarse a
frecuencias adecuadas para XR y es más adecuado para uso en robots con
frecuencias de menos de 10 Hz \autocite[Fig.
5]{rosinolKimeraOpenSourceLibrary2020}.

ORB-SLAM3\footnote{\url{https://github.com/UZ-SLAMLab/ORB_SLAM3}} es
desarrollado por la Universidad de Zaragoza con una licencia viral
GPL-3.0 \autocite{GNUGeneralPublic}. Fue publicado inicialmente en julio
de 2020 y su última actualización significativa\footnote{\url{https://github.com/UZ-SLAMLab/ORB_SLAM3/releases/tag/v1.0-release}}
ocurrió en diciembre de 2021. ORB-SLAM3 es la tercera iteración de una
línea de sistemas que dominan las tablas comparativas de SLAM desde hace
unos cuantos años y es por esto que se implementó incluso aunque no
posea una licencia permisiva. El sistema presenta varios métodos
novedosos que mejoran la precisión del tracking mientras que muestra ser
capaz de estimar poses a unos 20 Hz o 30 Hz \autocite[Tabla
6]{camposORBSLAM3AccurateOpenSource2021}. El sistema es el más versátil
del campo permitiendo ser ejecutado en configuraciones con una cámara
(monocular) o dos (estéreo), con o sin uso de la IMU e incluso con
cámaras de profundidad. Además, soporta la reconstrucción de múltiples
mapas y la capacidad de interconectarlos o incluso guardarlos en
almacenamiento persistente para ser reutilizados en sesiones de tracking
posteriores. Una desventaja es que la trayectoria que se da al construir
el mapa es particularmente ruidosa y difícil de utilizar en XR. Se
plantea como trabajo a futuro investigar más acerca de la funcionalidad
de reutilización del mapa y como funciona el tracking con mapas ya
construidos.

Finalmente, Basalt\footnote{\url{https://gitlab.com/VladyslavUsenko/basalt}}
es desarrollado por el Instituto Técnico de Múnich con una licencia
permisiva BSD-3 \autocite{3ClauseBSDLicense}. Fue publicado
originalmente en abril de 2019 y su última actualización
significativa\footnote{\url{https://gitlab.com/VladyslavUsenko/basalt/-/commit/24325f2a}}
\autocite{demmelBasaltSquareRoot2021} ocurrió en octubre de 2021. Basalt
solo es capaz de correr en tiempo real su sistema de VIO necesitando de
una pasada offline de su \emph{``mapper''} para lograr un mapa
consistente. A pesar de esto, mostró ser sorprendentemente preciso y
tener un gran desempeño. En particular, es un sistema que puede correr
fácilmente a 60 Hz consumiendo cuadros a mayores resoluciones que las
probadas en ORB-SLAM3 y Kimera y duplicando las frecuencias de muestreo
a 60 fps. Esta capacidad de soportar mayor cantidad de muestras aumenta
significativamente la precisión de la trayectoria a pesar de que no sea
un sistema de SLAM completo. Además de esto, el sistema es notablemente
más sencillo de compilar al tener un buen manejo de sus dependencias,
posee mejores prácticas de ingeniería de software y es en general mucho
más estable logrando fácilmente sesiones de tracking ininterrumpidas a
diferencia de los sistemas anteriores.

Los tres sistemas fueron integrados en Monado con distintos niveles de
éxito, pueden verse demostraciones de cómo funcionan en los videos
referenciados al pie de página \footnote{Kimera-VIO con Monado:
  \url{https://youtu.be/gxu3Ve8VCnI}} \footnote{ORB-SLAM3 con Monado:
  \url{https://youtu.be/kJwWY973b10}} \footnote{Basalt con Monado:
  \url{https://youtu.be/ajuqQ7E1MFw}}. Más adelante, en el
\Cref{evaluation}, se verán resultados y distintas métricas comparativas
entre los sistemas.

Las terminologías y jerga del área de SLAM/VIO pueden resultar
abrumadoras, pero se espera que introduciéndolas en el contexto de una
implementación concreta se puedan entender mejor los problemas que los
sistemas, en general, tienen que resolver. Por todas las razones
mencionadas anteriormente, Basalt es actualmente el sistema de
preferencia para ser utilizado con Monado y, si bien se estudió el
código fuentes de los tres sistemas, en esta parte del trabajo que sigue
a continuación \emph{nos vamos a enfocar en profundizar en la
implementación de Basalt}.

\hypertarget{problemuxe1ticas-de-un-sistema}{%
\subsection{Problemáticas de un
sistema}\label{problemuxe1ticas-de-un-sistema}}

Un problema central en este tipo de sistemas es el de poder generar un
mapa y una trayectoria que sean \emph{globalmente consistentes}. Con
esto nos referimos a que nuevas mediciones tengan en cuenta todas las
mediciones anteriores en el sistema. Una forma ingenua de encarar esto,
sería realizando \emph{bundle adjustment}\marginnote{%\
Introdujimos el término bundle adjustment (BA) en la \Cref{def:bundle-adjustment}
en el contexto de cuadrados mínimos. Este se refiere al refinamiento simultáneo de un conjunto de
poses de cámara, o vistas, y puntos 3D en el mapa que han sido observados por
estas vistas. Así, se intenta reducir el llamado “error de reproyección” del
conjunto actualizando tanto las poses de las cámaras como la posición de los
puntos observados. Este error hace referencia a la distancia entre las
posiciones de los puntos y en dónde las vistas esperan que estos
se encuentren \autocite{hartleyMultipleViewGeometry2004}.
} sobre todas las imágenes capturadas a lo largo de una corrida,
integrando además las mediciones provenientes de la IMU.
Desafortunadamente, este método excede rápidamente cualquier capacidad
de cómputo de la que dispongamos, y aún más teniendo en cuenta que
nuestro objetivo es localizar en tiempo real al dispositivo de XR.

Por esta razón, es usual recurrir a distintas formas de reducir la
complejidad del problema. Para realizar \emph{odometría visual-inercial
(VIO)}, es común que se ejecute la función de optimización sobre una
\emph{ventana local} de cuadros y muestras recientemente capturadas,
ignorando muestras históricas y acumulando error en las estimaciones a
lo largo del tiempo. Además, esta mirada tiene la desventaja adicional
de que una porción significativa de los fotogramas capturados podrían
tener posiciones similares que no añadirían demasiada información al
estimador, o incluso que algunos fotogramas puedan ser de baja calidad
por contener \emph{motion blur} u otro tipo de anomalías. Por otro lado,
soluciones que intentan hacer \emph{mapeo visual-inercial} realizan el
bundle adjustment sin utilizar todas las imágenes capturadas, sino que
se limitan a la utilización de algunos fotogramas clave, o
\emph{keyframes} elegidos mediante criterios que priorizan cuadros
nítidos y con distancias (\emph{baselines}) prudenciales entre ellos.

Como las muestras de IMU vienen a mayor frecuencia que las de la cámara,
es común que estas se \emph{preintegren} de forma tal de combinar
muestras simultáneas entre dos keyframes en una única entrada del
optimizador. Sin embargo, un problema en el que esta integración
incurre, es que las mediciones de las IMU son altamente ruidosas, y
acumularlas durante tiempos prolongados acumula también cantidades
significativas de error. Este factor nos limita el tiempo que puede
transcurrir entre dos keyframes; como ejemplo en
\textcite{mur-artalVisualInertialMonocularSLAM2017} se habla de
keyframes que no pueden tener más de medio segundo entre sí. Además,
tener keyframes a muy bajas frecuencias afecta la calidad de las
estimaciones de velocidad y \emph{biases}; estos últimos son offsets de
medición inherentemente variables de los acelerómetros y giroscopios a
los que es necesario reestimar de forma constante para compensar por
ellos en la medición final.

\hypertarget{propuesta-de-basalt}{%
\subsection{Propuesta de Basalt}\label{propuesta-de-basalt}}

La novedad de Basalt \autocite{usenkoBasaltVisualInertialMapping2020} es
que formula el mapeo visual-inercial como un problema de bundle
adjustment y utiliza, de una forma específica, todas mediciones visuales
e inerciales a altas frecuencias. Usa un \emph{grafo de
factores}\marginnote{%\
Los grafos de factores son una muy buena forma de representar problemas con
muchas variables aleatorias interdependientes y muestras que las relacionan (factores).
En general, el uso de estos grafos trae beneficios
computacionales interesantes y son de gran importancia para el área de SLAM.
Sin embargo, no nos adentraremos demasiado en el tema en este trabajo y dirigimos
al lector interesado a \textcite{dellaertFactorGraphsRobot2017}.
} \footnote{Artículo introductorio a los grafos de factores:
  \url{https://gtsam.org/2020/06/01/factor-graphs.html}} de forma
similar a otros sistemas, también llamado \emph{grafo de poses} en este
contexto por contener poses a estimar como nodos. En lugar de utilizar
todos los fotogramas se propone realizar la optimización en dos capas.
La capa de VIO, emplea un sistema de odometría visual-inercial, que ya
de por sí supera a otros sistemas del mismo tipo, proveyendo
estimaciones de movimiento a la misma frecuencia que el sensor de la
cámara provee imágenes. Luego, se seleccionan keyframes y se agregan
\emph{factores no-lineales} entre estos que estiman la diferencia de
posición relativa. Estos dos factores, keyframes y poses relativas, se
utilizan en la capa de bundle-adjustment global.

La capa de VIO, detecta \emph{features}\marginnote{%\
Las features son puntos de interés relevados en una imagen.
Son estos los puntos que triangularemos en el bundle adjustment.
} que son rápidas y buenas para seguir durante varios cuadros (esto es
el \emph{optical flow} que veremos en la sección \Cref{optical-flow}),
mientras que en la capa de mapeo se usan features adecuadas que son
indiferentes a las condiciones de luz o al punto de vista de la
cámara\marginnote{%\
Esto es fundamental para entender cuándo el dispositivo está visitando un lugar
por el que ya pasó. Esto se denomina “loop closing”.
}. De esta forma tenemos un sistema que es capaz de utilizar las
mediciones a alta frecuencias de los sensores y al mismo tiempo tiene la
capacidad de detectar cuando se está en ubicaciones ya visitadas,
obteniendo así un mapa que es globalmente consistente. Además, el
problema de optimización se reduce, ya que a diferencia de otros
sistemas, no es necesario estimar velocidades ni biases (de la IMU).

\hypertarget{implementaciuxf3n}{%
\section{Implementación}\label{implementaciuxf3n}}

A continuación describiremos la arquitectura e implementación de Basalt
de una manera más detallada. Estas secciones surgen directamente de la
lectura del código fuente del sistema e intentan proveer detalles más
bien pragmáticos que se encuentran en el mismo, pero que pueden quedar
escondidos en las publicaciones de más alto nivel que presentan estos
sistemas. A su vez, se toman ciertas licencias literarias que deberían
ayudar al entendimiento y que no son posibles a la hora de escribir
código.

Cómo vimos anteriormente, el funcionamiento de Basalt se divide en dos
etapas. La primera etapa de odometría visual-inercial (VIO), en el cual
se emplea un sistema de VIO que supera a sistemas equivalentes de
vanguardia mientras que la segunda etapa de mapeo visual-inercial (VIM),
toma keyframes producidos por la capa de VIO y ejecuta un algoritmo de
bundle adjustment para obtener un mapa global consistente. Estas dos
capas son completamente independientes. En una corrida usual, se ejecuta
inicialmente el sistema de VIO y es este el que decide y almacena
persistentemente qué cuadros y con qué información el sistema de VIM, de
ejecutarse, debería utilizar al realizar el proceso de bundle
adjustment.

Esto significa que, por defecto, no contamos con la capacidad de
utilizar el VIM en tiempo real para XR, solo el VIO. Por ende solo este
fue integrado con Monado. Se plantea como trabajo a futuro la
paralelización del VIM en un hilo separado para poder correrlo en tiempo
real\footnote{Discusión sobre como adaptar Basalt para poder correr el
  VIM en tiempo real:
  \url{https://gitlab.com/VladyslavUsenko/basalt/-/issues/69}}.
Exploraremos entonces, en esta parte del trabajo, los componentes
fundamentales de la capa de VIO: \emph{optical flow}, \emph{bundle
adjustment visual-inercial} y finalmente el proceso de
\emph{optimización y de marginalización parcial}.

\hypertarget{optical-flow}{%
\subsection{Optical flow}\label{optical-flow}}

El módulo de VIO toma dos tipos de entrada, una de ellas son las
muestras raw de la IMU; y la otra, contra intuitivamente, no son las
imágenes raw provenientes de las cámaras, sino que son los
\emph{keypoints} resultantes de ellas. Estos son la posición en dos
dimensiones sobre el plano de la imagen de las \emph{features}
detectadas. Las features a su vez son la representación de los puntos de
interés o \emph{landmarks} de la escena tridimensional proyectados sobre
las imágenes. El proceso de detectar features, computar su
transformación entre distintos cuadros, y producir los keypoints de
entrada para el módulo de VIO, está a cargo del módulo de \emph{optical
flow} (o \emph{flujo óptico}). Cabe aclarar que optical flow es el
nombre que recibe tanto el campo vectorial que representa el movimiento
aparente de puntos entre dos imágenes, como el proceso de estimarlo.
Este puede ser denso, si se considera el flujo de todos los píxeles, o
no (\emph{sparse}) si solo se computa el flujo de algunos keypoints como
es el caso que veremos.

El módulo de optical flow corre en un hilo separado y es por donde las
muestras del par de cámaras estéreo ingresan al pipeline de Basalt.
Inicialmente se genera una representación piramidal de las imágenes, o
también llamada de \emph{mipmaps}, esta es una forma tradicional
\autocite{williamsPyramidalParametrics1983} de almacenar una imagen en
memoria junto a versiones reescaladas de la misma como se ve en la
\figref{fig:mipmap}. Los mipmaps tienen múltiples utilidades en
computación gráfica (p.~ej. \emph{filtrado trilineal}, \emph{LODs},
reducción de \emph{patrones moiré}) pero en el caso de Basalt serán
utilizados para darle robustez al algoritmo de seguimiento de features o
\emph{feature tracking}.

\fig{fig:mipmap}{source/figures/mipmap.png}{Mipmaps}{%
Representación piramidal (mipmaps) de un cuadro del conjunto de datos estándar
EuRoC \autocite{burriEuRoCMicroAerial2016}.
}

Posteriormente se realiza la detección de features nuevas sobre las
imágenes utilizando el algoritmo \emph{FAST}
\autocite{rostenFasterBetterMachine2010a} para detección de esquinas, o
puntos sobresalientes de la imagen en general (keypoints), implementado
sobre \emph{OpenCV}\marginnote{%\
OpenCV es una de las bibliotecas de visión por computadora más populares y
utilizadas en este tipo de sistemas. Combina múltiples algoritmos y presenta
una licencia permisiva Apache 2 \autocite{ApacheLicenseVersion} (prev. BSD-3 \autocite{3ClauseBSDLicense}).
}. Resulta importante aclarar que Basalt es uno de los sistemas que
menos depende de OpenCV, siendo \passthrough{\lstinline!cv::FAST!} el
único algoritmo de la biblioteca en uso durante una ejecución usual. El
proyecto tiende a reimplementar muchas de las técnicas y algoritmia de
forma especializada y, como veremos en otros módulos, otras tareas
razonablemente complejas como la optimización de grafos de poses se
implementan también dentro del proyecto y sin recurrir a bibliotecas
externas. Esta es una de las varias razones por las que el sistema logra
tan buen rendimiento, ya que estas bibliotecas suelen necesitar de
campos y comprobaciones generales del problema que intenta solucionar,
mientras que Basalt puede prescindir de todas las que no apliquen a VIO.
Aunque, también es cierto que estas reimplementaciones añaden
complejidad al sistema.

Siguiendo con la detección de features, una heurística particular de
Basalt es la división del cuadro completo en celdas de, por defecto, 50
por 50 píxeles en donde se detectan los nuevos puntos de interés. Por
celda solo se conserva la feature de mejor calidad o con mejor
\emph{respuesta (response)}, aunque se puede configurar para detectar
más de una. Siempre que la celda tenga alguna localizada de cuadros
anteriores, no se intenta detectar nuevas. Esto contrasta con sistemas
como Kimera-VIO que corren la detección FAST sobre el cuadro entero y
evitan la redetección mediante el uso de \emph{máscaras} que le
instruyen al algoritmo a obviar esas secciones. Desafortunadamente la
construcción de tales máscaras suele ser costosa y la heurística de
Basalt, a pesar de desperdiciar espacio por no permitir la detección de
nuevas características entre celdas, es más eficiente, ya que en
situaciones comunes se logran detectar una cantidad razonable de
features sin problemas. Esta detección se realiza únicamente sobre la
primera cámara, usualmente la izquierda, mientras que en la otra cámara
se reutiliza el método de seguimiento de keypoints que se describe a
continuación.

En cada instante de tiempo que entran un nuevo par de imágenes se tiene
acceso a toda la información recolectada del instante anterior, en
particular a sus keypoints. Una suposición razonable es que las imágenes
correspondientes a este nuevo instante van a compartir muchos de los
keypoints con las imágenes anteriores y en posiciones similares. Con esa
suposición Basalt logra ahorrarse tener que volver a detectar features
de la imagen con FAST y en cambio el problema se transforma en, dado una
imagen anterior (inicial), sus keypoints y una imagen nueva (objetivo),
estimar donde ocurren esos mismos keypoints en la imagen nueva. Para
esto, por cada keypoint anterior, se genera un parche \(\Omega\)
alrededor de su ubicación de, por defecto, 52 puntos como se ve en el
ejemplo de la \figref{fig:patches}.

\fig{fig:patches}{source/figures/patches.png}{Parches}{%
Ejemplos de los parches de 52 puntos considerados para computar el optical flow
de distintos keypoints de una imagen.
}

Considerando entonces que este parche debería estar en la imagen nueva
en coordenadas cercanas a las del keypoint anterior, queremos encontrar
la transformación \(\mathbf{T} \in SE(2)\) que le ocurrió al parche, y
por ende al nuevo keypoint que se encontraría en el centro de este nuevo
parche. Basalt emplea entonces optimización por cuadrados mínimos
mediante el algoritmo iterativo de Gauss-Newton para encontrar
\(\mathbf{T}\) utilizando un residual \(r\) con:

\[
r_i =
  \frac{I_{t + 1}(\mathbf{T} \mathbf{x}_i)}{\overline{I_{t + 1}}} -
  \frac{I_{t}(\mathbf{x}_i)}{\overline{I_{t}}}
  \ \ \ \ \forall \mathbf{x}_i \in \Omega
\]

Y aquí siendo \(I_t(\mathbf{x})\) la intensidad de la imagen anterior en
el pixel ubicado en las coordenadas \(\mathbf{x}\), análogamente
\(I_{t + 1}(\mathbf{x})\) para la imagen objetivo; y siendo
\(\overline{I_{t}}\) la intensidad media del parche \(\Omega\) en la
imagen inicial, análogamente \(\overline{I_{t + 1}}\) para la imagen
objetivo y el parche transformado \(\mathbf{T}\Omega\). Notar que al
normalizar las intensidades obtenemos un valor que es invariante incluso
ante cambios de iluminación.

Los detalles del cálculo de gradientes y jacobianos están basados en el
método de \textcite{lucasIterativeImageRegistration1981} para tracking
de features (\emph{KLT}). El uso adicional de mipmaps sobre KLT fue
originalmente expuesto en
\textcite{bouguetPyramidalImplementationLucas1999}.

Para asegurar que la estimación fue exitosa se invierte el problema y se
intenta trackear desde la imagen nueva hacia la inicial y, si el
resultado está muy alejado de la posición inicial, el nuevo keypoint se
considera inválido (un \emph{outlier}) y se lo descarta. Otro detalle a
aclarar es que, recordando que la detección costosa de features con FAST
solo ocurría en las imágenes de una de las cámaras, es posible ahora
entender que las features en la segunda cámara pueden ser ``detectadas''
con este método menos costoso. Es decir, simplemente se computa el
optical flow desde la imagen de la cámara izquierda a la de la cámara
derecha en el mismo instante de tiempo.

Finalmente, el último de los pasos que ocurre cuando el módulo de
optical flow procesa un cuadro, es el de filtrado de keypoints, en el
cual se desproyectan los keypoints a posiciones en la escena
tridimensional y en caso de que el error de reproyección supere cierto
umbral, estos keypoints serán descartados por considerarse outliers.

\hypertarget{bundle-adjustment-visual-inercial}{%
\subsection{Bundle adjustment
visual-inercial}\label{bundle-adjustment-visual-inercial}}

En un hilo separado al módulo de optical flow, corre el estimador de VIO
encargado de realizar el bundle adjustment sobre los cuadros y muestras
de la IMU recientes para estimar la pose. Este toma como entrada las
muestras de la IMU junto a los keypoints 2D detectados para cada imagen,
o sea la salida del módulo de optical flow. Este módulo es el que
efectivamente realizará la integración y optimización con toda la
información recibida y producirá como salida en una cola, la estimación
de los estados del agente a localizar.

\hypertarget{basalt-preintegration}{%
\subsubsection{Inicialización y
preintegración}\label{basalt-preintegration}}

\begin{mdframed}[backgroundcolor=shadecolor]
TODO: En el párrafo que sigue hablo de una sección “Calibración de IMU” que todavía no
existe, tengo la mitad de esa sección escrita pero todavía estoy considerando si
agregarla al trabajo o simplemente referenciarla de un libro.
Cuando uno trabaja con sensores se tienen modelos matemáticos que describen las
distorsiones que se suelen dar en el sensor. Calibrar un sensor significa
encontrar los valores de los parámetros de este modelo para tu sensor particular.

Se usa también cuadrados mínimos no lineales para calibrar un sensor. Uno toma varias
mediciones y despues optimiza para encontrar los parámetros más adecuados.

Esto aplica exactamente igual para las cámaras (aunque ahí hay muchos más
modelos en uso). Y hay unos detalles muy copados de los que hablar por que
el equipo de Basalt desarrolló sus propios modelos de cámaras que son muy
rápidos por que están bien pensados para SLAM (la inversa de los modelos tienen expresión
cerrada que no es usual en los modelos de cámara estándar). Además tengo un merge request abierto
en donde contribuyo un nuevo modelo de cámara que se usa en unos cascos que les
hicimos ingeniería inversa (los de Windows Mixed Reality) y que Basalt no
soportaba anteriormente:
\url{https://gitlab.com/VladyslavUsenko/basalt-headers/-/merge_requests/21}.

Así que sí, me gustaría detallar más todo esto en el trabajo, pero el tiempo se
está acabando.
\end{mdframed}

Para comenzar, el hilo de procesamiento de este módulo espera a que la
primera muestra de la IMU arribe. Estas muestras son recibidas de forma
raw y antes de tratarlas, Basalt utiliza los parámetros de calibración
estáticos provistos por archivos de configuración para corregirlas. La
corrección, o calibración, ocurre como se explica en la sección de
``Calibración de IMU'' sin considerar todavía los parámetros de los
procesos aleatorios detallados en la misma. Una particularidad de
Basalt, es que el acelerómetro es utilizado como origen del agente
localizado y, fundándose en esto, se fija su orientación. Es decir, no
se aplica ningún tipo de corrección de orientación al calibrar las
muestras del acelerómetro. Esto hace que la matriz de alineamiento para
el acelerómetro tenga ceros en su triángulo superior (ver
\textcite{schubertBasaltTUMVI2018} secc. IV.B y \emph{discusión
relacionada \footnote{\url{https://gitlab.com/VladyslavUsenko/basalt-headers/-/issues/8}}}).

Luego de recibir esta primer muestra de la IMU se comienza la ejecución
del bucle principal, el cual espera indefinidamente por resultados
encolados por el módulo de optical flow para realizar una iteración. El
primer par de muestras de cámaras junto a la primera muestra de la IMU
posterior al par estéreo son utilizados para inicializar el estado del
agente en el primer cuadro. Esto es ya que a cada cuadro se le asigna un
estado que se compone de la posición, orientación, velocidad y biases
del giroscopio y acelerómetro que se estimaron para tal cuadro. Notar
que en Basalt, hablar de cuadros es equivalente a hablar de instantes de
tiempo, ya que los únicos puntos en el tiempo para los cuales el sistema
produce una pose estimada son las timestamps del par estéreo de imágenes
recibidas.

Para inicializar el primer estado se toman varias suposiciones. En
particular, se asume que el dispositivo comienza en la posición
\((0, 0, 0)\) y sin aceleración ni velocidad, esto permite utilizar el
vector de aceleración reportado por la muestra del acelerómetro como el
vector de gravedad y computar así la inclinación del agente. Notar que
esta inclinación no es capaz de informar la orientación de forma
completa al no poder contemplar uno de los ejes de rotación del cuerpo.
Por esta razón es recomendable iniciar la corrida con el agente rotado
con su eje vertical paralelo al vector gravedad, esto hará que Basalt
compute la orientación identidad. Tal rotación suele corresponder con la
posición de reposo pensada por el fabricante, o al menos este ha sido el
caso en los dispositivos utilizados en este trabajo. Además, es
conveniente posicionar el agente mirando hacia ``adelante'' ajustando el
\emph{yaw} \footnote{Los términos \emph{roll}, \emph{pitch} y \emph{yaw}
  provenientes de la aeronáutica son muy utilizados para hablar de la
  orientación de un cuerpo:
  \url{https://en.wikipedia.org/wiki/Aircraft_principal_axes}}, como el
usuario considere apropiado según su entorno.

En instantes posteriores, o sea al recibir nuevas imágenes, se realiza
la llamada preintegración de muestras consecutivas de la IMU.
Considerando que estas muestras arriban a mayores frecuencias que las de
las cámaras, preintegrarlas es un proceso que intenta resumir las
muestras entre los cuadros a una única pseudo-muestra que sucede en los
mismos instantes de tiempo que los cuadros como se ve en el ejemplo de
la \figref{fig:sample-frequencies}.

\fig{fig:sample-frequencies}{source/figures/sample-frequencies.pdf}{%
Ejemplo de frecuencias}{%
Frecuencia de distintos eventos para un ejemplo con cámaras a 30 fps y muestras
de la IMU a 150 Hz.
}

El proceso de preintegración es el siguiente. Dado el cuadro previo
\(i\) con timestamp \(t_i\) y el cuadro posterior \(j\) con timestamp
\(t_j\), se intenta computar una pseudo-muestra
\(\Delta \mathbf{s} = (\Delta \mathbf{R}, \Delta \mathbf{v}, \Delta \mathbf{p}) \in SO(3) \times \R^3 \times \R^3\)
que representa cambios de orientación, velocidad y posición
respectivamente según las mediciones de la IMU que ocurrieron desde
\(t_i\) hasta \(t_j\). Para cada timestamp \(t\) de la IMU tal que
\(t_i < t \leq t_j\) tenemos la muestra de aceleración lineal
\(\mathbf{a}_t\) y de velocidad angular \(\mathbf{\upomega}_t\).
Definimos entonces de forma recursiva la pseudo-muestra
\(\Delta \mathbf{s}\) de la siguiente manera:

\begin{align}
\label{eq:imu-preintegration}
(\Delta \mathbf{R}_{t_i}, \Delta \mathbf{v}_{t_i}, \Delta
\mathbf{p}_{t_i}) & := (\mathbf{I}, \mathbf{0}, \mathbf{0})
\\
\Delta \mathbf{R}_{t+1} & := \Delta \mathbf{R}_t exp(\mathbf{\upomega}_{t+1} \Delta t)
\\
\Delta \mathbf{v}_{t+1} & := \Delta \mathbf{v}_t + \Delta{\mathbf{R}_t}
\mathbf{a}_{t+1} \Delta t
\\
\Delta \mathbf{p}_{t+1} & := \Delta \mathbf{p}_t + \Delta \mathbf{v}_t \Delta t
\\
\Delta \mathbf{s}_{t} & := (\Delta \mathbf{R}_{t}, \Delta \mathbf{v}_{t}, \Delta
\mathbf{p}_{t})
\\
\Delta \mathbf{s} & := \Delta \mathbf{s}_{t_j}
\end{align}

Es destacable mencionar que este tipo de preintegración es también
utilizado por los otros sistemas estudiados Kimera y ORB-SLAM3. La
ventaja que presenta es que sus características son bien conocidas
gracias al trabajo de
\textcite{forsterOnManifoldPreintegrationRealTime2017} y las expresiones
necesarias para el cómputo de residuales, como sus jacobianos, son
cerradas y fueron derivadas de forma ejemplar en dicho trabajo.

Entonces, con esta muestra preintegrada junto a los datos del nuevo
cuadro (sus keypoints), se procede a la etapa de
\passthrough{\lstinline!measure!} del módulo de VIO. Aquí lo primero que
se hace es predecir que el estado de este nuevo instante estará basado
en el estado del instante anterior más la adición de las muestras
preintegradas de la IMU. El resto de la etapa de
\passthrough{\lstinline!measure!} se basa en el manejo y actualización
de la base de datos de los puntos de interés en 3D, o \emph{landmarks},
y sus \emph{observaciones}, junto a algo que, en Basalt, está
fuertemente ligado: la toma de cuadros clave, o \emph{keyframes}.

\hypertarget{base-de-datos-de-landmarks}{%
\subsubsection{Base de datos de
landmarks}\label{base-de-datos-de-landmarks}}

Recordemos que el módulo de optical flow encuentra keypoints en cada
cuadro, esto es, una landmark o punto de interés en la escena 3D
proyectada sobre el plano de la imagen 2D. Más aún este módulo era capaz
de hacer el seguimiento de keypoints similares mediante optical flow, es
decir, de keypoints que observan la misma landmark. Parte de nuestro
objetivo entonces será triangular las posiciones de estas landmarks
considerando las observaciones tomadas. Consideremos además la
naturaleza altamente ruidosa de estas observaciones, con landmarks que
aparecen y desaparecen de la visión de los cuadros por múltiples razones
como: ser ocluidas por objetos en el entorno, ser distorsionadas por el
ángulo del observador, distorsiones inherentes de los sensores
ópticos\marginnote{%\
Detallaremos en la \Cref{sec:data-characteristics} este tipo de distorsiones.
} como el motion blur, la sobreexposición, el ruido introducido por la
ganancia del amplificador de señal digital, o simplemente porque dejan
de estar en el campo de visión de las cámaras. Por estas razones
entonces, será fundamental la correcta gestión de la información de las
landmarks y sus observaciones. En Basalt, la clase que se encarga de
esto es la \passthrough{\lstinline!LandmarkDatabase!} con la estructura
como se define en el \Cref{lst:basalt-defs}.

\begin{lstlisting}[language={C++}, caption={Estructuras de Basalt}, label=lst:basalt-defs]
class LandmarkDatabase {
  map<LandmarkId, Landmark> landmarks;
  map<FrameId, map<FrameId, Keypoint>> observations;
}

class Landmark {
 FrameId keyframe;
 Vector2 direction;
 double inverse_distance;
}

class Keypoint {
  LandmarkId landmark_id;
  Vector2 position;
}
\end{lstlisting}

En la implementación de este Módulo, los términos de keypoint y landmark
tienden a utilizarse de forma intercambiable, significando keypoint algo
distinto a lo que era en el módulo de optical flow. Para evitar
confusión, se han diferenciado los términos de manera explícita en el
pseudo código anterior. \passthrough{\lstinline!Landmark!} se utilizará
para referirse al punto de interés en la escena tridimensional, mientras
que \passthrough{\lstinline!Keypoint!} será un punto 2D que observa la
proyección de una landmark definida por
\passthrough{\lstinline!landmark\_id!}. Entonces la
\passthrough{\lstinline!LandmarkDatabase!} es simplemente una colección
de todas las \passthrough{\lstinline!landmarks!} de interés y de los
keypoints que la observaron en \passthrough{\lstinline!observations!}.

Una \passthrough{\lstinline!Landmark!} surge de un keypoint observado en
el módulo de optical flow en algún cuadro y comparte el mismo
identificador. Explicaremos los campos
\passthrough{\lstinline!direction!} e
\passthrough{\lstinline!inverse\_distance!} que determinan la posición
en 3D de la misma más adelante. Un cuadro se identifica tanto por la
timestamp en el que fue tomado como por cuál de las dos cámaras lo tomó;
\passthrough{\lstinline!FrameId!} tiene esta información. Una landmark
existe en referencia a un cuadro, y un cuadro que es referenciado por
landmarks debe ser un keyframe por la implementación; decimos que el
keyframe aloja estas landmarks. Una observación en
\passthrough{\lstinline!observations!} tiene entonces un mapeo de
keyframes a los \passthrough{\lstinline!Keypoints!} que observan a las
landmarks alojadas en el keyframe. Un \passthrough{\lstinline!Keypoint!}
es reconocido en un cuadro particular, y solo tiene sentido como
coordenadas en ese cuadro, es esto lo que representa el segundo
\passthrough{\lstinline!FrameId!} de la definición de
\passthrough{\lstinline!observations!}.

Habiendo dicho esto, al recibir un nuevo cuadro,
\passthrough{\lstinline!measure!} simplemente recorre todas las
observaciones o \passthrough{\lstinline!Keypoint!}s que este trae y se
añaden a la base de datos las observaciones de landmarks ya existente en
la misma. Observaciones de landmarks no registradas en la base de datos
se guardan para poder determinar si el módulo amerita la toma de un
nuevo keyframe.

\hypertarget{keyframes}{%
\subsubsection{Keyframes}\label{keyframes}}

En contraste con sistemas como Kimera y ORB-SLAM3 que tienen condiciones
más intrincadas, en Basalt, la heurística para decidir si el cuadro
actual será un keyframe es muy sencilla: si \emph{más del 30\% de las
observaciones del cuadro actual corresponden a landmarks no registradas
y han pasado más de, por defecto, 5 cuadros consecutivos que no fueron
keyframes}, el cuadro actual será tomado como un keyframe. La toma de
keyframes es lo que registra nuevas landmarks a la base de datos y
realiza la estimación inicial de su posición.

En primer lugar, todas las observaciones \emph{``desconectadas''}
correspondientes al 30\% o más anterior que no correspondían a ninguna
landmark, intentarán ser registradas en la base de datos. Para esto es
necesario poder encontrar una segunda observación con la que realizar la
triangulación y poder así estimar la posición tridimensional de la
landmark al registrarla.

Se tiene entonces que para cada observación desconectada
\(\mathbf{p}_h \in \R^2\) producida por la cámara \(h\) del keyframe se
recorrerán todos los cuadros desde el último keyframe, buscando por una
segunda observación \(\mathbf{p}_t \in \R^2\) de esta landmark producida
en alguna cámara \(t\). En caso de encontrarla se continúa de la
siguiente manera:\newline

\begin{itemize}
\item
  La cámara \(h\) del keyframe tiene una pose estimada para la IMU en
  esa timestamp que denominaremos \(\mathbf{T}_{i_h} \in SE(3)\).
  Similarmente tendremos \(\mathbf{T}_{i_t}\) para la cámara \(t\).
\item
  A su vez como los parámetros de calibración son estáticos y conocidos,
  tenemos la función de proyección \(\pi_h\), sus parámetros intrínsecos
  \(\mathbf{i}_h\) y la pose relativa \(\mathbf{T}_{i_h c_h} \in SE(3)\)
  de la cámara \(h\) respecto a la IMU. Similarmente con \(\pi_t\),
  \(\mathbf{i}_t\) y \(\mathbf{T}_{i_t c_t}\) para la cámara \(t\).
\item
  Podemos ahora desproyectar \(\mathbf{p}_h\) y \(\mathbf{p}_t\) a sus
  respectivas posiciones 3D estimadas
  \(\mathbf{p}'_h := \pi^{-1}_h(\mathbf{p}_h, \mathbf{i}_h)\) y
  \(\mathbf{p}'_t := \pi^{-1}_t(\mathbf{p}_t, \mathbf{i}_t)\) con
  \(\mathbf{p}'_h, \mathbf{p}'_t \in \R^3\).
\item
  Computamos ahora la transformación de la IMU de \(h\) a la IMU de
  \(t\) dada por
  \(\mathbf{T}_{i_h i_t} := \mathbf{T}_{i_h}^{-1} \mathbf{T}_{i_t}\)
\item
  Luego podemos tener la transformación de la cámara \(h\) a \(t\) con
  \(\mathbf{T}_{c_h c_t} := \mathbf{T}_{i_h c_h}^{-1} \mathbf{T}_{i_h i_t} \mathbf{T}_{i_t c_t}\)
\item
  Finalmente con estos tres datos \(\mathbf{p}'_h\), \(\mathbf{p}'_t\) y
  \(\mathbf{T}_{c_h c_t}\) es posible realizar la triangulación
  utilizando DLT con SVD como se explicó en la sección ``DLT con SVD''
  para obtener un punto en coordenadas homogéneas en \(\R^4\) con el
  cuarto componente representando el inverso de la distancia entre la
  cámara y el punto de interés, mientras que los tres primeros
  componentes corresponden al \emph{bearing vector}, es decir, un vector
  unitario que determina la dirección desde la cámara hacia la landmark.
\end{itemize}

\begin{mdframed}[backgroundcolor=shadecolor]
TODO: En el último punto hablo de una sección “DLT con SVD” que no existe en el
trabajo.
Estoy todavía viendo si agregarla o simplemente referenciar el capítulo del libro
de Multiple View Geometry de la bibliografía.

DLT es direct linear transform y SVD es single value decomposition.
Esto se usa para triangular puntos 3D dadas dos vistas 2D.
Pero los detalles no necesité profundizarlos.
\end{mdframed}

Un detalle a considerar es que Basalt comprueba que la distancia
relativa entre los dos cuadros, la baseline, sea lo suficientemente
grande para considerar la triangulación exitosa, de lo contrario se
busca un nuevo cuadro para comparar con el keyframe. Esta comprobación
se reduce a revisar que la norma del vector de traslación contenido en
\(\mathbf{T}_{c_h c_t}\) sea de, por defecto, más de 5 cm.

Como se ve en la definición de \passthrough{\lstinline!Landmark!}, la
posición 3D de estos puntos de interés no se almacena exactamente de la
misma forma que la triangulación los produce. En particular, no se
almacena el bearing vector directamente, sino que se utiliza un punto 2D
más compacto \passthrough{\lstinline!direction!} que lo codifica (para
esto se utiliza una \emph{proyección estereográfica} como se explica en
la \figref{fig:stereographic-projection}) junto a
\passthrough{\lstinline!inverse\_distance!}, la distancia inversa a este
punto producto de la triangulación, de esta forma la posición de la
landmark queda ligada al keyframe que la aloja.

\fig{fig:stereographic-projection}{source/figures/stereographic-projection.png}{Proyección estereográfica}{%
Interpretación geométrica de la proyección estereográfica utilizada para
representar bearing vectors. Las coordenadas definidas por la propiedad \mono{Vector2
direction} en \mono{Landmark} definen un punto en el plano $XY$ ($Z=0$) mostrado en azul. Para
obtener el vector unitario correspondiente, se traza una línea desde el punto
$(0,\ 0,\ -1)^T$ hacia \mono{direction} en el plano $XY$. El vector en el que esta línea
interseca a la esfera unitaria será el bearing vector codificado. Se muestran
tres ejemplos en rojo, verde y amarillo, con lineas punteadas que representan
las líneas trazadas y flechas representando los bearing vectors obtenidos.
}

Si todos los procedimientos relacionados a la triangulación de estos dos
cuadros fueron correctos, se almacena la landmark nueva en la base de
datos. De haber otras observaciones de esta landmark no se utilizan
todavía para añadir información a su posición, si no que simplemente se
añaden las observaciones para uso futuro.

\hypertarget{optimizaciuxf3n-y-marginalizaciuxf3n-todo}{%
\subsection{Optimización y marginalización
(TODO)}\label{optimizaciuxf3n-y-marginalizaciuxf3n-todo}}

\begin{mdframed}[backgroundcolor=shadecolor]
TODO: Esta última parte es un poco
compleja y me ha estado costando terminar de cerrar como explicarla. Basicamente
lo que se hace es plantear la función de error a optimizar $E(x)$ que combina un
montón de términos de error. Y se le quiere aplicar gauss newton. El problema es
que al haber tantos términos que optimizar, el cálculo de los jacobianos se hace muy rebuscado
y se arma un callstack de funciones muy profundo que lo único que hace es
computar valores parciales de estos jacobianos.\\
\\
Por otro lado la idea de
“marginalización“ es importante por que habla de uno de los términos que aparece
en $E(x)$ y que es un poco la novedad que trae Basalt, ellos logran eliminar
(marginalizar) los efectos de muchas mediciones en unas pocas distribuciones
aproximadas.\\
\\
Para sumar a la complejidad de esta sección, está, un poco escondido en la
implementación, el concepto de “grafo de factores“ para armar antes de la
optimización. No lo he profundizado lo suficiente como para explicarlo pero a
grandes rasgos, es un grafo en donde se tienen nodos que son variables
aleatorias mientras que los lados son las mediciones que uno toma (los factores). Por ejemplo
un nodo puede ser la variable que representa la posición de una landmark que
quiero estimar,
mientras que un factor puede ser el vector distancia que el agente trianguló con
sus cámaras hacia esa landmark (también tiene ruido).\\
La idea de estos grafos es hacer una estimación
“máxima a posteriori“ (MAP) con todo lo que tienen adentro. MAP es un concepto
muy parecido al de estimación de máxima verosimilitud. Básicamente lo que se
tiene es que minimizar la función de error $E(x)$ lo que está haciendo es
acomodando los resultados de las variables de interés para que sean los de mayor
probabilidad (“de mayor verosimilitud“) dadas las restricciones impuestas por las mediciones tomadas (los
factores del grafo).
\end{mdframed}

\cleardoublepage
\ctparttext{Describiremos algunas de las contribuciones clave que fueron
producto de este trabajo. En el proceso se obtendrá un mejor entendimiento de
los problemas comunes de implementación por los que estos sistemas se ven
afectados. Además, veremos las soluciones que se les ha dado a otras
problemáticas que son propias de la localización en tiempo real aplicada a XR.}

\hypertarget{contribuciones}{%
\part{Contribuciones}\label{contribuciones}}

\hypertarget{tracking-por-slam-para-monado}{%
\chapter{Tracking por SLAM para
Monado}\label{tracking-por-slam-para-monado}}

\hypertarget{sec:thesis-context}{%
\section{Contexto}\label{sec:thesis-context}}

Por su naturaleza, el área de XR involucra una gran cantidad de partes
interconectadas y de dispositivos muy diversos con configuraciones
difíciles de generalizar, y más aún de predecir. Por esta razón hasta
hace muy poco tiempo no existían estándares razonables en el área lo
cual agravaba la situación con un ecosistema altamente fragmentado en
soluciones propietarias que causaban grandes problemas a los
desarrolladores de aplicaciones finales. En el mejor de los casos, la
carga de soportar los distintos SDK propietarios recaía sobre frameworks
como \emph{WebXR}\footnote{\url{https://www.w3.org/TR/webxr}} o motores
de juegos (p.~ej. \emph{Unreal Engine}, \emph{Unity},
\emph{Godot}\footnote{\url{https://www.unrealengine.com},
  \url{https://unity.com} y \url{https://godotengine.org}}) y esto
forzaba a los desarrolladores a elegir alguna de estas soluciones para
realizar su aplicación de XR. En caso de no querer hacerlo, se vería
obligado a realizar un esfuerzo significativo para portar su aplicación
a cada uno de estos SDK, y eso sin considerar el manejo de
características especiales que algunas plataformas exponen y otras no.
Este escenario se puede ver en la \figref{fig:before-openxr}.

\fig{fig:before-openxr}{source/figures/before-openxr.pdf}{Antes de OpenXR}{%
Antes de OpenXR las aplicaciones y motores necesitaban código
propietario separado para cada SDK de los dispositivos que quisieran soportar.
}

Luego de unos años de sufrir esta fragmentación, en julio de 2019 se
presenta la primera versión de \emph{OpenXR}
\autocite{thekhronosgroupinc.OpenXRSpecification} de la mano del
\emph{Khronos Group}\footnote{\url{https://www.khronos.org}}. Este es un
consorcio abierto y sin fines de lucro compuesto de, a la fecha, 170
organizaciones que desarrolla estándares en distintas áreas de la
industria como computación gráfica (\emph{OpenGL}, \emph{Vulkan}),
computación paralela (\emph{OpenCL}, \emph{SYCL}) y, ahora con OpenXR,
realidad virtual y aumentada entre otras. OpenXR provee una API
estandarizada con soporte para extensiones que permiten añadir
características peculiares de ser necesarias por algún fabricante en
particular. El estándar ha tenido un gran éxito al haber sido adoptado
por una gran cantidad de compañías \footnote{Compañías respaldando
  públicamente el estándar OpenXR:
  \url{https://www.khronos.org/assets/uploads/apis/2019-openxr-logo-field_1_15.jpg}}
como reemplazo a sus antiguos SDK propietarios. De esta forma, los
motores de juego y desarrolladores solo necesitan interactuar con una
única API que además les permite aprovechar cualquier característica
especial ofrecida por alguna extensión. Se puede ver la simplificación
del esquema de integración con OpenXR en la \figref{fig:after-openxr}.

\fig{fig:after-openxr}{source/figures/after-openxr.pdf}{OpenXR}{%
OpenXR provee una única interfaz (API) multiplataforma de alta performance entre las aplicaciones
y todos los dispositivos compatibles. Esto es una mejora en comparación a la
situación presentada en la \figref{fig:before-openxr}.
}

OpenXR es exclusivamente la especificación
\autocite{thekhronosgroupinc.OpenXRSpecification} de una API y por lo
tanto requiere una implementación, o \emph{runtime}, sobre el que
ejecutarse. Las implementaciones son provistas por los distintos
fabricantes interesados en soportar el estándar, en la imagen
referenciada en la nota al pie \footnote{Algunas de las compañías que
  implementan un runtime de OpenXR:
  \url{https://www.khronos.org/assets/uploads/apis/OpenXR-After_3.png}.}
se pueden ver algunas de las implementaciones ya desarrolladas. En este
trabajo nos enfocaremos en una de ellas, \emph{Monado}. Monado es un
runtime de la especificación OpenXR de código abierto, por ahora el
único con esta característica, licenciado bajo la \emph{Boost Software
License 1.0} \autocite{BoostSoftwareLicensea}. La plataforma principal
sobre la que Monado corre y se desarrolla es GNU/Linux, pero es capaz de
ser ejecutarse en otras como Android y Windows. Su desarrollo está
soportado por \emph{Collabora Ltd.}, quien es parte del grupo de trabajo
de OpenXR desde sus inicios. Las contribuciones a Monado listadas en
esta sección fueron realizadas durante una pasantía de seis meses
realizada por el autor en Collabora.

Además de proveer una implementación de OpenXR, Monado es altamente
modular e implementa distintos componentes reutilizables para XR como un
compositor especializado para realidad virtual; controladores para una
gran variedad de dispositivos, incluyendo hardware propietario y de
consumo masivo sobre los que la comunidad ha realizado ingeniería
inversa para poder utilizar; herramientas varias de calibración y
configuración de hardware; así como también distintos sistemas de fusión
de sensores para tracking; recientemente incluso incorpora un módulo de
localización de manos mediante visión por computadora y aprendizaje
automático.

Una característica faltante en Monado era la posibilidad de realizar
localización visual-inercial mediante sistemas de SLAM/VIO. Este tipo de
tracking ha cobrado gran popularidad en los últimos años por resultar
sumamente convenientes al no requerir sensores externos al dispositivo
de XR. Sistemas de este tipo son empleados en productos como el
\emph{Meta Quest}, los cascos \emph{Windows Mixed Reality} o incluso los
SDK \emph{ARCore} y \emph{ARKit} presentes en dispositivos móviles.
Desafortunadamente, todas estas soluciones son privativas y, por lo
tanto, no es posible obtener acceso a sus códigos fuentes para
reusarlos, modificarlos o simplemente estudiarlos sin obtener licencias
especiales de sus fabricantes. Más aún, existen compañías que se
especializan en desarrollar soluciones comerciales de SLAM como
\emph{SLAMCore\footnote{\url{https://www.slamcore.com/}}},
\emph{Arcturus\footnote{\url{https://arcturus.industries/}}} y
\emph{Spectacular AI\footnote{\url{https://www.spectacularai.com/}}}
entre otras.

Este trabajo se concentró entonces en el estudio de implementaciones de
código abierto de sistemas de tracking visual-inercial (ya sea mediante
SLAM o solamente VIO) y en la integración de estos sobre Monado. Se
necesitó armar la infraestructura para soportar una interacción modular
con los sistemas mediante el desarrollo de interfaces, herramientas,
controladores, y mejoras principalmente en Monado, pero también en los
sistemas a integrar o en sus \emph{forks} (clones específicos de las
implementaciones de SLAM/VIO para uso en Monado).

\hypertarget{monado}{%
\section{Monado}\label{monado}}

En este capítulo explicaremos los distintos módulos y funcionalidad
implementada en Monado para permitirle proveer tracking mediante
SLAM/VIO a distintos controladores de dispositivos y de esta forma, a
las aplicaciones OpenXR que utilizan estos dispositivos como medios de
interacción y corren sobre Monado. Cabe aclarar que a partir de ahora
usaremos de forma intercambiable los términos SLAM y VIO, ya que, a
pesar de ser distintos, para los fines prácticos de la implementación
son equivalentes: \emph{piezas de software que consumen imágenes y
muestras de IMU y devuelven poses estimadas como resultado}. Es por esto
que hablaremos de un localizador SLAM o \emph{SLAM tracker} que se ha
implementado para referirnos a una única funcionalidad que soporta tanto
sistemas externos de SLAM y como de VIO.

Antes de comenzar, vale la pena entender la arquitectura general de
Monado como se muestra en la \figref{fig:monado-arch}. Una aplicación
que hace uso de OpenXR mediante Monado puede ser corrida en dos
modalidades dependiendo de como se compiló el runtime. De la forma
tradicional, se compila como una biblioteca de manera \textbf{autónoma},
es decir que cuando la aplicación intenta cargar dinámicamente alguna
implementación de OpenXR provista por el sistema, se enlaza la
biblioteca compartida, o \emph{shared library}, de Monado, y al
interactuar con cualquier interfaz a OpenXR, se correría tal
procedimiento en el mismo hilo que la aplicación consumidora. Esta
modalidad puede tener ciertos beneficios para depuración de código, pero
el modo más usual de compilar y correr Monado es mediante
\emph{Inter-Process Communication} o \emph{IPC}. De esta manera el
runtime se corre en un proceso independiente que debe ser lanzado con
anterioridad a cualquier aplicación OpenXR. Lo que se termina enlazando
a la aplicación final es simplemente un cliente IPC que es capaz de
comunicarse con Monado. Esto ofrece ventajas como la posibilidad de
correr múltiples aplicaciones sobre la misma instancia del runtime.
Además, se quita la necesidad de ejecutar ambos procesos sobre el mismo
nodo de cómputo, lo cual facilita la posibilidad de tener realidad
virtual renderizada en la nube.

\fig{fig:monado-arch}{source/figures/monado-arch.pdf}{Arquitectura de Monado}{%
Arquitectura de Monado.
}

\FloatBarrier

No profundizaremos en los aspectos de comunicación con la
\emph{plataforma} ni del \emph{compositor} de Monado, pero es bueno
saber que estos son partes significativas del runtime. Son los
\emph{controladores} (o \emph{drivers}) específicos en Monado los que
interactúan con el hardware especializado como cascos, controles,
cámaras, entre otros dispositivos que quieran utilizarse para la
experiencia de XR. Sensores como las IMU y cámaras sobre los que
adquirir las muestras necesitarán ser implementados en este nivel.
Finalmente, hay una gran variedad de módulos auxiliares reutilizables
que exponen funcionalidades matemáticas, de interacción con el sistema
operativo, con OpenGL o Vulkan, entre otras. Es en estos módulos
auxiliares en donde encontramos componentes relacionados exclusivamente
al problema tracking y a tipos específicos de tracking. Se implementan
herramientas de calibración de cámaras, de depuración, filtros de
distinto tipo (p.~ej. \emph{Kalman} y \emph{low-pass}), entre otra
algoritmia genérica de fusión de sensores. Conforman este módulo también
sistemas completos de tracking para \emph{PlayStation Move},
\emph{PlayStation VR}, y tracking de manos en general. Es aquí entonces
en donde se comienza la implementación del SLAM tracker presentado en
este trabajo.

La implementación de un pipeline en Monado que permita la comunicación
entre dispositivos, sistemas de SLAM y la aplicación OpenXR requirió
desarrollar la infraestructura y herramientas necesarias dentro de
Monado. El pipeline en cuestión está esquematizado en la
\figref{fig:slam-tracker-dataflow}. Este, requiere un gran nivel de
modularidad, ya que sus componentes fundamentales necesitan poder ser
intercambiables: dispositivos, aplicaciones y el sistema que provee el
SLAM en sí. El resto de esta subsección está dedicada a explicar la
infraestructura y los distintos componentes que se necesitaron
implementar y adaptar para obtener un pipeline de SLAM modular corriendo
en Monado como se muestra en la \figref{fig:slam-tracker-dataflow}.

\fig{fig:slam-tracker-dataflow}{source/figures/slam-tracker-dataflow.pdf}{Flujo de datos de la implementación}{%
Diagrama esquemático de como ocurre el flujo de los datos desde que se generan
las muestras en los dispositivos XR hasta que la aplicación OpenXR obtiene una
pose utilizable. Notar que las flechas esconden múltiples hilos de ejecución,
colas de procesamiento, y desfasajes temporales con los que hay que lidiar en la
implementación.
}

\FloatBarrier

\hypertarget{interfaz-externa}{%
\section{Interfaz externa}\label{interfaz-externa}}

Desde un principio se entendió que se necesitaría utilizar sistemas ya
desarrollados como punto de partida. Estos sistemas son complejos y
suelen utilizar conceptos teóricos de significativa profundidad, por lo
que su creación suele estar limitado a grupos de investigación expertos
que toman gran tiempo de desarrollo. Los tres sistemas estudiados por
ejemplo, promedian las 25.000 líneas de código (o 25 \emph{KLOC}) cada
uno.

Ahora bien, en muchos casos, haber intentado integrar el código del
sistema directamente dentro de un componente de Monado no era una
opción. Dejando de lado las dificultades técnicas, los problemas de
compatibilidad de licencia fueron de particular interés. La gran mayoría
de sistemas SLAM producidos en la academia son liberados bajo licencias
abiertas ``virales'' como \emph{GPL} \autocite{GNUGeneralPublic} que
obligan a desarrolladores que utilizan código del sistema a liberar y
licenciar su código de la misma manera. Esto contrasta con la licencia
abierta y permisiva de Monado, la \emph{BSL-1.0}
\autocite{BoostSoftwareLicensea}, que no impone restricciones sobre como
los usuarios deben licenciar su código.

Estas fueron algunas razones para intentar desacoplar el sistema a
utilizar lo más que se pueda de Monado. Además, considerando la
naturaleza experimental de este trabajo, la posibilidad de que más de un
sistema necesitase ser integrado era razonable.

Monado está desarrollado principalmente en C, pero gran parte de su
código de tracking está implementado en C++ al igual que todos los
sistemas de SLAM contemplados. Adicionalmente tanto Monado como estos
sistemas suelen hacer un uso extensivo de la biblioteca \emph{OpenCV}, y
en particular su clase contenedora de imágenes y matrices
\passthrough{\lstinline!cv::Mat!}. Es por esto que se terminó optando
por el uso de un archivo \emph{header} C++, en el cual se declara la
clase \passthrough{\lstinline!slam\_tracker!} que será utilizada por
Monado como punto de comunicación con sistemas de SLAM arbitrarios y se
utilizan \passthrough{\lstinline!cv::Mat!} como contenedor de imágenes.
Luego de varias iteraciones de diseño, la clase
\passthrough{\lstinline!slam\_tracker!} tiene una interfaz que, quitando
detalles de tipos de C++, se puede resumir en algo como lo que se
muestra en el \Cref{lst:slam-tracker-def}.

\clearpage

\begin{lstlisting}[language={C++}, caption={Interfaz a implementar por sistemas de SLAM}, label=lst:slam-tracker-def]
class slam_tracker {
public:
  // (1) Constructor y funciones de inicio/fin
  slam_tracker(string config_file);
  void start();
  void stop();

  // (2) Métodos principales de la interfaz
  void push_imu_sample(timestamp t, vec3 accelerometer, vec3 gyroscope);
  void push_frame(timestamp t, cv::Mat frame, bool is_left);
  bool try_dequeue_pose(timestamp &t, vec3 &position, quat &rotation);

  // (3) Características dinámicas opcionales
  bool supports_feature(int feature_id);
  void* use_feature(int feature_id, void* params);

private:
  // (4) Puntero a la implementación (patrón PIMPL)
  void* impl;
}
\end{lstlisting}

Este header está presente en Monado, pero su implementación no. Esta
debe ser provista por el sistema externo, lo cual implica tener que
mantener una copia, o \emph{fork}, levemente modificado de los distintos
sistemas que se quieran utilizar, ver \figref{fig:slam-tracker-hpp}.

\fig{fig:slam-tracker-hpp}{source/figures/slam-tracker-hpp.pdf}{Interfaz de SLAM tracker}{%
Interacción entre Monado y sistemas SLAM mediante la interfaz en C++.
}

La versión actual de esta clase es el resultado de varias iteraciones y
generaliza adecuadamente los tres sistemas en uso. Algunas
consideraciones de los puntos marcados en el código:

\begin{enumerate}
\def\labelenumi{\arabic{enumi}.}
\item
  El parámetro \passthrough{\lstinline!config\_file!} del constructor es
  necesario, ya que todos los sistemas con los que se trató requieren
  proveer información de calibración y puesta a punto de parámetros
  previo a la corrida mediante un archivo de configuración. Además estos
  sistemas suelen tener etapas de creación e inicialización de recursos,
  así como de liberación de los mismos que quedan representados en el
  par de métodos
  \passthrough{\lstinline!start()!}/\passthrough{\lstinline!stop()!}.
\item
  Los sistemas corren en hilos separados de Monado y es por esto que es
  fundamental implementar colas concurrentes a la hora de intercambiar
  datos. Monado ingresa muestras mediante los métodos
  \passthrough{\lstinline!push\_imu\_sample!} y
  \passthrough{\lstinline!push\_frame!} mientras que sondea si hay poses
  ya estimadas por el sistema y las obtiene mediante
  \passthrough{\lstinline!try\_dequeue\_pose!}.
\item
  Algo que surge del desarrollo de una interfaz a un tipo de sistemas
  que aún se desconocen, es que va a haber varios cambios en la misma
  durante su creación. Si además esta interfaz es compartida por
  múltiples sistemas y repositorios, mantener todas las versiones
  sincronizadas se vuelve insostenible. Una forma de aliviar este
  problema fue la implementación de características dinámicas. En ellas,
  Monado evalúa si el sistema implementa alguna característica
  específica en tiempo de ejecución antes de utilizarla. Uno de sus
  usos, fue la automatización del envío de datos de calibración sin
  pasar por el archivo \passthrough{\lstinline!config\_file!}. La forma
  de añadir nuevas características de este tipo se debe reservar un
  entero \passthrough{\lstinline!feature\_id!} que la identifique en una
  nueva versión del header \passthrough{\lstinline!slam\_tracker!} de
  Monado y del fork que la va a implementar. En Monado tenemos cuidado
  de solo utilizarla si el sistema la reporta como disponible en tiempo
  de ejecución, ya que otros forks podrían no implementarla. De esta
  forma permitimos la extensión de sistemas específicos sin tener que
  adaptar la versión de la interfaz en todos ellos.
\item
  El miembro \passthrough{\lstinline!impl!} es una forma de ligar la
  definición compacta de \passthrough{\lstinline!slam\_tracker!} con una
  clase que implementa el estado y métodos privados necesarios para
  proveer la funcionalidad requerida. Este patrón de desarrollo es
  usualmente conocido como \emph{pointer to implementation} (o
  \emph{PIMPL}), a \passthrough{\lstinline!impl!} se lo denomina un
  \emph{puntero opaco} \autocite{lakosLargeScaleSoftwareDesign1996}.
\end{enumerate}

Esta interfaz no es perfecta: no contempla magnetómetros, asume una
configuración de a lo sumo dos cámaras, asume que el sistema utiliza
OpenCV y es difícil de extender con cambios no contemplados por el
concepto de características dinámicas. A pesar de estos problemas, ha
sido suficientemente buena para generalizar todos los sistemas
propuestos y correrlos con una performance adecuada \marginnote{%\
ILLIXR \autocite{HuzaifaDesai2020} es un proyecto de XR de código libre desarrollado
por la Universidad de Illinois, que también formó recientemente el Consorcio ILLIXR
para intentar mejorar el ecosistema de código libre para XR. Ha habido discusiones entre ILLIXR y Monado
para plantear una interfaz a sistemas de SLAM que le sirva a ambos proyectos y que quizás en un
futuro pueda transformarse en un estándar que los sistemas implementen por elección propia.
Para esta nueva interfaz, se tomarán ideas basadas en la experiencia obtenida desarrollando el header \mono{slam_tracker}.
}.

\FloatBarrier

\hypertarget{implementaciones-de-la-interfaz}{%
\section{Implementaciones de la
interfaz}\label{implementaciones-de-la-interfaz}}

A la hora de implementar la interfaz
\passthrough{\lstinline!slam\_tracker!} se documentan los tres métodos
principales \passthrough{\lstinline!push\_imu\_sample!},
\passthrough{\lstinline!push\_frame!} y
\passthrough{\lstinline!try\_dequeue\_pose!} con las precondiciones que
el usuario de la interfaz, Monado en este caso, y su implementación
deben cumplir:

\begin{enumerate}
\def\labelenumi{\arabic{enumi}.}
\tightlist
\item
  Debe haber un único hilo productor llamando a los métodos
  \passthrough{\lstinline!push\_*!}
\item
  Las muestras deben tener marcas de tiempo (\emph{timestamps})
  monótonas crecientes.
\item
  La implementación de los métodos \passthrough{\lstinline!push\_*!} no
  debe ser bloqueante.
\item
  Del punto anterior se desprende que las muestras deben ser procesadas
  en un hilo consumidor separado.
\item
  Las muestras de imágenes estéreo deben tener la misma timestamp.
\item
  Las muestras de imágenes estéreo deben ser enviadas de forma
  intercalada en orden izquierda-derecha.
\item
  Debe haber un único hilo consumidor llamando a
  \passthrough{\lstinline!try\_dequeue\_pose!}.
\end{enumerate}

Estas condiciones fueron desarrollándose mientras la interfaz
evolucionaba y nuevos sistemas se iban adaptando. Intentan ser
indicaciones que evitarían implementaciones deficientes mientras que son
lo suficientemente generales como para que todas las soluciones
estudiadas puedan cumplirlas.

La precondición \(\mathbf{1}\) para el usuario le permite a la
implementación utilizar colas enfocadas al caso de un único productor,
de las cuales puede obtenerse mayor rendimiento comparado a las
versiones que no tienen esta limitación. Más aún, se suelen utilizar
colas libres de \emph{locks} (primitiva de sincronización, también
conocido como \emph{mutex}). El punto \(\mathbf{2}\) facilita el trabajo
a la implementación mientras que suele ser trivial de cumplir para el
usuario, ya que los dispositivos tienden a reportar las muestras del
mismo tipo en orden temporal creciente. El punto \(\mathbf{3}\) obliga a
la implementación de los métodos \passthrough{\lstinline!push\_*!} a
terminar lo antes posible sin realizar ningún tipo de procesamiento en
esas funciones. Esto es fundamental, ya que Monado también recibe
muestras de dispositivos mediante colas y simplemente las redirige al
\passthrough{\lstinline!slam\_tracker!}, y demorarse en estos métodos
hace que las colas internas de Monado se llenen y muestras se pierdan.
El inciso \(\mathbf{4}\) aclara en dónde se procesarían las muestras;
dado que no pueden procesarse en los métodos
\passthrough{\lstinline!push\_*!} se necesita un hilo consumidor que
esté esperando por muestras nuevas en las distintas colas de la
implementación para procesar. El punto \(\mathbf{5}\) es un
requerimiento razonable de algunos sistemas, ya que es común que cámaras
estéreo pensadas para SLAM tengan sincronización por hardware de sus
timestamps. El punto \(\mathbf{6}\) facilita algunos procedimientos y
chequeos en la implementación. Y finalmente, la precondición
\(\mathbf{7}\) es similar al punto 1 en donde se le permite a la
implementación utilizar colas de un único consumidor que sean de mayor
rendimiento comparado a otras colas concurrentes.

La interfaz \passthrough{\lstinline!slam\_tracker!} se implementó en los
tres forks presentados en este trabajo; Kimera, ORB-SLAM3 y Basalt. La
idea general de las implementaciones es similar: permitir a Monado, el
usuario de la interfaz, utilizar el sistema de SLAM/VIO realizando las
conversiones adecuadas para poder comunicarse con la solución
subyacente. Sin embargo hay algunas consideraciones a remarcar en cada
versión.

Respecto al manejo de muestras entrantes, Kimera y Basalt encolan las
muestras directamente a las colas de entrada de sus respectivos
pipelines que corren en hilos separados. Para ORB-SLAM3, al ser
necesario ejecutar explícitamente el paso de estimación con el par de
imágenes estéreo y muestras IMU nuevas, se necesitó implementar colas
dedicadas para muestras de cámara izquierda, derecha e IMU, junto a un
hilo consumidor que ejecuta la estimación cuando llegan cuadros nuevos.
Cabe mencionar que Basalt no utiliza el tipo
\passthrough{\lstinline!cv::Mat!} de OpenCV como contenedor para sus
imágenes a diferencia de los otros sistemas. Más aún, espera imágenes
con intensidad de píxeles de 16 bits para soportar conjuntos de datos
como TUM-VI que provee imágenes de este tipo. Sin embargo, las imágenes
monocromáticas usadas comúnmente en estas aplicaciones suelen utilizar
un rango de 8 bits y por ende los dispositivos generan imágenes de este
tipo. Al no haber una forma sencilla de utilizar 8 bits en Basalt, es
necesario realizar una copia a un tipo de imagen de 16 bits cada vez que
se presentan nuevos cuadros. Este es un punto a mejorar en el futuro
sobre Basalt que no se trató de resolver en este trabajo, ya que no se
encontraron problemas de performance incluso con esta copia innecesaria.

Respecto a la cola de poses estimadas, Kimera implementa un sistema de
\emph{callbacks} que llaman a una función encargada de encolar los
resultados del pipeline tan pronto como estén listos. Basalt por otra
parte, expone una cola que se va llenando con los resultados de la
estimación; la implementación de \passthrough{\lstinline!slam\_tracker!}
la utiliza directamente, ya que también expone una interfaz de cola
mediante el método \passthrough{\lstinline!try\_dequeue\_pose!}.
ORB-SLAM3 difiere de los dos anteriores ya que el procesamiento sucede
en una llamada a una función bloqueante que al terminar devuelve la pose
estimada. Por lo tanto, nuestra implementación se encarga de armar la
maquinaria necesaria de hilos productores, hilos consumidores y colas de
entrada y salida para hacer que Monado no sea bloqueado al enviar nuevas
muestras de los dispositivos.

Respecto al archivo de configuración referenciado por el parámetro
\passthrough{\lstinline!config\_file!}, Basalt por defecto utiliza dos
archivos, uno para parámetros de calibración y el otro para
configuraciones del sistema en sí; se necesitó entonces crear un tipo de
archivo nuevo que referencie a estos dos y sea utilizado como
\passthrough{\lstinline!config\_file!}. La misma idea se utilizó para
Kimera que también posee múltiples archivos de configuración. ORB-SLAM3
por su parte utiliza un único archivo por defecto y no necesito de
mayores cambios.

Respecto a las interfaces gráficas, todos estos sistemas presentan la
posibilidad de visualizar, al menos, la trayectoria estimada junto a las
features detectadas en los cuadros entrantes como se muestra en la
\figref{fig:trackers-ui}. En todos los casos, siempre se intenta
permitir la posibilidad de utilizar estas herramientas de visualización
de forma opcional al utilizar la clase
\passthrough{\lstinline!slam\_tracker!}. Para Kimera y ORB-SLAM3 las
interfaces funcionan automáticamente al habilitarlas en los archivos de
configuración, mientras que Basalt utiliza sus herramientas de
visualización de manera más ad hoc y por lo tanto fue necesario
implementar una clase de visualización dedicada para el
\passthrough{\lstinline!slam\_tracker!}.

\fig{fig:trackers-ui}{source/figures/trackers-ui.pdf}{Visualizadores de SLAM trackers}{%
Las distintas interfaces gráficas y formas de visualizar presentadas por cada
uno de los sistemas adaptados. De arriba a abajo: Kimera, ORB-SLAM3 y Basalt.
}

\FloatBarrier

\hypertarget{clase-adaptadora}{%
\section{Clase adaptadora}\label{clase-adaptadora}}

Para interactuar con la interfaz que detallamos en la sección anterior,
se implementa en Monado una clase adaptadora en el módulo auxiliar de
tracking (recordar \figref{fig:monado-arch}) y que sirve de nexo entre
los controladores de dispositivos de hardware y la interfaz
\passthrough{\lstinline!slam\_tracker!}. Este adaptador recibe el
nombre, un poco confuso, de \passthrough{\lstinline!TrackerSlam!}
siguiendo las convenciones de Monado para con los trackers ya
existentes.

El funcionamiento de \passthrough{\lstinline!TrackerSlam!} es sencillo,
los controladores que quieran ser localizados por SLAM y puedan proveer
imágenes y muestras de IMU deben instanciar este adaptador e
inicializarlo. Por detrás, esto simplemente llama a los métodos
adecuados de la interfaz \passthrough{\lstinline!slam\_tracker!} con el
sistema externo. Ahora bien, se aprovecha esta clase
adaptadora\marginnote{%\
Es debatible si el añadido de estas funcionalidades
haría que \mono{TrackerSlam} deje de ser considerada una clase adaptadora, ya que
como veremos, ambas pueden ser deshabilitadas en tiempo de ejecución.
} para proveer dos funcionalidades fundamentales que escapan al alcance
de los sistemas de SLAM/VIO y serán explicados a continuación:
\textbf{predicción} y \textbf{filtrado} de poses.

\hypertarget{predicciuxf3n-de-poses}{%
\subsection{Predicción de poses}\label{predicciuxf3n-de-poses}}

\hypertarget{el-problema}{%
\subsubsection{El problema}\label{el-problema}}

Las aplicaciones XR requieren poder localizar constantemente a los
distintos dispositivos de entrada y salida soportados que son utilizados
por el usuario. En la especificación de OpenXR
\autocite{thekhronosgroupinc.OpenXRSpecification} las dos principales
funciones que le permiten a la aplicación pedirle al runtime las poses
de estos dispositivos son \passthrough{\lstinline!xrLocateSpace!} y
\passthrough{\lstinline!xrLocateViews!}. La primera se utiliza para
solicitar poses de dispositivos ``comunes'' (suelen ser distintos tipos
de mandos) mientras que la segunda es un poco más compleja, ya que
concierne a la ubicación de las pantallas que renderizan la escena
(p.~ej. la ubicación de las pantallas de un casco VR). Para nuestro
propósito, podemos resumir las signaturas de ambas funciones a:

\begin{lstlisting}[language={C++}]
XrPosef xrLocateSpace(XrSpace id, XrTime time);
XrPosef xrLocateViews(XrSpace id, XrTime time);
\end{lstlisting}

Ambas devuelven una pose a un tiempo \passthrough{\lstinline!time!} para
el ``\emph{espacio}'' identificado por \passthrough{\lstinline!id!}. Un
espacio o \emph{espacio de referencia} es un término utilizado en la
especificación para diferenciar cualquier punto que nos interese
trackear desde la aplicación y que en definitiva identifica a un sistema
de referencia inercial con rotaciones.

Para obtener estos espacios, la aplicación solicita las características
que desea. Si se solicitara el espacio de un control o mando el runtime,
Monado en este caso, intentará conseguir el más adecuado dentro de los
disponibles en el sistema. Entonces, la aplicación OpenXR es indiferente
a qué dispositivos están siendo utilizados, ni siquiera se asumen que
estos espacios sean dispositivos, podrían ser cualquier otro objeto de
interés que está siendo localizado por mecanismos externos (por visión
por computadora por ejemplo) o incluso dispositivos emulados por
software\footnote{Una de las primeras contribuciones realizadas para
  familiarizarse con el código fuente de Monado fue la implementación
  del controlador \passthrough{\lstinline!qwerty!} que le permite a los
  usuarios emular de forma modular un casco y/o mandos mediante teclado
  y ratón.
  \url{https://gitlab.freedesktop.org/monado/monado/-/merge_requests/714}}.
Los espacios que aplican a nuestro caso son aquellos que representaran
dispositivos que posean sensores IMU y cámaras que puedan utilizarse en
nuestros sistemas de SLAM/VIO.

Otro importante aspecto a considerar es que el punto en el tiempo
\passthrough{\lstinline!time!} para la cual la pose debe ser estimada es
provisto por el usuario, y como tal, resulta arbitrario para el runtime.
Sin embargo los sistemas de SLAM/VIO suelen ser sistemas de tiempo
discreto, es decir solo dan estimaciones para puntos en el tiempo para
los cuales tienen muestras. En el mejor de los casos esto implica que
pueden dar una estimación por cada muestra de IMU, las cuales vienen a
altas frecuencias (p.~ej. 200hz); Kimera cumple con esto. En el caso más
usual sin embargo, y el que más nos afecta en este trabajo, las
estimaciones vienen a la misma frecuencia que los cuadros de las
cámaras. Esta frecuencia suele ser al menos un orden de magnitud menor
(p.~ej. 20hz); ORB-SLAM3 y Basalt funcionan de esta manera.

Para complicar más las cosas, los tiempos en los que el programador
solicita las poses suelen estar ligados a momentos en los cuales hay que
renderizar un nuevo cuadro en la pantalla del usuario. Al ser este un
proceso que (en el caso usual) ocurre enteramente en el mismo nodo de
cómputo, suelen ser tiempos muy cercanos al presente. Para el pipeline
de SLAM sin embargo, siempre nos encontramos levemente en el pasado por
tener que lidiar con las demoras inherentes a la captura de muestras,
las transmisiones de datos y los tiempos de cómputo significativos de
los sistemas de SLAM.

\hypertarget{anuxe1lisis-de-un-ejemplo}{%
\subsubsection{Análisis de un ejemplo}\label{anuxe1lisis-de-un-ejemplo}}

\fig{fig:prediction-timeline}{source/figures/prediction-timeline.png}{Línea de tiempo de predicción}{%
Línea de tiempo con timestamps normalizadas. Las barras representan las
siguientes duraciones:
\mono{REQUEST\_TO\_PREDICTION}: Del momento en que una predicción es solicitada hasta la timestamp de la predicción.
\mono{SHOT\_TO\_RECEIVED}: Captura de imagen en el dispositivo hasta su recepción en Monado.
\mono{RECEIVED\_TO\_PUSHED}: Transferencia de Monado a Basalt.
\mono{PUSHED\_TO\_PROCESSED}: Cómputo de la estimación de pose.
}

En la \figref{fig:prediction-timeline} se puede apreciar una captura de
pantalla de la interfaz de \emph{Perfetto}\footnote{Biblioteca de
  perfilación Perfetto: \url{https://perfetto.dev}} que, en conjunto con
\emph{Percetto}\footnote{Wrapper de Perfetto para C:
  \url{https://github.com/olvaffe/percetto}}, son las herramientas de
medición de tiempos preferidas para Monado. Esta captura es sobre una
corrida en tiempo real con Monado, Basalt y una cámara RealSense D455
con imágenes estéreo de resolución 640x480 a 30 cuadros por segundo. La
figura muestra un tramo de unos 35 ms con la particularidad de que el
tiempo en el que el usuario pide una predicción y el par estéreo de
imágenes son capturadas por la cámara coinciden en
\passthrough{\lstinline![A]!}. Como la tarea de renderizado para esta
captura ocurría a unos 60 cuadros por segundo, tenemos que
aproximadamente dos predicciones (en \passthrough{\lstinline![A]!},
\passthrough{\lstinline![D]!}) son solicitadas entre cada muestra de la
cámara (\passthrough{\lstinline![A]!} y \passthrough{\lstinline![G]!}).

\FloatBarrier

A tiempo \passthrough{\lstinline![C]!} las imágenes llegan al host luego
de haber sido transferidas por un cable USB 3.2 con una demora de unos
13,5 ms representada por la barra
\passthrough{\lstinline!SHOT\_TO\_RECEIVED!}. En ese momento ocurre una
pequeña copia de Monado hacia Basalt
\passthrough{\lstinline!RECEIVED\_TO\_PUSHED!}\marginnote{Esta copia es necesaria por un detalle en la
implementación del \mono{slam\_tracker} para Basalt. Si bien es solucionable,
como se ve en el gráfico, no afecta significativamente al rendimiento del
pipeline}, y luego de unos 12 ms representados por
\passthrough{\lstinline!PUSHED\_TO\_PROCESSED!}, a tiempo
\passthrough{\lstinline![F]!}, la pose estimada para el tiempo
\passthrough{\lstinline![A]!} está computada. Es decir, tenemos una
demora de unos 25,5 ms desde que la muestra es capturada hasta que el
sistema de SLAM/VIO es capaz de estimar la pose correspondiente al
momento de captura de la muestra. Cabe aclarar que estos tiempos son muy
variables incluso en la misma corrida, al depender de la calidad de la
conexión USB, el sistema utilizado, y las propias muestras que pueden
complicar las iteraciones de los algoritmos de optimización que ocurren
en el sistema.

Por otro lado, tenemos que a tiempo \passthrough{\lstinline![A]!}, la
aplicación OpenXR solicita a Monado una predicción de a dónde el runtime
piensa que el dispositivo se va a encontrar 7 ms en el futuro, en
\passthrough{\lstinline![B]!}. Cabe aclarar que la solicitud debe ser
respondida de forma inmediata y la barra
\passthrough{\lstinline!REQUEST\_TO\_PREDICTION!} no implica ninguna
espera hasta \passthrough{\lstinline![B]!}, es solo una forma de
visualizar los 7 ms. Notar que en ese punto, todavía faltan 25,5 ms para
tener la predicción de Basalt para \passthrough{\lstinline![A]!}, más
aún, una predicción dada por Basalt para el futuro
\passthrough{\lstinline![B]!} ni siquiera existirá, ya que el sistema
solo estima poses para los tiempos de las muestras, es decir la próxima
estimación correspondería al tiempo \passthrough{\lstinline![G]!}.
Además de esto, tenemos que mientras todavía se está procesando la
muestra, otra petición a tiempo \passthrough{\lstinline![D]!} para
\passthrough{\lstinline![E]!} debe ser respondida.

En conclusión, tenemos un \textbf{desfasaje temporal} que hace que las
poses estimadas siempre estén levemente en el pasado; en el tramo
seleccionado fue de 25,5 ms (más los 7 ms de predicción), pero es
variable durante la corrida. Además, las poses se estiman para
\textbf{puntos discretos} de tiempo, mientras que el usuario puede pedir
una predicción para cualquier punto arbitrario. Entonces, si queremos
ser capaces de proveer al usuario una pose para el tiempo solicitado,
necesitemos implementar formas de interpolar y \textbf{\emph{predecir}}
la ubicación de un espacio. Es decir, no utilizamos las poses devueltas
por el sistema de SLAM/VIO de forma directa, sino como parte fundamental
de un procedimiento constante de predicción que se realiza del lado de
Monado. En este se intenta utilizar todos los datos disponibles en el
runtime para proveer la pose más razonable en el punto de tiempo
requerido.

\hypertarget{implementaciuxf3n-1}{%
\subsubsection{Implementación}\label{implementaciuxf3n-1}}

Se encara la solución a estos problemas de manera progresiva y haciendo
uso de algunas de las herramientas e ideas ya presentes en Monado.

Dentro de Monado, el concepto de espacio de referencia o
\passthrough{\lstinline!XrSpace!} de de OpenXR, es nombrado como una
\emph{relación espacial} y se representa con un
\passthrough{\lstinline!struct!} muy similar al siguiente.

\begin{lstlisting}[language={C++}]
struct xrt_space_relation {
  struct vec3 position;
  struct quat orientation;
  struct vec3 linear_velocity;
  struct vec3 angular_velocity;
};
\end{lstlisting}

En él, no solo se guarda la información de la pose
(\passthrough{\lstinline!position!} y
\passthrough{\lstinline!orientation!}) sino que además tenemos el estado
de la velocidad lineal y angular de este espacio. Esto es muy útil, ya
que si sabemos que a tiempo \(t_1\) tenemos cierto espacio con pose
\(T_1 \in SE(3)\) y velocidades \(v_1, \omega_1 \in \R^3\), podemos
predecir que a un tiempo \(t_2 = t_1 + \Delta t\) tendremos el espacio
con pose \(T_2 = \Delta T \ T_1\) con

\begin{align}
\label{eq:predicted-space-delta}
\Delta T = \begin{bmatrix}
exp(\Delta t \ \hat\omega) & \Delta t \ v \\
0 & 1
\end{bmatrix} \in \R^{4x4}
\end{align}

Monado provee varias herramientas que facilitan tareas que suelen ser
recurrentes en diversos sistemas de tracking. La tarea de estimar
espacios futuros basándose en uno dado con sus velocidades es una de
estas funcionalidades ya incluidas. Más aún, Monado implementa una
estructura de datos que permite almacenar un ``historial'' de estos
espacios en una cola circular y generar las interpolaciones y
extrapolaciones, tanto a futuro como a pasado, necesarias para cualquier
timestamp requerida por un usuario. Para las interpolaciones se utiliza
una simple interpolación lineal\footnote{En el caso de la orientación es
  la interpolación esférica lineal \(slerp\) que definimos en la
  \Cref{eq:slerp-def}} o \emph{lerp} de a trozos entre cada par de
espacios del historial. Para extrapolar hacia el futuro, se usa la pose
y velocidad almacenada en el espacio más reciente del historial para
realizar el cómputo con \(\Delta T\) como se definió en la
\Cref{eq:predicted-space-delta}. Simétricamente para extrapolar hacia el
pasado lejano\marginnote{%\
La especificación de OpenXR tiene una sección dedicada a
las restricciones y condiciones a los que el runtime está sujeto respecto a
solicitudes en el pasado y en el futuro por parte del usuario. Ver
\autocite{thekhronosgroupinc.OpenXRSpecification}, secc. 2.14.
}, o sea fuera del registro del historial, se utilizará el espacio más
antiguo almacenado y \(\Delta T^{-1}\).

Sería razonable, como primera aproximación a nuestro problema, utilizar
este historial de espacios. Esto nos garantiza que podamos proveerle al
usuario una pose para cualquier punto arbitrario en el tiempo que
solicite, basándonos en las estimaciones del sistema de SLAM que
asumimos son precisas. Sin embargo, un leve problema que surge es que no
existe ningún requerimiento para los sistemas sobre la estimación de
velocidades, ya que no todos estiman esta variable. La interfaz
\passthrough{\lstinline!slam\_tracker!} solo garantizan la estimación de
la posición y orientación del dispositivo como puede verse en su
definición en el \Cref{lst:slam-tracker-def}. Lo que haremos para
solventar esto es computar las velocidades con base a los pares de poses
adyacentes que tengamos en el historial. Estas poses tienen su timestamp
correspondiente, entonces es sencillo computar la diferencia entre las
mismas respecto a la unidad de tiempo, dando como resultado una
estimación de la velocidad del espacio. La
\figref{fig:prediction-with-space-history} muestra como funcionaría este
tipo de predicción en un ejemplo simplificado en 2D y en el que asumimos
que las poses estimadas por el sistema de SLAM/VIO coinciden
perfectamente con la trayectoria real del dispositivo.

\fig{fig:prediction-with-space-history}{source/figures/prediction-with-space-history.pdf}{Predicción con historial de espacios}{%
Ejemplo de predicción con el historial de espacios. Se asume que las poses
estimadas por el sistema de SLAM son perfectas por simplicidad.
}

\FloatBarrier

Esto es una buena primera solución al problema, y es la opción más
básica que la clase adaptadora \passthrough{\lstinline!TrackerSlam!} le
ofrece a los usuarios de Monado. Desafortunadamente, si vemos el ejemplo
estudiado en la \figref{fig:prediction-timeline} notaremos que la
frecuencia de poses que se computan es muy baja en comparación a la
cantidad de veces que la aplicación OpenXR requiere una nueva prdicción.
En el ejemplo se tenían muestras (y por ende estimaciones) a 30 cuadros
por segundo mientras que se renderizaba a 60. En ese ejemplo se está
utilizando un monitor de computadora estándar como pantalla, pero en
cascos de realidad virtual sin embargo, las frecuencias de renderizado
alcanzan fácilmente los 90 o 120 cuadros por segundo, haciendo que la
cantidad de predicciones que hacemos entre estimación y estimación
crezca. Esto empeora la calidad de las predicciones significativamente,
generando movimientos imprecisos y ruidosos, y agravando los efectos de
\emph{motion sickness} o \emph{cinetosis}\footnote{\url{https://es.wikipedia.org/wiki/Cinetosis}}
que estas experiencias pueden producir.

Para mejorar nuestras predicciones, visto el problema de que las
estimaciones computadas se encuentran muy espaciadas, respecto a las
peticiones de predicción, vamos a utilizar las muestras de IMU. Estas
vienen usualmente a frecuencias mucho mayores que las de renderizado;
250 Hz en el caso del ejemplo estudiado en la
\figref{fig:prediction-timeline}. Además, a pesar de sufrir de severos
problemas de drift, al utilizarlas en ventanas cortas de tiempo (unos
pocos milisegundos), estos se ven reducidos en gran medida y la
odometría que sus sensores proveen resulta suficientemente precisa.

La clase adaptadora tiene acceso a estas muestras, ya que las intercepta
para redirigirlas hacia los sistemas de SLAM. Utilizaremos un concepto
similar al de la pre-integración de muestras de IMU explorado en la
\Cref{basalt-preintegration} con algunas ideas relacionadas con las
ecuaciones presentadas a partir de la \cref{eq:imu-preintegration}.
Consideremos que no vamos a querer interferir de ninguna manera con las
estimaciones generadas por los sistemas de SLAM, ya que no querremos
distorsionar dichas poses. El proceso de pre-integración correrá de
forma completamente aislada de los sistemas. Más aún, será un proceso
mayormente efímero que solo tendrá consecuencias sobre puntos de tiempos
para los cuales el sistema de SLAM todavía no tenga una estimación
posterior. En particular, nos limitaremos a refinar las estimaciones de
la velocidad lineal y angular, dejando que las herramientas de
predicción de Monado que se basan en el último espacio del historial,
computen la pose adecuada para la timestamp requerida. Acumularemos
promedios de las mediciones recientes para reducir el impacto del ruido
presente en las muestras de la IMU.

Tenemos entonces que el algoritmo utilizado finalmente es similar al
pseudo código que se presenta a continuación. Cabe aclarar que la
función \passthrough{\lstinline!predict\_pose(t)!} es llamada cuando el
usuario quiere una predicción a tiempo \passthrough{\lstinline!t!}.
Además las herramientas de Monado son representadas por
\passthrough{\lstinline!relation\_history!} (el historial de relaciones)
y \passthrough{\lstinline!predict\_from\_space!} (la función de
predicción en base a un espacio).

\clearpage

\begin{lstlisting}[language={C++}, caption={Predicción de poses en Monado}, label=lst:predict-pose]
struct xrt_space_relation predict_pose(timestamp t) {
   if (relation_history.is_empty()) return {0};

   // Espacio más reciente estimado con SLAM/VIO
   struct xrt_space_relation r = relation_history.get_latest();
   timestamp rt = timestamp_of(r);

   // Variables configuradas por el usuario en tiempo de ejecución
   bool pred_on = /* predicción habilitada por usuario? */
   bool gyr_on = /* giroscopio habilitado por usuario? */
   bool acc_on = /* acelerómetro habilitado por usuario? */
   bool imu_on = gyr_on or acc_on;

   // Flujo condicional de la predicción según la configuración
   if (not pred_on) return r;
   if (not imu_on or t <= rt) relation_history.predict(t);
   if (gyr_on) {
      vec3 avg_gyr = gyr_average_between(rt, t);
      vec3 world_gyr = rotate_angular_velocity(r.orientation, avg_gyr);
      r.angular_velocity = world_gyr;
   }
   if (acc_on) {
      vec3 avg_acc = acc_average_between(rt, t);
      vec3 world_acc = rotate_linear_acceration(r.orientation, avg_acc);
      world_acc += gravity_vector;
      double dt = last_imu_timestamp - rt;
      r.linear_velocity += world_acc * dt;
   }

   return predict_from_space(r, t);
}
\end{lstlisting}

Es interesante notar la naturaleza modular del algoritmo representada
por las condiciones de los \passthrough{\lstinline!if!}, en donde
distintos componentes que proveen información a la predicción pueden
deshabilitarse. Notar que la velocidad angular provista por el
giroscopio es local y necesita ser ajustada a coordenadas globales; esto
se realiza con la orientación estimada en \passthrough{\lstinline!r!},
la pose más reciente computada por el sistema de SLAM. Similarmente el
acelerómetro debe ser corregido, y a este además se le suma una
corrección con el vector de la gravedad\footnote{El vector de gravedad
  puede computarse dinámicamente con la IMU, detectando momentos en los
  que el dispositivo está quieto y registrando el vector medido por el
  acelerómetro.}. A diferencia del giroscopio, el acelerómetro solo nos
puede proveer información sobre los cambios de velocidad, y no sobre la
velocidad inicial; para esta usamos la dada por
\passthrough{\lstinline!r!} que es computada con la diferencia de las
dos poses en \passthrough{\lstinline!relation\_history!} más recientes.

Mostramos en la \figref{fig:prediction-with-imu} un ejemplo simplificado
del algoritmo implementado utilizando esta idea de promediar muestras de
odometría de la IMU para predecir puntos de tiempo posteriores a
\passthrough{\lstinline!B!} cuando \passthrough{\lstinline!C!} todavía
no pertenece al historial de espacios. En el ejemplo se muestra como
iría actualizándose el vector promediado (en verde) al integrar las tres
muestras (en azul) que ocurren entre \(t=1\) y \(t=2\). Además, se
muestra como estos vectores promedios se utilizan para extrapolar
linealmente para los tiempos requeridos por un usuario
\(t \in \{1,45; 1,75; 1,95\}\) (en anaranjado). Por simplicidad, el
ejemplo asume que tanto las muestras de odometría de la IMU a tiempos
\(t \in \{1,3; 1,6; 1,9\}\) como las estimaciones
\passthrough{\lstinline!B!} y \passthrough{\lstinline!C!} coinciden
perfectamente con la trayectoria real del dispositivo.

\fig{fig:prediction-with-imu}{source/figures/prediction-with-imu.pdf}{Predicción con promediado de muestras IMU}{%
Ejemplo de predicción para tiempos $t \in \{1,45; 1,75; 1,95\}$ utilizando la
idea de promediar muestras de la IMU posteriores a la pose más reciente del
historial (\mono{B} en este caso, se considera que \mono{C} no pertenece al historial
aún). Se asume que las poses estimadas por el sistema de SLAM y las muestras de
la IMU son perfectas por simplicidad.
}

Para contrastar esto con el caso en el que se ignoran las muestras de la
IMU y solo se utiliza el historial de poses, se muestra en la
\figref{fig:prediction-without-imu} los errores de esta predicción para
el mismo escenario.

\fig{fig:prediction-without-imu}{source/figures/prediction-without-imu.pdf}{Predicción sin uso de muestras IMU}{%
Ejemplo de predicción para tiempos $t \in \{1,45; 1,75; 1,95\}$ ignorando las
muestras de IMU y utilizando unicamente el historial de poses con \mono{A} y \mono{B} como
últimas muestras (esto es antes de que llegue \mono{C}). De vuelta, asumimos que las
estimaciones y no contienen error por simplicidad.
}

Finalmente, cabe aclarar que la trayectoria presentada en estas figuras
es particularmente desafiante. Considerando que estos dispositivos XR
están usualmente sujetos a partes como la cabeza o las manos del
usuario, es improbable encontrar trayectorias demasiado abruptas entre
los tiempos de estimación de las poses de SLAM. Estos, al coincidir con
la frecuencia de muestreo de las cámaras, suelen ser menores a 50ms
(p.ej. a 30 cuadros por segundo, tenemos \textasciitilde33ms entre pose
y pose). Por otro lado, es usual tener una mayor cantidad de muestras de
IMU entre estimaciones y no solo las 3 que se muestran en las figuras
(p.ej. a 30 cuadros por segundo con la IMU a 250 hz, tenemos unas
\textasciitilde8 muestras de IMU entre cada par de cuadros
consecutivos), lo cual mejora la precisión de la predicción aún más.

\FloatBarrier

\hypertarget{filtrado-de-poses}{%
\subsection{Filtrado de poses}\label{filtrado-de-poses}}

Continuando con la descripción de la funcionalidad presente en la clase
adaptadora \passthrough{\lstinline!TrackerSlam!} y luego de haber
presentado el método de predicción que se emplea en Monado, veremos
ahora la siguiente funcionalidad que
\passthrough{\lstinline!TrackerSlam!} implementa: el \textbf{filtrado}
de poses.

También llamado \emph{smoothing}, alisado o suavizado, el filtro de
señales, curvas o, en este caso poses, es una forma de minimizar o
\emph{filtrar} el ruido presente en alguna fuente de información y
suavizar los cambios abruptos en las poses. Lo que sucede en nuestro
caso es que por la naturaleza del problema en cuestión y por algunas de
las decisiones tomadas, se pueden notar la presencia de micro
movimientos de alta frecuencia que existen en la trayectoria estimada.
Estas alteraciones no son más que ruido proveniente de distintos
factores.

Por un lado, los sistemas de SLAM/VIO en el área no necesariamente
priorizan el caso de uso de XR y suelen estar mayormente interesados en
aplicaciones de robótica. En este campo, la ``sensación'' que produciría
el tracking no es un factor de importancia comparado a la precisión
absoluta del mismo. En particular, las métricas más utilizadas en estos
trabajos son las que evalúan el error \emph{absoluto} de la trayectoria
en su totalidad estimada sobre conjuntos de datos estándar. Es intuitivo
pensar que estas métricas, a diferencias de sus versiones
\emph{relativas}, incentivan correcciones abruptas en la trayectoria
cuando los sistemas determinan que su pose actual no es la mejor posible
(p.~ej. en momentos de loop closure). El trabajo de
\textcite{zhangTutorialQuantitativeTrajectory2018} es una buena
explicación de estas métricas y de la forma de evaluar y reportar el
rendimiento de estos sistemas.

De la misma forma, minimizar el ruido en la trayectoria no es uno de los
objetivos usualmente considerados, ya que la aplicación artificial de
filtros que suavicen la trayectoria podría tender a empeorar el puntaje
obtenido en estos conjuntos de datos. Todo esto hace que muchas veces
las trayectorias computadas por los sistemas en condiciones no óptimas
tiendan a presentar una gran cantidad de movimientos bruscos. Estos
pueden resultar particularmente notables en sistemas VR en dónde las
pantallas ocupan todo el campo de visión del usuario; movimientos de
este tipo pueden fácilmente inducir cinetosis.

Por sobre esto, el método de predicción presentado en la sección
anterior en la \figref{fig:prediction-with-imu} presenta problemas que
no fueron tratados. En particular, la predicción que realizamos está
basada en una única pose del sistema SLAM y las muestras nueva de la
IMU, pero no comparte ningún otro estado en común. Entonces podría
esperarse que la diferencia entre las poses que se fueron prediciendo
con la nueva pose que el sistema estimará sea significativa. Más aún, en
la realidad el error acumulado por la IMU, a pesar de verse limitado a
un intervalo corto de unas decenas de milisegundos, es un término más
que afectaría esta diferencia. Todo esto puede verse en el error marcado
en la \figref{fig:prediction-offset-jump}. Este tipo de desfasaje sería
reintroducido de forma repetitiva afectando las predicciones cada vez
que una nueva pose es estimada, causando así constantes micro saltos
(ruido), especialmente en los momentos en los que la pose estimada por
SLAM y la pose predicha por nuestro método difieran significativamente.

\fig{fig:prediction-offset-jump}{source/figures/prediction-offset-jump.pdf}{Predicción sin uso de muestras IMU}{%
Desfasaje esperado entre el método de predicción por muestras promediadas y la
pose \mono{C} estimada de SLAM cuando \mono{C} aún no pertenece al historial de espacios.
Se muestra la incidencia de la acumulación de errores de la IMU. Este desfasaje
causará micro saltos que serían disruptivos en una experiencia de VR.
}

Es entonces por estas razones, que se implementaron tres métodos
sencillos de filtrado combinables para suavizar las trayectorias
estimadas y que intentan minimizar los problemas expuestos. Se deja como
trabajo futuro el uso de filtros más elaborados como los filtros Kalman
\autocite[ap. B]{welchSCAATIncrementalTracking1997}. Cabe aclarar que en
general, el uso de filtros también presenta una desventaja fundamental.
Esta es, la introducción de latencia artificial al sistema, ya que para
poder suavizar este ruido, los filtros tienden a reducir cambios
abruptos en la trayectoria, incluso cuando estos sean realmente los
movimientos que el usuario realizó físicamente.

Presentaremos tres filtros, con el más sofisticado basado en
\textcite{casiezFilterSimpleSpeedbased2012}. Además, dicho trabajo cubre
el resto de filtros presentados, muestra gráficas comparativas, y es
también una buena lectura para contextualizar el uso de filtros para
tracking de movimientos humanos. Los filtros presentados aquí y en el
trabajo mencionado pueden visualizarse interactivamente en la aplicación
web referenciada al pie de página\footnote{\url{https://cristal.univ-lille.fr/~casiez/1euro/InteractiveDemo/}.}.

El funcionamiento esquemático de un filtro es muy sencillo: es un
contenedor de \textbf{estado}, con un método de \textbf{actualización}
que recibe un nuevo \textbf{dato} con una \textbf{timestamp} posterior a
las provistas hasta el momento; con esto, el filtro computa el valor de
la \textbf{señal filtrada} para una timestamp requerida. En nuestro
caso, esta timestamp coincidirá con la del dato de entrada. En Monado,
ubicaremos el filtro al final del pipeline de
\passthrough{\lstinline!TrackerSlam!}, justo antes de devolverle la pose
a la aplicación OpenXR, es decir luego de haber pasado por el
procedimiento de predicción. Los filtros habilitados interceptan la pose
con sus métodos de actualización respectivos y devuelven en su lugar una
nueva pose filtrada con la misma timestamp. Los tres filtros se
encuentran uno detrás del otro y son combinables.

Expresado más formalmente, lo que tendremos es un conjunto de poses
\(X_0, ..., X_k \in \R^7\) que ocurren a tiempos \(t_0 < ... < t_k\)
respectivamente. Para \(i = 0, ..., k\) definimos \(X_i = (p_i, q_i)\)
con \(p_i \in \R^3\) definido por los primeros tres componentes de
\(X_i\) para representar la posición, mientras que definimos el
cuaternión \(q_i\) para representar la orientación con coeficientes
dados por los últimos cuatro componentes. Tenemos además las versiones
filtradas de las poses anteriores \(\hat{X}_0, ..., \hat{X}_{k - 1}\) y
querremos ahora computar la versión filtrada de la pose actual
\(\hat{X}_k\).

\hypertarget{media-muxf3vil}{%
\subsubsection{Media móvil}\label{media-muxf3vil}}

El primer filtro es una \emph{media móvil}. El estado de este filtro
consiste de un historial de poses de los últimos \(m\) segundos. El
método de actualización registra la pose de entrada y su timestamp en
este historial interno. Para computar la pose de salida en esa
timestamp, se toma el promedio de las poses de los últimos \(w < m\)
segundos (por defecto 66 ms). El parámetro \(w\) es configurable por el
usuario.

Cabe aclarar, que si bien la forma de calcular el promedio para las
posiciones \(p_i\) debería resultar clara, promediar las orientaciones
\(q_i\) expresadas mediante cuaterniones no es trivial. En este caso,
nos aprovecharemos del hecho de que \(w\) suele ser pequeño y por ende
contiene orientaciones que no cambian significativamente. Esto nos
permite utilizar el resultado expuesto por
\textcite{gramkowAveragingRotations2001} que muestra que, para
rotaciones de menos de 40 grados, calcular la media usual (y
posteriormente normalizarla) es una muy buena aproximación con un error
de menos del \(1\%\).

Tenemos entonces que el filtro queda definido por la siguiente ecuación.
\begin{align}
\hat{X}_k &= \frac{1}{|W|} \sum_{i \in W}{X_i}
\quad \text{con} \\
W &= \{ i : t_k - w \leq t_i \leq t_k \} \quad \text{(ventana del filtro)}
\end{align}

Notar que podemos definir el concepto de sumatoria componente a
componente en \(\R^7\) gracias a la aproximación de los cuaterniones
mencionada anteriormente.

\hypertarget{suavizado-exponencial}{%
\subsubsection{Suavizado exponencial}\label{suavizado-exponencial}}

Este filtro codifica en un único valor los datos históricos e integra
nuevos datos con una intensidad dada por un \emph{factor de suavizado}
\(\alpha \in [0, 1]\) configurable por el usuario (0,1 por defecto). En
el caso de tener un único escalar \(x_k\) que filtrar, el suavizado
exponencial queda definido de esta forma: \begin{align}
\hat{x}_0 &= x_0 \\
\hat{x}_k &= \alpha x_k + (1 - \alpha) \hat{x}_{k-1} \label{eq:exp-smoothing-scalar}
\end{align}

Reformulemos la \Cref{eq:exp-smoothing-scalar} de la siguiente manera:
\begin{align}
\hat{x}_k &= \alpha x_k + (1 - \alpha) \hat{x}_{k-1} \\
&= \alpha x_k + \hat{x}_{k-1} - \alpha \hat{x}_{k-1} \\
&= \hat{x}_{k-1} + \alpha (x_k - \hat{x}_{k-1})
\end{align}

Esto nos deja ver el paso de actualización del filtro como una
interpolación (lineal) de \(\hat{x}_{k-1}\) hacia \(x_k\) con paso
\(\alpha\). Utilizaremos esta idea para interpolar esféricamente la
orientación de \(\hat{X}_k\). Tenemos entonces que el paso de
actualización para este filtro queda definido como: \begin{align}
\hat{X}_k &= (\hat{p}_k, \hat{q}_k) \\
\hat{p}_k &= lerp(\hat{p}_{k-1}, p_k, \alpha) = \hat{p}_{k-1} + \alpha (p_k - \hat{p}_{k-1}) \\
\hat{q}_k &= slerp(\hat{q}_{k-1}, q_k, \alpha) = \hat{q}_{k-1}
(\hat{q}_{k-1}^{-1} q_k)^\alpha
\end{align}

\hypertarget{filtro-1}{%
\subsubsection{Filtro 1€}\label{filtro-1}}

El filtro 1€ \autocite{casiezFilterSimpleSpeedbased2012} se basa en el
\hyperref[suavizado-exponencial]{suavizado exponencial}, pero utiliza un
factor \(\alpha\) dinámico que se adapta automáticamente con base a la
tasa de cambio de la señal. Reusamos también la idea de interpolar
esféricamente para la orientación presentada en el filtro anterior. El
filtro queda definido para la posición \(p_k\) de la siguiente manera:
\begin{align}
\hat{p}_0 &= p_0 \\
\hat{p}_k &= lerp(\hat{p}_{k-1}, p_k, \alpha)
\end{align}

Con \(\alpha\) que se adapta con la velocidad de la señal: \begin{align}
\alpha &= \frac{1}{1 + \frac{\tau}{\Delta t_k}} \\
\Delta t_k &= t_k - t_{k-1} \\
\tau &= \frac{1}{2 \pi f_C}
\end{align}

A continuación, \(f_C\) es la llamada \emph{frecuencia de corte}
\marginnote{%\
Algunos lectores reconocerán que el término “frecuencia de corte”
proviene de los filtros low-pass y efectivamente, notarán que el filtro 1€ es
de este tipo con la particularidad de tener una frecuencia de corte dinámica.
} y posee un mínimo ajustable por el usuario \(f_{C_{min}}\) y un
parámetro de intensidad de actualización \(\beta\) también configurable
\marginnote{%\
Frecuencias de corte bajas reducen el ruido de las poses a
costa de aumentar la latencia \autocite{stauffertKaBoomVisuallyExploring2021}. La forma en la que $f_c$ es definida resulta en
valores altos (y por ende con baja latencia) cuando se presentan cambios
significativos en las poses, mientras que para movimientos más suaves se reduce
$f_c$ para a su vez disminuir el ruido.
}. \begin{align}
f_C &= f_{C_{min}} + \beta | \hat{\dot{p}}_k |
\end{align}

La velocidad de la señal es a su vez ajustada con otro filtro de
suavizado exponencial con factor de suavizado fijo \(f_{C_d}\), también
configurable por el usuario. \begin{align}
\hat{\dot{p}}_0 &= 0 \\
\hat{\dot{p}}_k &= lerp(\hat{\dot{p}}_{k-1}, \dot{p}_k, f_{C_d})
\end{align}

Con velocidades instantáneas. \begin{align}
\dot{p}_0 &= 0 \\
\dot{p}_k &= \frac{p_k - \hat{p}_{k-1}}{\Delta t_k}
\end{align}

La versión del filtro utilizada para la orientación \(q_k\) es análoga a
la definición anterior con algunas aclaraciones:

\bigbreak

\begin{itemize}
\item
  En lugar de \(lerp\) se utiliza \(slerp\).
\item
  Sobrecargamos el operador de resta de cuaterniones de la siguiente
  forma: \(q_a - q_b = q_b^{-1} q_a\).
\item
  La norma de un cuaternión es equivalente a la norma euclídea en
  \(\R^4\), es decir \(|q| = \sqrt{q_x^2 + q_y^2 + q_z^2 + q_w^2}\).
\end{itemize}

\bigbreak

\hypertarget{recapitulando}{%
\section{Recapitulando}\label{recapitulando}}

Con la \figref{fig:slam-tracker-dataflow} vista al principio de la
sección en mente, hemos visto que se ha implementado una clase
adaptadora \passthrough{\lstinline!TrackerSlam!} en Monado que funciona
como punto central de comunicación con los sistemas de SLAM/VIO
integrados. Esta clase además implementa funcionalidades de predicción y
filtrado de poses que son configurables en runtime.

Los controladores de dispositivos físicos serán los encargados de
instanciar y comunicarse con esta clase. Veremos en la siguiente sección
los dos controladores (y por ende dispositivos) que necesitaron ser
expandidos para este trabajo y que constituyen una contribución
importante al runtime.

\hypertarget{drivers}{%
\chapter{Controladores en Monado}\label{drivers}}

En Monado, la interacción con la gran variedad de dispositivos que el
runtime soporta es realizada mediante \emph{drivers} o
\emph{controladores}. Estos, como se mostró en la
\figref{fig:slam-tracker-dataflow}, le permiten a Monado interactuar con
sistemas XR físicos mediante abstracciones derivadas de los requisitos
de OpenXR. Un sistema XR en este contexto hace referencia a un conjunto
de dispositivos XR. Un dispositivo XR, en forma intuitiva, es algún tipo
de hardware que permite la entrada y/o salida de información para
intercambio con aplicaciones de XR. Un caso paradigmático de un sistema
XR podría considerarse al conjunto de casco y un par de mandos provistos
por un fabricante.

El concepto que se termina implementando en Monado es un poco más
general, ya que además de sistemas físicos, soporta sistemas simulados
que proveen distintas funcionalidades. Entre estas se incluyen la
capacidad de conectar dispositivos de forma remota, emulación de
dispositivos con teclado y ratón\footnote{\url{https://gitlab.freedesktop.org/monado/monado/-/merge_requests/714}}
o mediante otros dispositivos como placas Arduino\footnote{\url{https://www.arduino.cc/}}.
En este trabajo se diferenció el concepto de \emph{fuente de datos} del
de \emph{dispositivo} ya que es en definitiva esto en lo que estaremos
interesados para SLAM, obtener fuentes de datos de IMU y cámaras.

\fig{fig:devices-ody-d455}{source/figures/devices-ody-d455.jpg}{Dispositivos XR utilizados}{%
Dispositivos XR utilizados en este trabajo. A izquierda un casco Samsung
Odyssey+ y a derecha una cámara Intel RealSense D455.
}

Se utilizaron los dos dispositivos que se muestran en la
\figref{fig:devices-ody-d455} como principales fuentes de datos para
SLAM. A la derecha de la imagen tenemos una cámara de profundidad
\emph{Intel RealSense D455}\footnote{\url{https://www.intelrealsense.com/depth-camera-d455}}
mientras que a izquierda tenemos un casco \emph{Samsung Odyssey+}. La
línea de cámaras y módulos RealSense de Intel se enfoca en aplicaciones
de robótica y visión por computadora, presentan distintos modelos con
múltiples sensores especializados; en nuestro caso nos limitaremos a
utilizar su IMU y cámaras estéreo. Estos vienen precalibrados, y además
se tiene un SDK\footnote{\url{https://github.com/IntelRealSense/librealsense}}
de código abierto en C/C++, con \emph{bindings} para otros lenguajes,
que facilita la obtención y manipulación de datos. En contraste con
esto, el casco de Samsung es un casco ligado a la plataforma privativa
\emph{Windows Mixed Reality (WMR)}\footnote{\url{https://www.microsoft.com/en-us/mixed-reality/windows-mixed-reality}}
que solo es soportada en sistemas operativos Windows\footnote{\url{https://www.microsoft.com/en-us/windows/}}.
WMR incluye algoritmos propietarios de tracking por SLAM desarrollados
por Microsoft.

Mientras que la cámara D455 funcionó como un dispositivo sumamente
versátil para el prototipado y experimentación con sistemas de SLAM, el
Odyssey+ presenta serios desafíos que requirieron trabajo con la
comunidad y métodos de ingeniería inversa para poder obtener acceso a
las fuentes de datos necesarias para el tracking por SLAM. Cabe aclarar,
que anterior a este trabajo, y según mi mejor entendimiento, no existía
forma de utilizar este tipo de cascos con tracking basado en SLAM/VIO
para correr aplicaciones OpenXR sobre sistemas operativos basados en
GNU/Linux \marginnote{%\
A pesar de que aún queda mucho por hacer para que el tracking
presentado llegue a niveles de calidad comparables a la versión
privativa de WMR, la contribución presentada deja asentada en upstream una
infraestructura sobre la que extender y mejorar el ecosistema de VR para
GNU/Linux.
} \marginnote{%\
Cuando hablamos de GNU/Linux nos referimos a distribuciones
enfocadas a computadoras personales como Ubuntu o Manjaro.
Técnicamente, los dispositivos autónomos Oculus Quest de Meta (y otros), corren
sobre sistemas operativos basados en Android, que a su vez está basado en
GNU/Linux.
}.

\hypertarget{sec:data-characteristics}{%
\section{Características de los datos}\label{sec:data-characteristics}}

Lo primero que se necesita para poder utilizar estos dispositivos para
SLAM es conseguir el acceso a los datos que generan; o sea los flujos de
imágenes y muestras de IMU. La forma y protocolos necesarios para
comunicarse con estos dispositivos se realiza de maneras específicas
para cada uno y como veremos en las secciones dedicadas, este es uno de
los trabajos fundamentales de un controlador. Además, un controlador
tiene que ser capaz, no solo de obtener esta información sino que
también de redirigirla a los módulos adecuados de Monado. En la
implementación de estas ideas surgen naturalmente los conceptos de
\emph{fuentes (sources)} y \emph{sumideros (sinks)} de datos. En el
runtime, estos conceptos existían exclusivamente para el tratado de
imágenes, y fueron extendidos para también soportar el flujo de muestras
de IMU. Los dispositivos entonces funcionan como fuentes de datos, que
instancian un \passthrough{\lstinline!TrackerSlam!} el cual les provee
sumideros a los que redirigir las muestras de sus sensores.

A la hora de implementar un controlador para una familia de dispositivos
enfocado a SLAM, hay una serie de cuestiones a tener en cuenta. Desde el
punto de vista de sistemas de SLAM/VIO, lo único que nos interesa es un
flujo adecuado de imágenes de cámaras y muestras de IMU. Es en la
definición de \emph{adecuado} en donde se esconden varios detalles
importantes. Intentaremos clarificar los las características que deben
considerarse en las muestras para que los sistemas de SLAM puedan
utilizarlas apropiadamente:

\hypertarget{calibraciuxf3n-de-paruxe1metros-intruxednsecos}{%
\subsection{Calibración de parámetros
intrínsecos}\label{calibraciuxf3n-de-paruxe1metros-intruxednsecos}}

En primer lugar, es necesario acceder a la información de calibración de
los sensores en cuestión, para poder comunicársela de alguna forma a los
sistemas de SLAM. Esto incluye la calibración para los cuatro sensores
usuales: giroscopio, acelerómetro y el par de cámaras estéreo. La
información de calibración describe valores que toman los parámetros
\emph{intrínsecos} (propios del sensor) de un modelo de calibración
especificado. Es usual que, para evitarle el proceso de calibración al
usuario de los sensores, estos provean valores de fábrica. La precisión
de tales valores ha probado ser de suficiente calidad para los sistemas
de SLAM integrados en este trabajo, pero es necesario aclarar que los
valores reales irán cambiando con el tiempo y el desgaste, haciendo que
la recalibración sea un punto importante en el uso de estos
dispositivos. Monado provee algunas herramientas básicas de calibración
para unos pocos modelos de cámara, pero será necesario profundizar en
este aspecto como trabajo a futuro.

Habiendo dicho esto, vale aclarar que existen sistemas como OpenVINS
\autocite{genevaOpenVINSResearchPlatform2020} que son capaces de
realizar el proceso de calibración de forma \emph{online} durante
corridas normales. Esto es una característica realmente importante que
ninguno de los tres sistemas integrados provee. Todos ellos asumen
parámetros de calibración fijos que deben ser provistos previos a la
ejecución.

Como intentamos evitarle la calibración manual al usuario, querremos ser
capaces de utilizar los valores de fábrica. Para esto es vital que los
sistemas de SLAM soporten los mismos modelos de calibración para los que
el fabricante especificó los valores\marginnote{%\
Mientras que en la práctica es conveniente recalibrar con modelos soportados,
podría resultar razonable estudiar expresiones analíticas que ``traduzcan'',
en alguna forma significativa, parámetros de un modelo a otro.
}. De lo contrario se necesitarán correcciones que usualmente terminan
siendo mucho más costosas que si los sistemas soportaran los modelos
nativamente. Ejemplos de esto son la \emph{desdistorsión}\footnote{Ver
  \passthrough{\lstinline!initUndistortRectifyMap!} en
  \url{https://docs.opencv.org/3.4/da/d54/group__imgproc__transform.html}}
y \emph{rectificación}\footnote{Ver
  \passthrough{\lstinline!stereoRectify!} en
  \url{https://docs.opencv.org/3.4/d9/d0c/group__calib3d.html}} de
imágenes sobre la cual no nos explayaremos aquí, pero basta con entender
que es una transformación costosa aplicada sobre todos los píxeles de
las imágenes para ``normalizarlas''. Esta normalización facilita su uso
cuando los sistemas no implementan los modelos de calibración en los que
el fabricante comparte los parámetros del dispositivo.

\hypertarget{calibraciuxf3n-de-paruxe1metros-extruxednsecos}{%
\subsection{Calibración de parámetros
extrínsecos}\label{calibraciuxf3n-de-paruxe1metros-extruxednsecos}}

El conjunto de sensores a utilizar: un par de cámaras estéreo y una IMU
con acelerómetro y giroscopio, deben estar sujetos a un cuerpo rígido.
Es decir, las transformaciones en \(SE3\) que describen como alterar la
pose de un sensor a otro se mantienen constantes sin importar los
movimientos del dispositivo en conjunto. Los valores que describen estas
transformaciones relativas son lo que se denominan \emph{parámetros
extrínsecos} de la calibración en contraste con los íntrinsecos. Los
tres sistemas integrados requieren estos valores antes de comenzar la
corrida.

\hypertarget{sincronizaciuxf3n-temporal}{%
\subsection{Sincronización temporal}\label{sincronizaciuxf3n-temporal}}

Otro problema escencial que debe solucionarse es el de la sincronización
de timestamps internas. Cualquier dispositivo que quiera utilizarse para
SLAM, debería tener las muestras de todos sus sensores sincronizadas por
hardware. Es decir, con timestamps que reflejen lo más precisamente
posible los tiempos internos en los que las muestras se tomaron.

Siguiendo con esta línea de problemas de marcas de tiempo, necesitamos
un mecanismo de sincronización de timestamps con el host que traiga los
relojes del dispositivo y los del host, el cual está corriendo Monado,
al mismo \emph{dominio temporal}. Es decir, que una timestamp en uno,
represente el mismo momento de tiempo en el otro. Como las aplicaciones
OpenXR utilizan el tiempo del host al solicitar predicciones de poses,
siempre intentaremos trabajar en ese dominio. Veremos que esta
sincronización no es trivial y es una fuente de errores común.

\hypertarget{muestras-de-imu-unificadas}{%
\subsection{Muestras de IMU
unificadas}\label{muestras-de-imu-unificadas}}

Entrando a problemas que afectan a sensores particulares, todos los
sistemas de SLAM estudiados esperan muestras de IMU unificadas que
combinen en una única timestamp los valores del acelerómetro y los del
giroscopio. En el caso de tener sensores con frecuencias diferentes o
timestamps que no coinciden (pero siempre sincronizadas), será necesario
realizar algun tipo de transformacion sobre las muestras para
unifcarlas.

\hypertarget{obturador}{%
\subsection{Obturador}\label{obturador}}

Por el lado de los sensores de la cámara, será importante que el
obturador o \emph{shutter}\marginnote{%\
El término obturador proviene de los sensores ópticos tradicionales.
Un obturador era una pieza mecánica en movimiento que controlaba el tiempo
durante el cual la película fotográfica era expuesta a la luz de la escena.
Actualmente el proceso de control de exposición se realiza con interrupciones
digitales.
} de las cámaras sea un \emph{global shutter}, es decir que todos los
píxeles sean capturados en el mismo instante. Por el contrario, las
cámaras que se utilizan habitualmente en dispositivos de consumo como
teléfonos inteligentes, presentan un \emph{rolling shutter} en dónde los
píxeles son capturados en un ``barrido'', fila por fila. Esto genera
distorsiones\marginnote{%\
Esta distorsión es homónima al tipo de cámara y se
denomina el efecto de rolling shutter.
} significativas en presencia de movimientos suficientemente rápidos
como se ve en la \figref{fig:rolling-shutter}. Estas distorsiones
afectan negativamente la capacidad de los algoritmos de SLAM para
reconocer y trackear features de forma estable.

\figw{fig:rolling-shutter}{source/figures/rolling-shutter.jpg}{Efecto rolling shutter}{%
Ejemplo de la distorsión que se genera en cámaras con rolling shutter que
capturan la imagen barriendo por filas de píxeles. Esto genera distorsiones en
los objetos rígidos que son particularmente notables en presencia de movimientos
de alta velocidad como los de la hélice. Esto no sucede con cámaras con global
shutter que capturan todos los píxeles en el mismo momento.\\
\\
\tiny Recorte de fotografía por Soren Ragsdale CC BY 2.0 \url{https://creativecommons.org/licenses/by/2.0/}
}{0.5\linewidth}

\hypertarget{exposiciuxf3n-y-ganancia}{%
\subsection{Exposición y ganancia}\label{exposiciuxf3n-y-ganancia}}

Otro aspecto muy importante a la hora de presentar muestras de imágenes
a sistemas de SLAM es el de utilizar valores adecuados de
\emph{exposición (exposure)} y \emph{ganancia (gain)}. La exposición o
exposure es la cantidad de tiempo que el obturador de la cámara habilita
la entrada de luz a los sensores ópticos por cada cuadro. Por otro lado,
la ganancia controla el nivel de amplificación, usualmente digital, que
ocurrirá sobre la señal original. El control de estos parámetros, y el
de la iluminación del entorno, es vital para asegurar imágenes que
posean un brillo adecuado. Imágenes muy oscuras pierden detalles, y por
ende la posibilidad de generar features. Imágenes \emph{sobreexpuestas}
que tienen demasiada iluminación en el entorno y valores de exposición y
ganancia altos, presentan el mismo problema al saturar los receptores
ópticos, causando que porciones significativas de la imagen se
transformen en manchas blancas. Estos problemas pueden verse en la
\figref{fig:under-over-exposure}. Además, los valores de exposición y
ganancia afectan el \emph{ruido} y el \emph{motion blur}, esto es la
difuminación que se genera por movimientos rápidos en la
imagen\marginnote{%\
El motion blur ocurre cuando los movimientos que se realizan modifican
sustancialmente la imagen a la que los sensores ópticos están expuestos mientras
el obturador sigue abierto.
}; ejemplos esto se muestran en la \figref{fig:expgain-grids}.

\fig{fig:under-over-exposure}{source/figures/under-over-exposure.pdf}{Poca y sobre-exposición}{
Arriba: imagen con valores de exposición y ganancia adecuados a las condiciones
de iluminación de la toma. Izquierda: imagen oscura. Derecha: imagen sobre
expuesta. Los histogramas muestran la cantidad de píxeles que toman los valores
del eje X de 0 a 255. En las imágenes de abajo hay pérdida de información hacia
ambos extremos del histograma.
}

\fig{fig:expgain-grids}{source/figures/expgain-grids.pdf}{Ruido y motion blur}{%
Efectos de la exposición y ganancia sobre el ruido y motion blur de la imagen.
El ejemplo fue tomados con la cámara D455 con distintas configuracions de
exposición y ganancia. A izquierda se observan múltiples capturas de la misma
región de 256 píxeles cuadrados que tiene un objeto en movimiento. A derecha
tenemos una porción más pequeña de 32 píxeles cuadrados provenientes de las
mismas imágenes. Se puede observar en ambas grillas que aumentar ambos valores
también incrementa el brillo de la imagen. Por su parte en la figura derecha
podemos ver que aumentar la ganancia incrementa el ruido, mientras que en la
figura izquierda se observa que valores altos de exposición incrementan el
motion blur del objeto en movimiento.
}

Un último problema a considerar que se presenta por el parámetro de
exposición, es el llamado \emph{parpadeo} o \emph{flicker}. Este ocurre
en entornos iluminados por lámparas artificiales. Estas presentan
oscilaciones de intensidad a altas frecuencias que no son visibles a
simple vista, pero que si esas frecuencias no están alineadas con las de
la exposición de las cámaras, producen una secuencia de imágenes con
niveles de brillo intermitentes. Este punto debe tenerse en cuenta, ya
que no todos los sistemas son capaces de trackear features de forma
eficiente cuando estas cambian las intensidades de sus píxeles.

Para evitar este abanico de problemas, se emplean algoritmos de
\emph{ajuste automático} de exposición y ganancia. Estos intentan
balancear los valores de estos parámetros en tiempo real; usualmente con
técnicas de análisis de histogramas. Hay que ser cuidadoso de que los
sistemas empleados soporten tales cambios en la naturaleza de los datos
introducidos. En este trabajo se evaluaron algunas posibilidades para
emplear algoritmos de este tipo pero ninguna se desarrolló lo suficiente
como para realizar una contribución a upstream. Implementar ideas como
las propuestas por \textcite{zhangActiveExposureControl2017} es una
tarea pendiente como trabajo a futuro. Por ahora, los controladores
recaen en ajustes manuales en los cuales se usa la heurística de que si
la imagen se ve razonablemente bien a los ojos del usuario, entonces
esta es probablemente adecuada para el sistema de SLAM.

\hypertarget{comunicaciuxf3n-con-el-dispositivo}{%
\subsection{Comunicación con el
dispositivo}\label{comunicaciuxf3n-con-el-dispositivo}}

La interfaz de comunicación con el hardware será diferente en cada
controlador, pero será usual tener que tratar con hilos consumidores y
callbacks asíncronos. En estos habrá colas de datos que deberemos evitar
saturar con procesamientos bloqueantes. Para el caso particular del
manejo de imágenes, al ser estos recursos de gran tamaño, Monado utiliza
mecanismos de \emph{reference counting} mediante su estructura
\passthrough{\lstinline!xrt\_frame!} para la gestión de memoria. Además,
como se mencionó anteriormente, \passthrough{\lstinline!slam\_tracker!}
utiliza la estructura de datos \passthrough{\lstinline!cv::Mat!} de
OpenCV como contenedora de imágenes la cual también provee conteo de
referencias. Finalmente, las interfaces con hardware mediante
bibliotecas específicas puede traer un contenedor extra de conteo de
referencias. Esto puede dar lugar a muchos problemas con el manejo de la
memoria y especial cuidado debe tenerse a la hora de adquirir y liberar
estos recursos.

\hypertarget{configuraciones-de-captura}{%
\subsection{Configuraciones de
captura}\label{configuraciones-de-captura}}

Los sensores suelen venir con distintos modos de captura de muestras.
Algunos habilitan mayores frecuencias a costa de menor precisión, o
mayor precisión a costa de un mayor consumo energético, y otras
soluciones de compromiso similares. La capacidad de acceso a la
configuración de estos sensores dependerá principalmente de lo que los
fabricantes del dispositivo decidan exponer al programador. En caso de
existir más de una forma de captura, será necesario poder seleccionar en
el controlador la configuración deseada.

\hypertarget{sistemas-de-coordenadas}{%
\subsection{Sistemas de coordenadas}\label{sistemas-de-coordenadas}}

Finalmente, el último punto que traeremos a atención, es la convivencia
de múltiples sistemas de coordenadas que deben unificarse: el de la IMU,
el de la implementación de SLAM y el de Monado.

Por un lado, será necesario en algunos casos alinear las muestras de la
IMU al sistema de coordenadas de la implementación de SLAM antes de
transferirlas. De lo contrario, se pueden presentar situaciones como una
implementación que considera que está moviéndose hacia adelante por las
muestras ópticas, mientras que las muestras de la IMU le indican que
está yendo hacia la derecha; esto causa inconsistencias que suelen
terminar en divergencias de los algoritmos de optimización. Además, es
usual que también se necesite aplicar un cambio de coordenadas a las
poses que el sistema le devuelve a Monado.

\hypertarget{realsense-todo}{%
\section{RealSense (TODO)}\label{realsense-todo}}

\begin{mdframed}[backgroundcolor=shadecolor]
TODO: Esta sección con esta cámara, y la siguiente con el casco VR deberían
simplemente terminar de concretar un poco los problemas que listé arriba con
algunos matices particulares que hayan aparecido. Escribí un poco de RealSense
pero no lo terminé.
\end{mdframed}

Estamos ahora en condiciones de entender las contribuciones a
controladores realizadas en este trabajo. Comencemos por las del
controlador para dispositivos RealSense.

Para soportar la cámara D455, se extendió significativamente en Monado
el controlador de dispositivos RealSense. Hasta el momento, la única
cámara de esta línea soportada por Monado era la T265\footnote{\url{https://www.intelrealsense.com/tracking-camera-t265/}}.
Esta cámara es curiosa, ya que presenta un algoritmo de SLAM privativo
que corre dentro del dispositivo sin necesidad de interactuar con el
host. Este controlador se encarga únicamente de inicializar el módulo
interno de SLAM de la cámara y obtener las poses computadas por la misma
para uso en las aplicaciones OpenXR. Uno de sus usuarios clave era el
casco libre del proyecto North Star \footnote{El proyecto North Star de
  UltraLeap (prev. LeapMotion) es un casco AR ``DIY'' que puede
  fabricarse con piezas impresas en 3D y la compra de algunos
  componentes. \url{https://developer.leapmotion.com/northstar}} que las
utilizaba, en iteraciones anteriores, como principal forma de tracking.
En la web pueden encontrarse imágenes \footnote{https://www.collabora.com/assets/images/blog/ProjectNorthStar.jpg}
que muestran estos cascos con la cámara T265 sujeta en su parte
superior.

Al no haber ningún tipo de manejo de las imágenes o muestras de IMU
provistas por la cámara, se tuvo que implementar la gestión de estos
sensores. Es ahora posible utilizar cualquier cámara de la línea
RealSense que posea un par de cámaras estéreo y una IMU. Como trabajo
futuro, sería interesante comparar el tracking interno de una T265 con
los distintos sistemas externos integrados en este trabajo.

La cámara D455 tiene un par de sensores con líneas de visión paralelas y
por ende las dos imágenes comparten la mayoría de sus respectivos campos
de visión. Las cámaras tienen una separación de unos 10cm entre ellas.
Presentan imágenes estéreo \emph{prerectificadas}, es decir que las
imágenes llegan al sistema de SLAM virtualmente sin ningún tipo de
distorsión, esto evita que se necesiten modelos de cámaras particulares
implementados en los sistemas externos. Poseen un \emph{obturador
global} y el sensor en sí es de \emph{buena calidad} siendo capaz de
producir imágenes con buena cantidad de información incluso en
situaciones con poca luz y con valores de exposición y ganancia bajos.
Cuenta con la capacidad de \emph{ajustar automáticamente} la ganancia y
exposición, aunque no pareciera que el algoritmo en uso beneficie a los
sistemas de SLAM y por ende se decidió desactivar esta característica
para en su lugar configurar estos valores manualmente desde la interfaz
gráfica de Monado.

Intel provee un SDK en C/C++ de código abierto y con una licensia
permisiva Apache 2.0 \autocite{ApacheLicenseVersion}. El SDK facilita la
interacción entre Monado y las cámaras; se encarga de la gestión de
colas y provee estructuras para manejo automático de la memoria de los
cuadros recibidos. Con el SDK es posible configurar múltiples formas de
ejecución de los sensores. Las cámaras soportan frecuencias de hasta 90
fps y resoluciones de hasta 1280x720 píxeles cada una, mientras que los
sensores que conforman la IMU se encuentra explicitamente divididos en
acelerómetro que puede ejecutarse a 60hz y 250hz y giroscopio que logra
alcanzar frecuencias de 200hz y 400hz.

Los \emph{relojes de la IMU y cámaras} están sincronizados por hardware,
pero es necesario aplicar un parche al kernel de la versión de Linux
instalada lo cual dificulta el uso. Sin este parche, se intenta utilizar
los tiempos de arribo al host de las muestras, pero es usual que estos
causen que los sistemas divergan por lo poco preciso que son. Una vez
que se tiene este parche, el SDK estima automáticamente la
sincronización entre los \emph{relojes del host y de la cámara} y por
ende es capaz de traducir timestamps de un dominio temporal al otro y
viceversa.

\hypertarget{windows-mixed-reality-todo}{%
\section{Windows Mixed Reality
(TODO)}\label{windows-mixed-reality-todo}}

\begin{mdframed}[backgroundcolor=shadecolor]
TODO: Similar a lo que menciné en la sección de arriba. Lo interesante de esta
sección creo yo es que participé un montón con la comunidad y le estuvimos
haciendo ingeniería inversa a estos cascos privativos para ver que comandos USB
mandarles y como leerlos para tener acceso a los streams de la cámara y de la
IMU. Además, esto es eso que marco en el abstract, es el primer casco comercial
que tiene tracking de este tipo andando en Linux con el stack full open source.
Y aunque está genial eso, es cierto que no le llega ni a los talones todavía lo
que tenemos comparado a la versión privativa de Microsoft. Pero bueno por lo
menos se puede caminar y dar vuelta un poco para observar cosas sin problemas lo
cual es mucho mejor que nada.
\end{mdframed}

\hypertarget{otras-contribuciones-todo}{%
\chapter{Otras contribuciones (TODO)}\label{otras-contribuciones-todo}}

\begin{mdframed}[backgroundcolor=shadecolor]
TODO: Hubo otras cositas más pequeñas contribuídas también, entre ellos unas
utilidades para grabar y reproducir datasets de SLAM (en formato EuRoC) muy utiles para testear el
sistema sin tener que andar moviendo las cámaras siempre. La otra contribución
fue a Basalt, el modelo de cámara que usan los cascos WMR y Basalt no lo tenía.
Además me gustaría listar todos los merge requests contribuídos en algún lado.
\end{mdframed}

\hypertarget{utilidades-euroc-todo}{%
\section{Utilidades EuRoC (TODO)}\label{utilidades-euroc-todo}}

\hypertarget{modelo-radial-tangencial-todo}{%
\section{Modelo radial-tangencial
(TODO)}\label{modelo-radial-tangencial-todo}}

\cleardoublepage
\ctparttext{Para cerrar presentaremos algunos resultados que
intentan describir el rendimiento y la precisión logrados en la implementación
actual. Concluiremos con una revisión general de los temas tratados y posibles
líneas de trabajo a considerar para el futuro.}

\hypertarget{conclusiones}{%
\part{Conclusiones}\label{conclusiones}}

\hypertarget{evaluation}{%
\chapter{Resultados}\label{evaluation}}

El proceso de evaluar y obtener métricas en sistemas de SLAM/VIO
requiere de ciertas consideraciones. A grandes rasgos, podemos dividir
las métricas de interés usuales en medidas de \emph{precisión} y de
\emph{eficiencia} que describen, respectivamente, la exactitud de la
trayectoria estimada y el uso de recursos por parte de los sistemas.
Para la evaluación de sistemas se desarrollaron funcionalidades
\footnote{\url{https://gitlab.freedesktop.org/monado/monado/-/merge_requests/1152}}
\footnote{\url{https://gitlab.freedesktop.org/monado/monado/-/merge_requests/1172}}
y herramientas \footnote{\url{https://gitlab.freedesktop.org/mateosss/xrtslam-metrics}}
dedicadas a la evaluación de sistemas de SLAM en Monado. Para más
información sobre evaluación que la que presentaremos en esta sección,
referimos al lector al trabajo de
\textcite{kummerleMeasuringAccuracySLAM2009} que detalla en mayor
profundidad el proceso de evaluación y a la suite de herramientas
SLAMBench\footnote{\url{https://apt.cs.manchester.ac.uk/projects/PAMELA/tools/SLAMBench}}
\autocite{nardiIntroducingSLAMBenchPerformance2015} que intenta
generalizarlo para una gran variedad de sistemas.

Para la evaluación se utilizan \emph{conjuntos de datos} o
\emph{datasets} pregrabados con distintos dispositivos. Utilizaremos dos
datasets populares en el área: \emph{EuRoC}
\autocite{burriEuRoCMicroAerial2016} que es grabado con un
\emph{vehículo micro aéreo} (\emph{MAV} o \emph{drone}) y \emph{TUM-VI}
\autocite{schubertTUMVIBenchmark2018} con muestras provenientes de un
dispositivo que es sostenido con la mano (\emph{handheld}) lo cual lo
hace particularmente bueno para evaluar tracking en aplicaciones de XR.
Además, estos conjuntos de datos fueron tomados con sistemas de captura
de movimiento (\emph{MoCap}) externos de gran precisión pero también de
gran costo. Esto les permite presentar una trayectoria muy precisa que
se utiliza como punto de referencia y se la conoce como \emph{ground
truth}.

Además de estos datasets, se presentan datos tomados especialmente para
este trabajo\footnote{\url{https://drive.google.com/drive/folders/163KuF88viW_wPcVNZJ2Onxe7zHf2Qo7L?usp=sharing}}
con los dispositivos introducidos en el \Cref{drivers}: la cámara
RealSense D455 y el casco Odyssey+. Se tienen también datos monoculares
del celular móvil Poco X3 Pro capturados con el trabajo de
\textcite{huaiMobileARSensor2019} pero no llegaron a utilizarse en estos
resultados; es decir, todos los datos evaluados son con cámaras estéreo.
Estos conjuntos capturados con la D455 y el Odyssey+ no presentan ground
truth, pero ofrecen grabaciones especialmente pensadas para XR y
utilizan dispositivos que Monado ya soporta. En la figura
\figref{fig:datasets-preview} se pueden observar algunas imágenes de los
distintos datasets utilizados en al evaluación.

\fig{fig:datasets-preview}{source/figures/datasets-preview.pdf}{Datasets}{%
Imágenes para visualizar el tipo de entorno en el que los datasets de evaluación fueron capturados.
}

Como es usual en este tipo de estudios comparativos, utilizaremos
acrónimos para referirnos a los distintos conjuntos de datos con las
siguientes características: \newline

\begin{itemize}
\item
  C*: Datasets específicos para este trabajo (\emph{custom}). Cada uno
  presenta dos modalidades, la primera EASY tiene movimientos tranquilos
  similares a los que se verían al inspeccionar una habitación. Por otro
  lado la modalidad HARD contiene una sucesión de movimientos agitados e
  intenta simular momentos de acción en un juego. Los sensores
  utilizados son: \newline

  \begin{itemize}
  \item
    C6*: RealSense D455 en 640x480 a 30 fps e IMU a 250 Hz.
  \item
    C8*: RealSense D455 en 848x480 a 60 fps e IMU a 250 Hz.
  \item
    CO*: Odyssey+ en 640x480 a 30 fps e IMU a 250 Hz. Para Basalt en
    particular tenemos dos posibles modelos de cámara para utilizar con
    los lentes de este casco. Por defecto se utilizará el modelo
    radial-tangencial de 8 parámetros \footnote{\url{https://gitlab.com/VladyslavUsenko/basalt-headers/-/merge_requests/21}}
    dado de fábrica pero también se comparará con el modelo
    Kannala-Brandt de 4 parámetros (KB4) recalibrado y nativo en Basalt.
  \end{itemize}
\item
  E*: Los datasets EuRoC en dos habitaciones (V1 y V2) y una sala de
  máquinas (MH) en 752x480 a 20 fps e IMU a 200 Hz.
\item
  T*: TUM-VI en una habitación (R) en 512x512 a 20 fps e IMU a 200 Hz.
\end{itemize}

También utilizaremos acrónimos para los distintos sistemas evaluados,
que son variantes de Basalt, Kimera y ORB-SLAM3 dados actualizaciones
significativas que le ocurrieron a Basalt\footnote{\url{https://gitlab.com/VladyslavUsenko/basalt/-/commit/24325f2a}}
y ORB-SLAM3\footnote{\url{https://github.com/UZ-SLAMLab/ORB_SLAM3/releases/tag/v1.0-release}}
durante el desarrollo de este trabajo.

\begin{itemize}
\tightlist
\item
  K: Kimera-VIO
\item
  OO: ORB-SLAM3 original antes de la versión 1.0.
\item
  ON: ORB-SLAM3 nueva versión 1.0.
\item
  BO: Basalt original.
\item
  BNF: Basalt nuevo luego de la actualización presentada en
  \textcite{demmelBasaltSquareRoot2021} la cual incluye, entre otras
  cosas, la posibilidad de especificar la precisión de punto flotante
  utilizada. Esta versión utiliza \passthrough{\lstinline!float!}.
\item
  BND: Idéntico al punto anterior pero con precisión
  \passthrough{\lstinline!double!}. \newline
\end{itemize}

Todas las corridas fueron realizadas sobre un procesador Intel i7-1065G7
con consumo de hasta 15 W y memoria suficiente; ninguno de los sistemas
hace uso de GPU.

\hypertarget{completitud}{%
\section{Completitud}\label{completitud}}

Nuestro primer análisis se encuentra en la \autoref{tab:completion} y se
centra en la capacidad de los sistemas integrados de no presentar fallas
fatales durante las corridas de los conjuntos de datos probados. Esta
métrica muestra la tolerancia a fallas de las distintas
implementaciones. La tabla presenta a Basalt como el más estable
mientras que ORB-SLAM3 presenta algunas fallas ocasionales. Kimera por
el otro lado presenta severas fallas incluso en conjuntos estándares
como TUM-VI. Vale la pena aclarar que en la práctica, si bien Basalt
resulta más estable que los otros sistemas, existen situaciones que
pueden hacerlo fallar, particularmente ante la presencia de muestras de
baja calidad durante tiempos prolongados. Es notable que la introducción
de los datasets C* es particularmente desafiante para la estabilidad de
los sistemas en comparación a los conjuntos estándares EuRoC y TUM-VI
para los cuales las implementaciones suelen estar bien probados y
configurados por defecto.

\begin{table}[h]
\caption[Completitud de ejecución]{Completitud de ejecución}
\label{tab:completion}
\resizebox{\textwidth}{!}{
\begin{tabular}{ |l||l|l|l|l|l|l|  }
\hline
  \spacedlowsmallcaps{Dataset} & \spacedlowsmallcaps{BND} & \spacedlowsmallcaps{BNF} & \spacedlowsmallcaps{BO}  & \spacedlowsmallcaps{K} & \spacedlowsmallcaps{ON} & \spacedlowsmallcaps{OO}\\[3pt]
 \hline
C6EASY & ✓     & ✓     & ✓     & ✓       & ✓       & ✓       \\
C6HARD & ✓     & ✓     & ✓     & 35.89\% & ✓       & ✓       \\
C8EASY & ✓     & ✓     & ✓     & ✓       & ✓       & ✓       \\
C8HARD & ✓     & ✓     & ✓     & 52.25\% & 56.61\% & 55.86\% \\
COEASY & ✓     & ✓     & ✓     & ✓       & ✓       & ✓       \\
(KB4)  & ✓     & ✓     & ✓     &         &         &         \\
COHARD & ✓     & ✓     & ✓     & ✓       & ✓       & 95.71\% \\
(KB4)  & ✓     & ✓     & ✓     &         &         &         \\
EMH01  & ✓     & ✓     & ✓     & ✓       & ✓       & ✓       \\
EMH02  & ✓     & ✓     & ✓     & ✓       & ✓       & ✓       \\
EMH03  & ✓     & ✓     & ✓     & ✓       & ✓       & ✓       \\
EMH04  & ✓     & ✓     & ✓     & ✓       & ✓       & ✓       \\
EMH05  & ✓     & ✓     & ✓     & ✓       & ✓       & 96.48\% \\
EV101  & ✓     & ✓     & ✓     & ✓       & ✓       & ✓       \\
EV102  & ✓     & ✓     & ✓     & ✓       & ✓       & ✓       \\
EV103  & ✓     & ✓     & ✓     & ✓       & ✓       & ✓       \\
EV201  & ✓     & ✓     & ✓     & ✓       & ✓       & ✓       \\
EV202  & ✓     & ✓     & ✓     & ✓       & ✓       & ✓       \\
TR1    & ✓     & ✓     & ✓     & 40.25\% & ✓       & ✓       \\
TR2    & ✓     & ✓     & ✓     & 38.32\% & ✓       & ✓       \\
TR3    & ✓     & ✓     & ✓     & ✓       & ✓       & ✓       \\
TR4    & ✓     & ✓     & ✓     & 63.58\% & ✓       & ✓       \\
TR5    & ✓     & ✓     & ✓     & 52.67\% & 74.81\% & ✓       \\
TR6    & ✓     & ✓     & ✓     & 52.37\% & ✓       & ✓       \\
\hline
\textbf{Media} & \textbf{100\%} & \textbf{100\%} & \textbf{100\%} & \textbf{83.42\%} & \textbf{96.88\%} & \textbf{97.64\%} \\
\hline
\end{tabular}
}
\end{table}

\hypertarget{tiempos}{%
\section{Tiempos}\label{tiempos}}

En este trabajo nos limitaremos a presentar el tiempo promedio que le
toma a cada sistema devolver la estimación de la pose. Se plantea como
trabajo a futuro la posibilidad de automatizar la medición del consumo
de otros recursos como memoria, energía y capacidad de cómputo. Recordar
que los conjuntos de datos probados presentan cuadros a 20, 30 y 60 fps
respectivamente y para lograr tracking a tiempo real necesitaríamos
tiempos de estimación menores a 50, 33 y 16 ms respectivamente.

Viendo la \autoref{tab:timing} vemos que Basalt (BNF) sale ganador con
tiempos bien por debajo de los 16 ms. La versión BND se presenta como
curiosidad y muestra lo ineficiente que es el pipeline cuando se
utilizan números \passthrough{\lstinline!double!}; como veremos en las
siguientes tablas el tracking de BNF y BND dan los mismos resultados de
precisión gracias a la técnica introducida en la actualización
\autocite{demmelBasaltSquareRoot2021}. ORB-SLAM3 y Kimera tienen tiempos
de ejecución significativamente mayores indicando que convendría
utilizar sensores a menor frecuencia para estos sistemas si se quieren
evitar congestiones. Es necesario aclarar que los tiempos de ORB-SLAM3
podrían ser mejorados si se evitara utilizar las configuraciones que
vienen por defecto. De la misma manera, la construcción del mapa en
tiempo real de ORB-SLAM3 es pesada y en la nueva versión (ON) soporta la
capacidad construirlo de antemano para luego sólo ejecutar los módulos
de localización y no de mapeo, reduciendo así considerablemente los
tiempos de cómputo. Probar estas configuraciones con ORB-SLAM3 se deja
como trabajo a futuro.

\begin{table}[h]
\caption[Tiempos de ejecución]{Tiempos de ejecución medio por cuadro [ms]}
\label{tab:timing}
\begin{addmargin*}[-0.2\textwidth]{-0.2\textwidth}
\resizebox{1.4\textwidth}{!}{
\begin{tabular}{ |l||l|l|l|l|l|l|  }
\hline
  \spacedlowsmallcaps{Dataset} & \spacedlowsmallcaps{BND} & \spacedlowsmallcaps{BNF} & \spacedlowsmallcaps{BO}  & \spacedlowsmallcaps{K} & \spacedlowsmallcaps{ON} & \spacedlowsmallcaps{OO}\\[3pt]
 \hline
C6EASY & 826.60 ± 441.30   & 5.84 ± 1.38   & 9.05 ± 2.40    & 46.00 ± 6.10 & 36.11 ± 7.67  & 35.09 ± 11.81 \\
C6HARD & 668.75 ± 555.59   & 5.58 ± 1.38   & 9.83 ± 3.01    & 47.45 ± 7.94 & 30.66 ± 9.64  & 32.92 ± 12.41 \\
C8EASY & 940.54 ± 485.31   & 6.93 ± 2.10   & 12.89 ± 13.63  & 49.23 ± 7.37 & 33.69 ± 10.50 & 33.24 ± 10.22 \\
C8HARD & 716.58 ± 579.89   & 6.20 ± 2.46   & 12.60 ± 8.38   & 46.28 ± 7.89 & 35.83 ± 11.75 & 37.43 ± 12.04 \\
COEASY & 873.74 ± 404.36   & 6.17 ± 1.12   & 10.96 ± 2.94   & 37.25 ± 4.94 & 35.40 ± 9.35  & 29.41 ± 9.69  \\
(KB4)  & 734.47 ± 357.71 K & 6.24 ± 1.02 K & 10.92 ± 2.98 K &              &               &               \\
COHARD & 617.47 ± 419.17   & 5.71 ± 1.16   & 12.90 ± 3.83   & 37.31 ± 5.20 & 21.52 ± 7.61  & 23.69 ± 7.98  \\
(KB4)  & 592.62 ± 410.55 K & 5.81 ± 1.02 K & 12.81 ± 3.81 K &              &               &               \\
EMH01  & 2123.22 ± 1118.05 & 10.63 ± 3.22  & 14.17 ± 3.15   & 53.15 ± 7.00 & 30.29 ± 6.63  & 36.73 ± 12.68 \\
EMH02  & 2280.38 ± 1118.44 & 11.16 ± 4.28  & 15.33 ± 5.76   & 53.93 ± 6.05 & 29.29 ± 5.37  & 35.32 ± 10.40 \\
EMH03  & 2184.08 ± 940.11  & 11.02 ± 2.81  & 15.17 ± 3.65   & 53.83 ± 6.05 & 32.09 ± 5.71  & 37.08 ± 12.91 \\
EMH04  & 2117.07 ± 916.61  & 11.82 ± 3.73  & 15.83 ± 3.41   & 53.12 ± 7.19 & 29.77 ± 6.94  & 32.67 ± 11.86 \\
EMH05  & 2187.28 ± 902.72  & 11.17 ± 2.15  & 15.47 ± 3.63   & 53.40 ± 7.07 & 29.04 ± 6.17  & 34.06 ± 15.61 \\
EV101  & 1687.89 ± 524.66  & 10.23 ± 1.76  & 13.62 ± 2.21   & 54.37 ± 6.18 & 30.26 ± 5.95  & 35.43 ± 13.93 \\
EV102  & 1322.72 ± 624.59  & 10.18 ± 2.09  & 15.35 ± 3.63   & 55.48 ± 5.79 & 29.74 ± 6.06  & 32.71 ± 13.47 \\
EV103  & 844.55 ± 609.03   & 11.65 ± 2.56  & 17.31 ± 4.37   & 56.54 ± 6.47 & 34.74 ± 11.43 & 31.13 ± 10.70 \\
EV201  & 1628.73 ± 718.45  & 10.08 ± 1.89  & 15.53 ± 3.04   & 55.00 ± 5.66 & 36.63 ± 11.87 & 32.51 ± 10.04 \\
EV202  & 1296.74 ± 667.75  & 10.65 ± 3.49  & 17.57 ± 4.14   & 55.37 ± 5.24 & 37.77 ± 10.90 & 34.78 ± 11.45 \\
TR1    & 800.61 ± 327.45   & 6.37 ± 1.02   & 12.54 ± 2.70   & 21.34 ± 3.51 & 46.72 ± 11.81 & 44.95 ± 12.33 \\
TR2    & 767.92 ± 287.46   & 6.08 ± 0.93   & 11.47 ± 2.37   & 21.42 ± 2.53 & 44.78 ± 12.06 & 43.94 ± 13.21 \\
TR3    & 697.36 ± 285.12   & 5.96 ± 0.93   & 11.93 ± 2.47   & 23.31 ± 3.69 & 38.33 ± 9.31  & 41.27 ± 11.67 \\
TR4    & 857.84 ± 330.19   & 6.57 ± 1.19   & 11.74 ± 2.38   & 22.62 ± 6.31 & 39.54 ± 10.50 & 42.37 ± 11.85 \\
TR5    & 694.90 ± 308.46   & 6.09 ± 1.01   & 12.53 ± 2.98   & 20.44 ± 3.01 & 32.79 ± 5.37  & 42.42 ± 11.87 \\
TR6    & 1007.87 ± 269.40  & 7.00 ± 1.05   & 10.72 ± 1.80   & 22.33 ± 5.05 & 33.19 ± 6.98  & 44.83 ± 12.68 \\
\hline
\textbf{Media} & \textbf{1186.25 ± 566.77} & \textbf{8.13 ± 1.91} & \textbf{13.26 ± 3.86} & \textbf{42.69 ± 5.74} & \textbf{34.01 ± 8.62} & \textbf{36.09 ± 11.86}\\
\hline
\end{tabular}
}
\end{addmargin*}
\end{table}

\hypertarget{precisiuxf3n-de-la-trayectoria}{%
\section{Precisión de la
trayectoria}\label{precisiuxf3n-de-la-trayectoria}}

Para medir la precisión de la trayectoria completa estimada por el
sistema se la compara con las trayectorias ground truth provistas por
los datasets. La métrica usual para representar esta precisión es la
\emph{raíz del error cuadrático medio (RMSE)} del \emph{error de
trayectoria absoluto (ATE)} que se define de la siguiente manera
\autocite{sturmBenchmarkEvaluationRGBD2012}.

Una trayectoria será una secuencia de poses \(P_i \in SE(3)\) con \(i\)
siendo el tiempo o timestamp en el que la pose ocurrió. Tenemos dos
trayectorias que considerar; la estimada que denotaremos con
\(P^{est}_i\) y la de referencia o groud truth que denotaremos con
\(P^{ref}_i\) con \(i=N\) la última timestamp.

El \(ATE_i\) entre cada pose \(P^{est}_i\) y \(P^{ref}_i\) al momento
\(i\) es la distancia entre las posiciones de las poses es decir:
\begin{align}
ATE_i = ||pos(P^{est}_i) - pos(P^{ref}_i)||
\end{align}

Con \(pos(P) \in \R^3\) el componente de posición o traslación de la
pose \(P\). Notar que no se considera el componente de rotación de las
poses.

Finalmente utilizamos el RMSE de los ATE al cuadrado: \begin{align}
RMSE = \sqrt{\frac{1}{N} \sum_{i=1}^{N} ATE_i^2}
\end{align}

Y es esta la métrica que se utiliza en esta sección para medir la
precisión en las trayectorias estimadas. Cabe aclarar además que las
poses dadas por la ground truth y las dadas por las estimaciones, en
general, no van a coincidir en sus timestamps y se debe utilizar alguna
forma de relacionarlas. Este trabajo utiliza la biblioteca EVO
\autocite{grupp2017evo} y por defecto esta simplemente utiliza la
trayectoria con menos poses y las relaciona a las poses con timestamps
más cercanas de la otra trayectoria. Otro detalle en el proceso que es
estándar en la práctica es alinear previamente las dos trayectorias
minimizando con cuadrados mínimos el error de las mismas con el trabajo
de \textcite{umeyamaLeastsquaresEstimationTransformation1991}.

Ahora sí, podemos ver los resultados en la \autoref{tab:ate}. En general
Basalt es también superior en estas métricas, pero aquí hay una
aclaración importante que hacer. ORB-SLAM3 es usualmente considerado el
estado del arte porque las métricas que reporta ocurren luego de que los
datasets hayan terminado de procesarse. Esto hace que el mapa utilizado
sea el final ya construido, lo cual afecta retroactivamente a las poses
computadas anteriormente. En este caso como estamos corriendo ORB-SLAM3
en tiempo real y utilizamos las poses que reporta inmediatamente con el
mapa en construcción, estas son mucho peores. Se deja como trabajo
próximo la evaluación de ORB-SLAM3 con el mapa pre construido, creemos
que en este caso deberían verse valores más parecidos a los reportados
por la publicación original del sistema y probablemente sobrepasen a
Basalt. También hay que notar que todos los datasets tienen duraciones
menores a 5 minutos, esto hace que el desvío que se acumula en sistemas
como Basalt que no tienen noción de mapa global no se note tanto. Con
esto queremos decir que es usual en Basalt notar luego de sesiones de
uso más largas que el entorno simulado empieza a cambiar de lugar
lentamente a causa de la deriva (\emph{drift}) usuales en sistemas de
VIO y no de SLAM. Se plantea también como trabajo a futuro mejorar este
aspecto de Basalt (ver discusión relacionada\footnote{\url{https://gitlab.com/VladyslavUsenko/basalt/-/issues/69}}).

\begin{table}[h]
\caption[Error en las trayectorias]{Error absoluto de la trayectoria (ATE) [m]}
\label{tab:ate}
\begin{addmargin*}[-0.2\textwidth]{-0.2\textwidth}
\resizebox{1.4\textwidth}{!}{
\begin{tabular}{ |l||l|l|l|l|l|l|  }
\hline
  \spacedlowsmallcaps{Dataset} & \spacedlowsmallcaps{BND} & \spacedlowsmallcaps{BNF} & \spacedlowsmallcaps{BO}  & \spacedlowsmallcaps{K} & \spacedlowsmallcaps{ON} & \spacedlowsmallcaps{OO}\\[3pt]
 \hline
EMH01 & 0.061 ± 0.023 & 0.061 ± 0.023 & 0.087 ± 0.026 & 0.290 ± 0.568   & 0.173 ± 0.230  & 0.216 ± 0.306 \\
EMH02 & 0.043 ± 0.022 & 0.043 ± 0.022 & 0.049 ± 0.023 & 0.127 ± 0.051   & 0.151 ± 0.133  & 0.627 ± 0.811 \\
EMH03 & 0.059 ± 0.019 & 0.059 ± 0.019 & 0.075 ± 0.039 & 0.192 ± 0.056   & 1.797 ± 1.175  & 2.513 ± 1.797 \\
EMH04 & 0.107 ± 0.038 & 0.107 ± 0.038 & 0.099 ± 0.040 & 0.188 ± 0.081   & 0.815 ± 0.517  & 2.065 ± 1.132 \\
EMH05 & 0.139 ± 0.041 & 0.139 ± 0.041 & 0.120 ± 0.041 & 0.206 ± 0.071   & 1.797 ± 0.785  & 3.537 ± 1.868 \\
EV101 & 0.040 ± 0.017 & 0.040 ± 0.017 & 0.040 ± 0.016 & 0.071 ± 0.027   & 9.842 ± 10.408 & 0.179 ± 0.168 \\
EV102 & 0.043 ± 0.013 & 0.043 ± 0.013 & 0.053 ± 0.019 & 0.093 ± 0.039   & 0.600 ± 0.359  & 0.951 ± 0.393 \\
EV103 & 0.049 ± 0.020 & 0.049 ± 0.020 & 0.067 ± 0.026 & 0.182 ± 0.050   & 13.274 ± 9.972 & 0.127 ± 0.105 \\
EV201 & 0.036 ± 0.015 & 0.036 ± 0.015 & 0.031 ± 0.017 & 0.046 ± 0.024   & 0.141 ± 0.130  & 0.098 ± 0.096 \\
EV202 & 0.045 ± 0.021 & 0.045 ± 0.021 & 0.060 ± 0.022 & 0.120 ± 0.041   & 0.323 ± 0.351  & 0.471 ± 0.248 \\
TR1   & 0.096 ± 0.048 & 0.096 ± 0.048 & 0.093 ± 0.042 & 4264.6 ± 2534.0 & 0.081 ± 0.028  & 0.546 ± 0.567 \\
TR2   & 0.067 ± 0.040 & 0.067 ± 0.040 & 0.062 ± 0.030 & 4447.9 ± 2728.6 & 0.087 ± 0.075  & 0.061 ± 0.082 \\
TR3   & 0.110 ± 0.057 & 0.110 ± 0.057 & 0.123 ± 0.063 & 6916.5 ± 4071.1 & 0.076 ± 0.032  & 0.123 ± 0.127 \\
TR4   & 0.050 ± 0.029 & 0.050 ± 0.029 & 0.049 ± 0.022 & 4918.2 ± 2749.4 & 0.105 ± 0.059  & 0.211 ± 0.175 \\
TR5   & 0.160 ± 0.067 & 0.160 ± 0.067 & 0.121 ± 0.051 & 5417.1 ± 2905.8 & 0.159 ± 0.122  & 0.112 ± 0.086 \\
TR6   & 0.018 ± 0.011 & 0.018 ± 0.011 & 0.018 ± 0.009 & 5003.9 ± 2511.9 & 0.105 ± 0.059  & 0.122 ± 0.168 \\
\hline
\textbf{Media} & \textbf{0.070 ± 0.030} & \textbf{0.070 ± 0.030} & \textbf{0.072 ± 0.030} & \textbf{1935.6 ± 1093.8} & \textbf{1.845 ± 1.527} & \textbf{0.747 ± 0.508} \\
\hline
\end{tabular}
}
\end{addmargin*}
\end{table}

\hypertarget{precisiuxf3n-de-los-movimientos}{%
\section{Precisión de los
movimientos}\label{precisiuxf3n-de-los-movimientos}}

Finalmente presentamos en la \autoref{tab:rte} el RMSE del \emph{error
relativo de la trayectoria (RTE)}. Este error es capaz de representar
las imprecisiones en tramos cortos de la trayectoria. Una forma de
pensar esto en el contexto de XR es qué tan mal se sentirán los
movimientos individuales realizados por un usuario.

Para definir este error primero dividimos la trajectoria \(P_k\) en
vectores definidos entre pares de timestamps \(i\) y \(j\):
\begin{align}
\delta_{ij} = pos(P_j) - pos(P_i) \ \in \R^3
\end{align}

Luego, definimos el error de traslación entre las timestamps \(i\) y
\(j\) como: \begin{align}
RTE_{ij} = ||\delta^{ref}_{ij} - \delta^{est}_{ij}||
\end{align}

Y finalmente el RMSE RTE queda definido como: \begin{align}
RMSE = \sqrt{\frac{1}{N} \sum_{\forall i, j} RTE_{ij}^2}
\end{align}

Notar que aquí la elección de que tan largo son los vectores
\(\delta_{ij}\) puede variar. En nuestro caso, usando EVO, utilizamos
vectores con timestamps separadas por el equivalente tiempo a 6 cuadros
de cada dataset, es decir unos 0.3, 0.2, y 0.1 segundos para 20, 30 y 60
fps.

\begin{table}[h]
\caption[Error en los movimientos]{Error relativo de la trayectoria (RTE, intervalos de 6 cuadros) [m]}
\label{tab:rte}
\begin{addmargin*}[-0.2\textwidth]{-0.2\textwidth}
\resizebox{1.4\textwidth}{!}{
\begin{tabular}{ |l||l|l|l|l|l|l|  }
\hline
  \spacedlowsmallcaps{Dataset} & \spacedlowsmallcaps{BND} & \spacedlowsmallcaps{BNF} & \spacedlowsmallcaps{BO}  & \spacedlowsmallcaps{K} & \spacedlowsmallcaps{ON} & \spacedlowsmallcaps{OO}\\[3pt]
 \hline
EMH01 & 0.004 ± 0.003 & 0.004 ± 0.003 & 0.004 ± 0.003 & 0.069 ± 0.283     & 0.138 ± 0.113 & 0.137 ± 0.110 \\
EMH02 & 0.004 ± 0.002 & 0.004 ± 0.002 & 0.004 ± 0.003 & 0.019 ± 0.019     & 0.140 ± 0.094 & 0.147 ± 0.167 \\
EMH03 & 0.009 ± 0.008 & 0.009 ± 0.008 & 0.010 ± 0.008 & 0.038 ± 0.030     & 0.368 ± 0.398 & 0.385 ± 0.460 \\
EMH04 & 0.010 ± 0.008 & 0.010 ± 0.008 & 0.011 ± 0.009 & 0.043 ± 0.031     & 0.335 ± 0.281 & 0.341 ± 0.392 \\
EMH05 & 0.009 ± 0.006 & 0.009 ± 0.006 & 0.010 ± 0.007 & 0.041 ± 0.030     & 0.307 ± 0.308 & 0.365 ± 0.660 \\
EV101 & 0.011 ± 0.006 & 0.011 ± 0.006 & 0.011 ± 0.006 & 0.044 ± 0.024     & 0.222 ± 1.958 & 0.136 ± 0.080 \\
EV102 & 0.011 ± 0.005 & 0.011 ± 0.005 & 0.011 ± 0.005 & 0.040 ± 0.022     & 0.277 ± 0.183 & 0.276 ± 0.188 \\
EV103 & 0.011 ± 0.007 & 0.011 ± 0.007 & 0.014 ± 0.009 & 0.039 ± 0.025     & 0.358 ± 2.249 & 0.246 ± 0.173 \\
EV201 & 0.003 ± 0.002 & 0.003 ± 0.002 & 0.003 ± 0.002 & 0.015 ± 0.012     & 0.092 ± 0.064 & 0.097 ± 0.081 \\
EV202 & 0.007 ± 0.006 & 0.007 ± 0.006 & 0.012 ± 0.025 & 0.025 ± 0.018     & 0.219 ± 0.148 & 0.221 ± 0.160 \\
TR1   & 0.007 ± 0.005 & 0.007 ± 0.005 & 0.008 ± 0.006 & 384.484 ± 305.665 & 0.505 ± 0.288 & 0.524 ± 0.294 \\
TR2   & 0.006 ± 0.005 & 0.006 ± 0.005 & 0.007 ± 0.006 & 468.756 ± 475.490 & 0.492 ± 0.421 & 0.503 ± 0.421 \\
TR3   & 0.005 ± 0.004 & 0.005 ± 0.004 & 0.006 ± 0.005 & 262.503 ± 201.940 & 0.618 ± 0.488 & 0.624 ± 0.486 \\
TR4   & 0.005 ± 0.005 & 0.005 ± 0.005 & 0.005 ± 0.005 & 342.893 ± 179.226 & 0.295 ± 0.161 & 0.300 ± 0.164 \\
TR5   & 0.009 ± 0.007 & 0.009 ± 0.007 & 0.010 ± 0.008 & 341.326 ± 155.828 & 0.477 ± 0.284 & 0.483 ± 0.285 \\
TR6   & 0.003 ± 0.002 & 0.003 ± 0.002 & 0.003 ± 0.002 & 355.299 ± 219.485 & 0.268 ± 0.214 & 0.275 ± 0.227 \\
\hline
\textbf{Media} & \textbf{0.007 ± 0.005} & \textbf{0.007 ± 0.005} & \textbf{0.008 ± 0.007} & \textbf{134.727 ± 96.133} & \textbf{0.319 ± 0.478} & \textbf{0.316 ± 0.272} \\
\hline
\end{tabular}
}
\end{addmargin*}
\end{table}

\hypertarget{conclusiones-y-trabajo-futuro}{%
\chapter{Conclusiones y trabajo
futuro}\label{conclusiones-y-trabajo-futuro}}

En este trabajo se estudiaron distintos sistemas de SLAM/VIO en el
contexto de localización en tiempo real para XR. Vimos algunos de los
conceptos fundamentales que estos utilizan como el algoritmo de Gauss
Newton para resolver problemas de optimización no lineal, de los cuales
el área de visión por computadora está plagado, y SLAM no es la
excepción. También vimos las distintas formas de representar
transformaciones y rotaciones en dos y tres dimensiones: ángulos euler,
cuaterniones, ángulo axial, matrices de rotación y una mirada práctica
sobre los grupos de Lie \(SO(n)\) y \(SE(n)\) junto a sus álgebras de
Lie \(\so(n)\) y \(\se(n)\).

Posteriormente nos adentramos en la implementación de la capa de
odometría visual-inercial de Basalt. Esto permitió ver de primera mano
los distintos tipos de algoritmos que se reúnen en este tipo de
sistemas. Se integró Kimera-VIO, ORB-SLAM3 y Basalt a Monado, el runtime
OpenXR de código libre. Para esto hizo falta diseñar una interfaz
eficiente que generaliza de forma razonable estos sistemas. Se
analizaron los problemas de implementación particulares a considerar
para XR como la predicción y filtrado de poses, o como lidiar con la
imperfección de los sensores de cámara e IMU. Se contribuyeron a Monado
todas estas mejoras, incluyendo la extensión de dos controladores de
dispositivos que ahora son capaces de aprovechar este tipo de tracking.
Uno de ellos es una plataforma VR de producción comercial que ahora
puede ser utilizada por usuarios entusiastas que deseen utilizar este
tipo de hardware en GNU/Linux con un stack de software completamente
libre.

Este proyecto plantea las bases de infraestructura en Monado para este
tipo de sistemas, pero aún hay mucho por hacer y por mejorar para lograr
tracking con calidades similares a las que se encuentran en productos
comerciales. Se plantea como trabajo futuro:\newline

\begin{itemize}
\item
  Mejorar la experiencia de usuario para al utilizar el SLAM tracker en
  Monado. El trabajo realizado actualmente puede resultar un poco
  complejo de instalar y configurar para un usuario inexperto.
\item
  Permitir el uso de múltiples implementaciones de SLAM/VIO de forma
  dinámica. Es decir, poder tener distintas implementaciones corriendo
  en simultáneo y localizando a distintos dispositivos.
\item
  Hay espacio de mejora en el rendimiento de las implementaciones. En
  general, en este trabajo nos limitamos a hacer lo justo y necesario
  para que el tracking funcione a tiempos razonables y no retrase el
  pipeline de Monado. Más aún, parece existir poca cantidad de trabajos
  que apliquen unidades de cómputo masivamente paralelas, como lo son
  las GPU, al problema de localización visual-inercial. Creemos que
  existen posibles ganancias de eficiencia en esta línea de trabajo.
\item
  Sería bueno extender Basalt para soportar algún tipo de mapeo global
  en tiempo real que permita tener trayectorias consistentes que no
  tiendan a moverse lentamente con el tiempo. Discusiones de esto
  referenciada en una nota al pie\footnote{\url{https://gitlab.com/VladyslavUsenko/basalt/-/issues/69}}.
\item
  Sería bueno mejorar las formas de testeo y evaluación de sistemas SLAM
  en Monado, poder automatizarlas e integrarlas en los procesos de
  integración continua del proyecto. Esto permitiría el impacto que
  nuevos cambios traen al rendimiento y la precisión del sistema.
\item
  Existen pocos conjuntos de datos para SLAM aptos para XR (p.~ej.
  TUM-VI \autocite{schubertBasaltTUMVI2018}), y ante la dificultad de
  producirlos, sería ideal aprovecharse de las herramientas
  fotorrealistas que son fácilmente accesibles en la actualidad para la
  generación de datos sintéticos.
\item
  Existen métodos de predicción más eficientes que podrían adaptarse a
  Monado en lugar del método ad hoc desarrollado en este trabajo. En
  particular el mismo trabajo de preintegración de muestras de IMU
  utilizado en Basalt
  \autocite{forsterOnManifoldPreintegrationRealTime2017} puede ser un
  muy buen punto de partida para un método de predicción más preciso.
\item
  Finalmente, muchos de los módulos que forman parte de estos sistemas
  son útiles de manera individual. La integración de estos en Monado
  podría beneficiar a distintos controladores que quieran hacer uso de
  algoritmia de visión por computadora específica en otros contextos.
\end{itemize}

%********************************************************************
% Other Stuff in the Back
%*******************************************************
\cleardoublepage\include{FrontBackmatter/Bibliography}
% \cleardoublepage\include{FrontBackmatter/Declaration}
% \cleardoublepage\include{FrontBackmatter/Colophon}
% ********************************************************************
% Game Over: Restore, Restart, or Quit?
%*******************************************************
\end{document}
% ********************************************************************
